\documentclass[11pt]{article}
\usepackage[a4paper,margin=29mm]{geometry}
\usepackage[T1]{fontenc}
\usepackage{lmodern}
\usepackage{microtype}
\usepackage{amsmath,amssymb,amsthm,mathtools}
\usepackage{mathrsfs}
\usepackage{booktabs}
\usepackage{array}
\usepackage{enumitem}
\usepackage[dvipsnames]{xcolor}
\usepackage[colorlinks=true,linkcolor=MidnightBlue,citecolor=BrickRed,urlcolor=RoyalBlue]{hyperref}
\usepackage[nameinlink,noabbrev]{cleveref}
\emergencystretch=3em
\hypersetup{
  pdftitle={Moving-Level Trace Towers, Exact Ramanujan Transmutations, and Arithmetic Barrier Theorems, Version 2.3},
  pdfauthor={KEIO},
  pdfsubject={Exact zero-trace towers, Meijer-G/hypergeometric duality, Gamma-jet and probabilistic realizations, and sharp arithmetic barriers},
  pdfkeywords={Euler-Mascheroni constant, exponential integral, spectral trace, Ramanujan Challenge, Meijer G-function, hypergeometric function, Gamma deformation, arithmetic independence}
}

\newtheorem{theorem}{Theorem}[section]
\newtheorem{proposition}[theorem]{Proposition}
\newtheorem{corollary}[theorem]{Corollary}
\newtheorem{lemma}[theorem]{Lemma}
\newtheorem{conjecture}[theorem]{Conjecture}
\theoremstyle{definition}
\newtheorem{definition}[theorem]{Definition}
\newtheorem{problem}[theorem]{Problem}
\theoremstyle{remark}
\newtheorem{remark}[theorem]{Remark}

\newcommand{\Q}{\mathbb Q}
\newcommand{\Qbar}{\overline{\mathbb Q}}
\newcommand{\C}{\mathbb C}
\newcommand{\R}{\mathbb R}
\newcommand{\N}{\mathbb N}
\newcommand{\Z}{\mathbb Z}
\newcommand{\Ein}{\operatorname{Ein}}
\newcommand{\Cin}{\operatorname{Cin}}
\newcommand{\Ei}{\operatorname{Ei}}
\newcommand{\Ci}{\operatorname{Ci}}
\newcommand{\Chi}{\operatorname{Chi}}
\newcommand{\Eexp}{\mathcal E_{\mathrm{exp}}}
\newcommand{\trdeg}{\operatorname{trdeg}}
\newcommand{\Tr}{\operatorname{Tr}}
\newcommand{\Norm}{\operatorname{N}}
\newcommand{\LP}{\mathcal{LP}}
\newcommand{\ord}{\operatorname{ord}}
\newcommand{\Span}{\operatorname{span}}
\newcommand{\mult}{\operatorname{mult}}
\newcommand{\calG}{\mathcal G}
\newcommand{\lcm}{\operatorname{lcm}}
\newcommand{\diag}{\operatorname{diag}}
\newcommand{\dist}{\operatorname{dist}}
\newcommand{\FPTr}{\operatorname{FPTr}}
\newcommand{\Var}{\operatorname{Var}}
\newcommand{\Prob}{\mathbb P}
\newcommand{\stirling}[2]{\genfrac{[}{]}{0pt}{}{#1}{#2}}

\title{\bfseries Moving-Level Trace Towers,\\
Exact Ramanujan Transmutations,\\
and Arithmetic Barrier Theorems}
\author{KEIO}
\date{Version 2.3 --- 1 September 2026}

\begin{document}
\maketitle

\begin{abstract}
Let \(K=\Eexp:=\Qbar(e^\alpha:\alpha\in\Qbar)\), and use the algebraic
independence over \(K\) of the nonzero algebraic-point values of the entire
exponential integral \(\Ein\), established in a companion paper.  For every
nonzero level \(c\), the absolutely convergent inverse-power traces of the
divisor of \(\Ein(z)-c\) are universal polynomials
\[
 T_m(c)=P_m(c^{-1}),\qquad
 P_m(X)=-m[z^m]\log(1-X\Ein(z)),
\]
and the two even traces \(T_2,T_4\) recover \(c\) exactly.  Balanced
algebraic-point levels decreasing doubly exponentially to Euler's constant
therefore yield countable transversal towers of algebraically independent
exact traces.  At orders two and four their positive increments give
trace-class carriers for \(S_2\) and \(S_4\), whose traces recover
\(\gamma\); at order two the carrier additionally gives an order-zero
Laguerre--P\'olya determinant.  We also give finite-rank cancellation and
cofinite-independence theorems for multiplicative groups of algebraic
arguments.  The same argument yields cofinite algebraic independence for
the tails \(E_1(\alpha)\), with exactly one additional possible defect, and
an all-order version for regularized lower incomplete-Gamma jets.  We
classify the unique nontrivial collision of the pair
\((T_2,T_3)\), and exhibit, at every order, algebraic levels with vanishing
inverse trace but wholly transcendental spectral divisors.

The 2026 Ramanujan Challenge recurrence is developed into five explicitly
linked realizations.  Its terminating \({}_2F_2\) state and the official
Meijer--\(G\) state are related by an explicit rational \(SL_3\) coboundary;
the resulting conserved primal--dual pairing gives a finite-index
Meijer--\(G\)/\({}_2F_2\) bilinear identity.  An Euler--Beta--log transform
identifies the adjacent family with the cubic Pilehrood construction and
provides an explicit fixed-offset Beta intertwiner in the present
normalization.  A single rational deformation has
zero- and first-jets equal to the denominator and numerator, and its
normalization at \(s=0\) converges locally uniformly to
\[
 \Gamma(1-s)^3\Gamma(1+s)^2.
\]
Consequently the Challenge ratio is the first logarithmic jet of a fixed
rational deformation whose higher logarithmic jets recover every
\(\zeta(m)\).  The same
deformation produces a monic integer contact polynomial whose low
coefficients contain the two recurrence solutions while its high
coefficients impose \(2a+1\) exact Newton-moment identities.

Probabilistically, every finite rational approximant is the mean of an
explicit mixture of five exponential maxima.  The variables converge with
all moments to a signed sum of five independent Gumbel variables, whose
moment generating function is the preceding Gamma product.  The official
Meijer state is simultaneously a stop-loss expectation, a survival
probability, and a boundary density.  A separate Hilbert--Schmidt operator
realizes the full Gamma product as a regularized Fredholm determinant: its
odd traces of order at least three are the odd zeta values, and its even
traces are five times the even zeta values.  We determine the complete rational span of all
Ap\'ery limits of the recurrence, strengthen its forced common divisor, and
explain the factorial-six Casoratian through the denominator-cleared dual
gauge and the cyclic endpoint gauge.

The paper also proves simultaneous real and \(p\)-adic convergence along a
common Euler--Pad\'e subsequence and full unit-boundary convergence at
\(t=-1\) for \(p=2,3,5,13\).  The final part retains sharply scoped
obstruction theorems: the critical moving-point and coefficient-height
barriers, their two-base extension, the Pr\'evost fixed-place Ridout deficit,
universal saddle cancellation for a fixed-prime Rodrigues deformation, and
the separation between gamma-free unframed Stokes data and the framed
Euler--Tate extension.  A corrected shift calculation identifies every
finite multi-level trace field with the exact \(\mathbb G_a\)-invariant
field; analytic completion at the moving limit is precisely where that
formal symmetry breaks.

Finally, an explicit mixed arithmetic-Gevrey family gives a simultaneous
conditional dilation theorem under Fischler--Rivoal's Conjecture~3.  This is
an application, not a new proof of the already known conditional implication
from that conjecture to the transcendence of \(\gamma\).  The Euler--Tate
matched-line target is sharpened: positive lifting would force
\(1,\pi,\gamma\) to be linearly independent over \(\Qbar\), and hence both
\(\gamma\) and \(\gamma/\pi\) to be transcendental.

Priority is separated explicitly: the official Meijer--\(G\) solution, the
classical Bell--zeta tower, the Gumbel probability identity, and the
multiple-Laguerre background are credited to their sources.  No
unconditional theorem here proves the irrationality or transcendence of
\(\gamma\), \(\delta\), any
odd zeta value, or \(S_2\).
\end{abstract}

\tableofcontents

\section{Introduction and exact scope}

Write
\[
 \Ein(z)=\int_0^z\frac{1-e^{-t}}t\,dt
 =\sum_{r\ge1}\frac{(-1)^{r-1}z^r}{r\,r!},
\]
and let
\[
 \gamma=\lim_{n\to\infty}\left(\sum_{k=1}^n\frac1k-\log n\right),
 \qquad
 \delta=eE_1(1).
\]
Neither \(\gamma\) nor \(\delta\) is known to be irrational; see
Lagarias' survey \cite{Lagarias2013}.  The elementary positive-axis identity
\begin{equation}\label{eq:Ein-connection}
 \Ein(x)=\gamma+\log x+E_1(x)\qquad(x>0)
\end{equation}
places \(\gamma\) at the boundary between algebraic-point \(E\)-values and an
asymptotic connection constant.

The companion paper \cite{KEIO2026I} determined the algebraic relations among
all nonzero algebraic-point \(\Ein\)-values over the pure exponential field,
constructed several finite Euler witnesses, and computed the first two
inverse traces of the divisor of \(\Ein-\gamma\).  The present paper begins
where that finite theory stops.  Its organizing question is:
\begin{quote}
What remains true when the level of the divisor is allowed to move through a
countable algebraically independent family converging to \(\gamma\)?
\end{quote}

The results have four layers.  In the spectral layer, the two even traces
of orders two and four
recover the level exactly, so the entire fixed-level inverse-trace ring is
one-dimensional; moving the level restores one algebraically
independent coordinate per level.  The order-two row has a positive
telescoping structure and hence a genuine operator-theoretic realization.
Algebraic levels of a natural gamma-free inverse map have
transcendental preimages.  All finite linear witnesses are governed by
one normalized-certificate uniqueness lemma.  The dyadic tail is the
unique minimizer in a natural cumulative-positive augmentation cone, while
outside that cone three-point forms are already dense.  The resulting
coefficient-height transition identifies the exact scale at which the first
tail can be cancelled.  In the transmutation layer, the four-term Euler
recurrence from the 2026 Ramanujan Challenge is connected explicitly to the
official Meijer--\(G\) construction, the cubic Pilehrood family, a finite
probability law, a regularized Gamma determinant, and a full rational
zeta-jet tower.  Its denominator and numerator live in a flat dual-number
\(SL_3\) system, equivalently an ordinary \(SL_6\) system.  In the
obstruction layer, the particular residual-prime,
fixed-place, Rodrigues, natural-index P-recursive, and strict-Stokes
constructions studied here yield scoped saturation identities, budget
inequalities, or sharply stated missing inputs.

The principal results are as follows.

\begin{theorem}[Universal trace and transversal independence]
\label{thm:intro-universal}
For \(c\in\C^\times\), let \(Z_c\) be the zero multiset of
\(\Ein(z)-c\).  For every \(m\ge2\),
\[
 T_m(c):=\sum_{\rho\in Z_c}\rho^{-m}=P_m(c^{-1}),
\]
where \(P_m\in\Q[X]\) is monic of degree \(m\).  If
\(\{c_j:j\in J\}\) is algebraically independent over a characteristic-zero
field \(F\), and one integer \(m_j\ge2\) is selected for each \(j\), then
\(\{T_{m_j}(c_j):j\in J\}\) is algebraically independent over \(F\).
Moreover
\[
 c^{-1}=144T_4(c)-144T_2(c)^2-14T_2(c),
\]
so \(F[T_m(c):m\ge2]=F[c^{-1}]\).
The map \(c\mapsto(T_2(c),T_3(c))\) is injective except for the single
unordered pair
\[
 \{6-2\sqrt3,6+2\sqrt3\},
\]
whose common trace pair is \((-1/24,1/96)\).  For each \(m\ge2\), the
nonzero roots of \(P_m\) also give a nonempty finite set of algebraic levels
at which \(T_m\) vanishes and every point of the corresponding divisor is
transcendental.
\end{theorem}

\begin{theorem}[Finite-rank tails and incomplete-Gamma jets]
\label{thm:intro-tail-closures}
Let \(\Gamma\subset\Qbar_{>0}^{\times}\) have multiplicative rank \(r\).
For every finite \(A\subset\Gamma\),
\[
 \trdeg_KK(E_1(\alpha):\alpha\in A)\ge |A|-r-1.
\]
After deleting at most \(r+1\) values, the remaining full \(E_1\)-tail
family is jointly algebraically independent over \(K\).  The same bound
holds for the combined classical family
\(\{E_1(\alpha),\Ei(\alpha),\Ci(\alpha),\operatorname{Si}(\alpha)\}\).

For the regularized lower incomplete-Gamma jets
\[
 R_{\alpha,k}=[s^k]\{s^{-1}-\Gamma_{<}(s,\alpha)\},
\]
every finite \(S\subset\Gamma\times\N_0\) satisfies
\[
 \trdeg_KK(R_{\alpha,k}:(\alpha,k)\in S)\ge |S|-r.
\]
Thus at most \(r\) exceptions need be removed from the complete two-index
lower-jet family.  The corresponding upper jets have defect at most
\(r\) plus the number of distinct jet orders.
\end{theorem}

\begin{theorem}[Dyadic moving zero-trace tower]
\label{thm:intro-dyadic}
For \(N\ge1\), put
\[
 Y_N=(N+1)\Ein(2^N)-N\Ein(2^{N+1}).
\]
Then \(\{Y_N:N\ge1\}\) is countably algebraically independent over \(K\),
\(Y_N\downarrow\gamma\), and
\[
 Y_N-\gamma\sim \frac{N+1}{2^N}e^{-2^N}.
\]
For every choice \(m_N\ge2\), the exact traces
\[
 \tau_{m_N,N}:=\sum_{\Ein(\rho)=Y_N}\rho^{-m_N}
\]
are countably algebraically independent over \(K\).  For fixed \(m\),
\(\tau_{m,N}\to S_m\), the \(m\)-th inverse trace of \(\Ein-\gamma\).
\end{theorem}

\begin{theorem}[Inverse levels and critical coefficients]
\label{thm:intro-inverse}
The function
\[
 \mathscr R(q)=E_1(1)-2E_1(q)+E_1(q^2)
\]
maps \((1,\infty)\) increasingly onto \((0,E_1(1))\).  The inverse image of
every algebraic number in this interval is transcendental.  In particular,
the unique roots \(r_m>1\) of \(\mathscr R(r_m)=1/m\), \(m\ge6\), are all
transcendental.

For each fixed algebraic \(q>1\), the critical three-point cancellation
coefficients \(u_N(q)\) defined in \eqref{eq:uNq} satisfy, for every finite
\(I\),
\[
 \trdeg_KK(u_N(q):N\in I)\ge |I|-1.
\]
\end{theorem}

\begin{theorem}[Positive spectral realizations]
\label{thm:intro-positive}
Set
\[
 \tau_N=\frac1{Y_N^2}-\frac1{2Y_N},\qquad
 d_0=\tau_1,\qquad d_N=\tau_{N+1}-\tau_N\quad(N\ge1).
\]
Then \(d_N>0\), the \(d_N\) are countably algebraically independent over
\(K\), and
\[
 S_2=\sum_{N\ge0}d_N.
\]
The diagonal operator \(\mathsf A e_N=d_Ne_N\) belongs to every Schatten
class and has trace \(S_2\).  Moreover
\[
 \Delta_2(z)=\det\!\left(I-\frac{z^2}{2}\mathsf A\right)
 =\prod_{N\ge0}\left(1-\frac{d_N}{2}z^2\right)
\]
is an order-zero Laguerre--P\'olya entire function.  Its positive zeros are
countably algebraically independent over \(K\), and
\[
 \sum_{\rho\in Z(\Delta_2)}\rho^{-2}=S_2.
\]
The analogous order-four increments define a positive diagonal operator
\(\mathsf A_4\) with trace \(S_4\), and, writing
\(\mathsf A_2=\mathsf A\),
\[
 \gamma=
 \left\{144\Tr\mathsf A_4
 -144(\Tr\mathsf A_2)^2-14\Tr\mathsf A_2\right\}^{-1}.
\]
\end{theorem}

\begin{theorem}[Cross-index witness rigidity]
\label{thm:intro-witness}
Let \(q_*\in(0,1)\) be the real solution of \(\Ein(q_*)=\gamma\).  Let
\(\beta_n\in(0,1)\) be the positive real zero of
\[
 G_n(z)=n\Ein(z)-\Ein(z^n)-(n-1)\gamma,
\]
and, for odd \(n\ge3\), let \(\lambda_n<-1\) be its negative real zero.
Then at most one member of
\[
 \{q_*\}\cup\{\beta_n:n\ge2\}
 \cup\{\lambda_n:n\ge3,\ n\text{ odd}\}
\]
is algebraic.
\end{theorem}

\begin{theorem}[Critical moving-point barrier]
\label{thm:intro-barrier}
Any theorem of the natural form
\[
 |P(F(\xi))|\ge
 \exp\!\bigl(-C(1+|\xi|)^c\Phi(\deg P,\log H(P))\bigr)
\]
intended to hold uniformly for all fixed \(E\)-function systems and positive
integral sample points, or already for all nonzero polynomial values of the
displayed system, must have \(c\ge1\).  The obstruction occurs for the fixed
system
\[
 F=(\Ein(z),\Ein(2z),\Ein(4z))
\]
and the fixed linear form \(P=X_1-2X_2+X_3\).  General sufficiency at
\(c=1\) is not known; at that scale the same example forces at least the
polynomial loss \(\xi^{-1}\).  Within the augmentation-one stratum the
dyadic coefficients uniquely minimize the tail in the cumulative-positive
cone.  Outside that cone three points give dense tail values, while
cancellation of the first active tail requires at least
\(\log H\ge(q-1)q^N+O(1)\).  For consecutive augmentation-one triples this
transition has a matching construction; at natural point cost a normalized
lower bound for that affine template must pay height exponent at least two.

For the dyadic tower, algebraicity of \(S_2\) is
equivalent to the existence of one algebraic coefficient \(s\) for which
\[
 -\log|2sY_N^2+Y_N-2|\sim2^N.
\]
\end{theorem}

\begin{theorem}[Exact Ramanujan transmutation package]
\label{thm:intro-recurrence}
The two distinguished solutions of the Euler recurrence
\eqref{eq:challenge-recurrence} have the finite sums
\eqref{eq:qn-closed}--\eqref{eq:pn-closed}, satisfy the recurrence for every
\(n\ge0\), and obey \(p_n/q_n\to\gamma\).  Their terminating
\({}_2F_2\) state and the official Meijer--\(G\) state are joined by the
rational coboundary \eqref{eq:ram-coboundary}; the normalized conserved
pairing is the finite identity \eqref{eq:ram-bilinear-matrix}.  The
denominator system embeds in the flat \(SL_3\) field
\eqref{eq:Ma}--\eqref{eq:cmf-flatness}, and its zero--first jet lift is the
flat \(SL_6\) system \eqref{eq:SL6-block}.

The Euler--Beta--log transform
\eqref{eq:q-Beta-transform}--\eqref{eq:p-Beta-transform} identifies the
same solutions with the cubic Pilehrood family.  The rational carrier
\(\mathscr K_a\) has denominator and numerator as its zero- and first-jets;
\(\mathscr K_a(s)/\mathscr K_a(0)\) converges to
\(\Gamma(1-s)^3\Gamma(1+s)^2\) and yields the zeta limits
\eqref{eq:zeta-log-jets}.  Its scaled integer numerator is the monic contact
polynomial of \cref{thm:contact-polynomial}.  All recurrence solutions have
the limiting form \eqref{eq:all-initial-limits}.  Moreover, with
\(b=n+4\) and \(D_b=\lcm(1,\ldots,b)\),
\[
 \frac{b!}{D_b}\gcd\!\left(b,\frac{D_b}{b}\right)
 \mid\gcd(p_n,q_n),
\]
and the factorial-six part of the exact Casoratian
\eqref{eq:challenge-casoratian} comes from the determinant of the
denominator-cleared official dual transfer.
\end{theorem}

\begin{theorem}[Arithmetic and Stokes barriers]
\label{thm:intro-new-barriers}
The following statements hold unconditionally.
\begin{enumerate}[label=\textup{(\roman*)}]
 \item The dyadic level \((Y_N)\) and its order-two trace \((\tau_N)\) are
 not P-recursive.
 \item The two-base form of \cref{prop:two-base-obstruction} remains
 exponentially small with
 \[
  -\log\{Y_N(2)-Y_{M_N}(3)\}
  =2^N+N\log2-\log(N+1)+O(1).
 \]
 Thus multiplicative independence of the two bases does not remove this
 moving-point scale.
 \item Every finite multi-level trace field is exactly the
 \(\mathbb G_a\)-invariant field of the gauge shift
 \(G\mapsto G+t,H_i\mapsto H_i-t\).  The completed moving tower breaks this
 formal symmetry through its boundary limit.
 \item Every nondegenerate rank-one real Hilbert quotient obtained from the
 fixed residual-direction constraints is exactly saturated under a candidate
 \(\gamma=a/b\).  For the optimally saddled uncombined Pr\'evost form, the
 lcm-only height-one ledger has deficit
 \(2-\log(3+2\sqrt2)\).
 \item The fixed-prime Rodrigues family \eqref{eq:fixed-p-rodrigues}
 places \(p^{n-\ell}\) in the endpoint coefficient.  Even if this whole
 factor is credited as fixed-place gain, the resulting ledger and the
 archimedean saddle satisfy \eqref{eq:universal-fixed-p-bound}.
 \item The displayed unframed lateral transition of the elementary
 differential operator annihilating \(F=-\Ein\) is gamma-free.  The Euler
 constant occurs in the connection framing and, under
 \(\gamma\in\Qbar\), can be removed by an algebraic shear.  A framed
 Euler--Tate matched-line lifting theorem would imply the
 \(\Qbar\)-linear independence of \(1,\pi,\gamma\), and hence the
 transcendence of \(\gamma\), \(\gamma/\pi\), and \(S_2\).
\end{enumerate}
\end{theorem}

The qualification ``transversal'' in
\cref{thm:intro-universal,thm:intro-dyadic} is essential: one chooses one
trace order at each level.  At a fixed level all trace orders are algebraic
over the same one-dimensional level field.  Likewise, the infinite sum in
\cref{thm:intro-positive} has no automatic arithmetic consequence.  An
infinite sum of algebraically independent transcendental numbers may be
algebraic.

\paragraph{Organization.}
\Cref{sec:input} records the exact input from the companion paper.
\Cref{sec:rank-deficit} proves finite-rank cofinite independence for both
connection coordinates and classical tails, together with the complete
cancellation geometry; \cref{sec:external-realizations} exports the
engine to sine-integral graphs, generalized Dickman transforms, and
regularized incomplete-Gamma jets; the
characteristic-zero algebraic-matroid realization is isolated in
\cref{app:Ein-matroids} because it uses an external linearization theorem.
\Cref{sec:trace-polynomials} develops the universal trace calculus and
identifies its finite multi-level field with the exact shift quotient.
\Cref{sec:dyadic} proves the dyadic and all-order moving towers.
\Cref{sec:spectral} constructs the two positive operators and the
order-two Laguerre--P\'olya determinant.  \Cref{sec:other-towers} gives fixed-coefficient
power-orbit towers, transcendental inverse levels, and the logarithmic
lattice.  \Cref{sec:certificates} proves normalized certificate uniqueness,
cross-index rigidity, and the one-defect cancellation theorem.
\Cref{sec:relative} gives the relative five-class exclusion and the exact
finite-descent criterion.
\Cref{sec:ramanujan-recurrence,sec:cmf} solve the four-term Euler recurrence,
classify all of its limits, and construct the flat \(SL_3/SL_6\) field.
\Cref{sec:meijer-duality,sec:beta-log} give the exact official gauge and the
Euler--Beta--log offset ladder.  \Cref{sec:gamma-jets,sec:ram-probability}
develop the rational Gamma jets, contact polynomial, finite probability
law, and zero geometry; \cref{sec:gamma-determinant} gives the regularized
determinant realization.  \Cref{sec:cross-place-pade} proves the real and
\(p\)-adic Euler--Pad\'e convergence results.
\Cref{sec:barrier} proves the sharp inside/outside augmentation dichotomy and
isolates the critical moving-point obstruction.
\Cref{sec:sondow} replaces an invalid asymmetric-integral heuristic by an
exact evaluation.  \Cref{sec:prevost-grid,sec:fixed-prime-rodrigues}
establish the residual-prime, Ridout, and Rodrigues barriers.
\Cref{sec:stokes-framing,sec:mixed-gevrey,sec:euler-tate} separate unframed
Stokes data from the mixed arithmetic-Gevrey and canonical motivic
framings, prove the conditional dilation and Gamma-jet ranks, and isolate
the strengthened matched-line target.  The final
section records what remains open.

\paragraph{Dependency and status ledger.}
To keep imported input, internal proof, and conditional extrapolation
separate, the logical dependencies are summarized below.
\begin{center}
\small
\begin{tabular}{@{}
 >{\raggedright\arraybackslash}p{0.24\textwidth}
 >{\raggedright\arraybackslash}p{0.32\textwidth}
 >{\raggedright\arraybackslash}p{0.34\textwidth}@{}}
\toprule
Source & Imported statement & Status in this paper\\
\midrule
Companion paper \cite[Thms.~1.4, 5.1 and Lem.~11.2]{KEIO2026I}
& Algebraic independence of algebraic-point resonant \(\Ein_q\)-values
and of the balanced dyadic coordinates
& Standing arithmetic input; the tail, incomplete-Gamma, trace, carrier,
and rigidity deductions are proved here\\
Fischler--Rivoal \cite{FischlerRivoalMixed2024}
& Definitions of mixed functions; conjecture
\(\mathbf E\cap\mathbf D=\Qbar\); conditional linear specialization
& The dilation and all-order Gamma-jet families are constructed here;
every conclusion using the intersection conjecture is explicitly
conditional\\
Official Challenge solution \cite[\S5.2]{RamanujanOfficial2026}
& Official Meijer--\(G\) normalization and Miller-selected limiting direction
& Imported normalization; the rational coboundary, bilinear identity, and
transmutations are proved here\\
Rosen--Sidman--Theran \cite[Thm.~1.10]{RosenSidmanTheran2025}
& Linear representability over an algebraically closed characteristic-zero
field
& Used only in Appendix~\ref{app:Ein-matroids}\\
Powers \cite[\S5]{Powers2025}
& A concluding multisection algebraic-independence assertion
& Used only under an explicit hypothesis in
Remark~\ref{rem:Powers-conditional}; no unconditional result depends on it\\
\bottomrule
\end{tabular}
\end{center}

\section{Input from the companion paper}
\label{sec:input}

Throughout,
\[
 K=\Eexp=\Qbar(e^\alpha:\alpha\in\Qbar),
\]
and \(K^{\mathrm{alg}}\) denotes its relative algebraic closure in \(\C\).
We use the following theorem from \cite[Thm.~5.1]{KEIO2026I}.

\begin{theorem}[Algebraic-point \(\Ein\)-independence]
\label{thm:input-independence}
The countable family
\[
 \{\Ein(\alpha):\alpha\in\Qbar^\times\}
\]
is algebraically independent over \(K\), in the finite-subset sense.
Consequently it remains algebraically independent over
\(K^{\mathrm{alg}}\).
\end{theorem}

For the incomplete-Gamma application below we also use the resonant
extension of this input.  Put
\[
 \Ein_q(z)=\sum_{n\ge1}\frac{(-1)^{n-1}z^n}{n^q n!}
 \qquad(q\ge1).
\]
By \cite[Thm.~1.4]{KEIO2026I}, the family
\begin{equation}\label{eq:resonant-input}
 \{\Ein_q(\alpha):q\ge1,\ \alpha\in\Qbar^\times\}
\end{equation}
is algebraically independent over \(K\), in the finite-subset sense.

The final assertion is standard: algebraic independence is preserved under
algebraic extension of the base field.  We shall repeatedly use the following
linear consequence.

\begin{lemma}[Full-rank linear images]
\label{lem:linear-images}
Let \(x_1,\ldots,x_r\) be algebraically independent over a field \(F\), and
let \(A\in M_{s\times r}(F)\) have row rank \(s\).  Then the \(s\) coordinates
of \(A(x_1,\ldots,x_r)^{\mathsf T}\) are algebraically independent over \(F\).
\end{lemma}

\begin{proof}
Complete the rows of \(A\) to an invertible \(r\)-by-\(r\) matrix over \(F\).
The resulting linear change of variables is an automorphism of the rational
function field \(F(x_1,\ldots,x_r)\).
\end{proof}

For later comparison we also recall the Euler-level traces from
\cite[Thm.~14.5]{KEIO2026I}.  If \(Z_\gamma\) is the zero multiset of
\(\Ein(z)-\gamma\), then
\begin{align}
 S_2&:=\sum_{\rho\in Z_\gamma}\rho^{-2}
 =\frac1{\gamma^2}-\frac1{2\gamma},\label{eq:S2}\\
 S_3&:=\sum_{\rho\in Z_\gamma}\rho^{-3}
 =\frac1{\gamma^3}-\frac3{4\gamma^2}+\frac1{6\gamma},\label{eq:S3}
\end{align}
and
\begin{equation}\label{eq:gamma-recovery}
 \gamma=\frac{24S_2+1}{24S_3+6S_2}.
\end{equation}
In particular \(S_2\) is algebraic if and only if \(\gamma\) is algebraic.

\section{Finite-rank independence and cancellation geometry}
\label{sec:rank-deficit}

The algebraic-point theorem has a useful export form when the arguments lie
in a multiplicative group of finite rank.  We restrict here to positive
algebraic arguments, so that the real logarithm is single-valued and no
additional \(2\pi i\)-coordinate occurs.  Put
\begin{equation}\label{eq:c-alpha}
 c_\alpha=\Ein(\alpha)-\log\alpha=\gamma+E_1(\alpha)
 \qquad(\alpha\in\Qbar_{>0}).
\end{equation}

\begin{theorem}[Finite-rank deficit and cofinite independence]
\label{thm:cofinite-independence}
Let \(\Gamma\subset\Qbar_{>0}^{\times}\) have multiplicative rank \(r\).
For every finite \(A\subset\Gamma\),
\begin{equation}\label{eq:finite-rank-deficit}
 \trdeg_KK(c_\alpha:\alpha\in A)\ge |A|-r.
\end{equation}
Consequently there is a set \(E\subset\Gamma\), \(|E|\le r\), such that
\begin{equation}\label{eq:cofinite-family}
 \{c_\alpha:\alpha\in\Gamma\setminus E\}
 \quad\hbox{is algebraically independent over }K.
\end{equation}

Fix \(\alpha_0\in\Gamma\).  There is likewise a set
\(E_{\alpha_0}\subset\Gamma\setminus\{\alpha_0\}\),
\(|E_{\alpha_0}|\le r\), for which
\begin{equation}\label{eq:cofinite-differences}
 \{E_1(\alpha)-E_1(\alpha_0):
 \alpha\in\Gamma\setminus(E_{\alpha_0}\cup\{\alpha_0\})\}
\end{equation}
is algebraically independent over \(K\).
\end{theorem}

\begin{proof}
Choose \(\lambda_1,\ldots,\lambda_r\) as a basis of the \(\Q\)-vector space
spanned by \(\{\log\alpha:\alpha\in\Gamma\}\); torsion causes no difficulty
because all elements are positive.  For every finite \(A\), all
\(\log\alpha\), \(\alpha\in A\), lie in the \(\Q\)-span of the
\(\lambda_j\).  By \eqref{eq:c-alpha},
\[
 K(\Ein(\alpha):\alpha\in A)
 \subset K(c_\alpha:\alpha\in A,\lambda_1,\ldots,\lambda_r).
\]
The left side has transcendence degree \(|A|\) by
\cref{thm:input-independence}; this proves \eqref{eq:finite-rank-deficit}.

Choose a maximal algebraically independent subfamily of
\(\{c_\alpha:\alpha\in\Gamma\}\).  If its complement contained
\(r+1\) elements, the algebraic dependence of each of them on a finite part
of the maximal family would produce a finite set with deficit at least
\(r+1\), contradicting \eqref{eq:finite-rank-deficit}.  This proves
\eqref{eq:cofinite-family}.

For distinct \(\alpha_1,\ldots,\alpha_m\ne\alpha_0\), the differences
\(\Ein(\alpha_i)-\Ein(\alpha_0)\) are algebraically independent over \(K\):
they are full-rank linear images of the independent coordinates at
\(\alpha_0,\alpha_1,\ldots,\alpha_m\).  Subtracting the logarithmic
differences costs at most the same \(r\) logarithmic coordinates.  The first
argument and the maximality argument therefore apply verbatim.
\end{proof}

The common Euler term costs at most one further transcendence degree.  This
turns the shift bookkeeping into an unconditional theorem about the tails
themselves.

\begin{theorem}[Finite-rank tail defect and cofinite independence]
\label{thm:E1-cofinite}
Let \(\Gamma\subset\Qbar_{>0}^{\times}\) have multiplicative rank \(r\).
For every finite \(A\subset\Gamma\),
\begin{equation}\label{eq:E1-finite-rank-defect}
 \boxed{\displaystyle
 \trdeg_KK(E_1(\alpha):\alpha\in A)\ge |A|-r-1.}
\end{equation}
There is a set \(E\subset\Gamma\), \(|E|\le r+1\), such that
\begin{equation}\label{eq:E1-cofinite}
 \{E_1(\alpha):\alpha\in\Gamma\setminus E\}
\end{equation}
is algebraically independent over \(K\).  In particular at most \(r+1\)
members of the full tail family are algebraic over \(K\), and at most
\(r+1\) members of
\[
 \{\gamma\}\cup\{E_1(\alpha):\alpha\in\Gamma\}
\]
are algebraic over \(K\).  If \(\gamma\in K^{\mathrm{alg}}\), every
occurrence of \(r+1\) in these tail assertions can be replaced by \(r\).
\end{theorem}

\begin{proof}
By \cref{thm:cofinite-independence},
\[
 |A|-r\le\trdeg_KK(c_\alpha:\alpha\in A).
\]
Since \(c_\alpha=\gamma+E_1(\alpha)\), the field on the right is contained
in \(K(\gamma,E_1(\alpha):\alpha\in A)\).  Adjoining \(\gamma\) costs at
most one transcendence degree, which proves
\eqref{eq:E1-finite-rank-defect}.  Choose a maximal algebraically
independent subfamily of the tails.  If its complement contained \(r+2\)
elements, adjoining the finitely many members of the maximal family needed
to algebraize them would contradict \eqref{eq:E1-finite-rank-defect}.
This proves \eqref{eq:E1-cofinite} and the first algebraicity count.

If \(\gamma\) is algebraic over \(K\), adjoining it costs nothing and the
same proof gives defect \(r\).  If it is transcendental, it is not counted
among the algebraic members.  The assertion for the union follows in both
cases.
\end{proof}

\begin{corollary}[The two-prime triangular grid]
\label{cor:two-prime-grid}
For \(n\ge1\), put
\[
 A_n=\{2^k3^\ell:k,\ell\ge0,\ 1\le k+\ell\le n\},
 \qquad |A_n|=\frac{n(n+3)}2.
\]
Among \(\{c_\alpha:\alpha\in A_n\}\), all but at most two can be chosen
jointly algebraically independent over \(K\), and at most two are
algebraic.  Among \(\{E_1(\alpha):\alpha\in A_n\}\), all but at most three
can be chosen jointly algebraically independent, and at most three are
algebraic.  Thus at \(n=20\) the grid has \(230\) entries: at least \(228\)
of the first family, and at least \(227\) of the tail family, belong to
jointly algebraically independent subfamilies.
\end{corollary}

\begin{proof}
The group generated by \(2\) and \(3\) has multiplicative rank two.  Apply
\cref{thm:cofinite-independence,thm:E1-cofinite}; the cardinality is
\(\sum_{j=1}^n(j+1)=n(n+3)/2\).
\end{proof}

\begin{corollary}[Four classical special-function families]
\label{cor:four-classical-cofinite}
Let \(\Gamma\subset\Qbar_{>0}^{\times}\) have multiplicative rank \(r\).
From the combined family
\begin{equation}\label{eq:four-classical-family}
 \{E_1(\alpha),\Ei(\alpha),\Ci(\alpha),
   \operatorname{Si}(\alpha):\alpha\in\Gamma\}
\end{equation}
delete at most \(r+1\) members.  The remaining family is algebraically
independent over \(K\); in particular at most \(r+1\) members of the full
family are algebraic over \(K\).

For \(\alpha=1,\ldots,M\), the \(4M\) displayed values have transcendence
degree at least
\begin{equation}\label{eq:four-classical-integer-count}
 4M-\pi(M)-1
\end{equation}
over \(K\), and at most \(\pi(M)+1\) of them are algebraic over \(K\).
\end{corollary}

\begin{proof}
For each positive \(\alpha\), the four independent coordinates
\[
 \Ein(\alpha),\quad \Ein(-\alpha),\quad
 \frac{\Ein(i\alpha)+\Ein(-i\alpha)}2,\quad
 \frac{\Ein(i\alpha)-\Ein(-i\alpha)}{2i}
\]
are respectively
\[
 \gamma+\log\alpha+E_1(\alpha),\quad
 \gamma+\log\alpha-\Ei(\alpha),\quad
 \gamma+\log\alpha-\Ci(\alpha),\quad
 \operatorname{Si}(\alpha).
\]
Across distinct positive \(\alpha\), the underlying points
\(\{\pm\alpha,\pm i\alpha\}\) are disjoint; hence all these coordinates
are algebraically independent over \(K\) by
\cref{thm:input-independence,lem:linear-images}.  For any finite selection
from \eqref{eq:four-classical-family}, adjoining \(\gamma\) and at most
\(r\) logarithmic generators recovers the same number of independent
coordinates.  The resulting defect is at most \(r+1\), and maximality gives
the cofinite claim.  The integers up to \(M\) generate a multiplicative
group of rank \(\pi(M)\), which proves
\eqref{eq:four-classical-integer-count}.
\end{proof}

The finite-dimensional form identifies exactly where the deficit occurs.
Let \(\alpha_1,\ldots,\alpha_n\) be positive algebraic numbers contained in
a rank-\(r\) group.  Choose logarithmic generators
\(\lambda=(\log g_1,\ldots,\log g_r)^{\mathsf T}\) and an exponent matrix
\(L\in M_{n\times r}(\Q)\) such that
\begin{equation}\label{eq:log-exponent-matrix}
 (\log\alpha_1,\ldots,\log\alpha_n)^{\mathsf T}=L\lambda.
\end{equation}

\begin{theorem}[Linear-image rank deficit]
\label{thm:linear-rank-deficit}
Let \(M\in M_{m\times n}(\Qbar)\) have row rank \(s\), and put
\(t=\operatorname{rank}(ML)\).  With
\(c=(c_{\alpha_1},\ldots,c_{\alpha_n})^{\mathsf T}\), one has
\begin{equation}\label{eq:linear-rank-deficit}
 \boxed{\displaystyle
 \trdeg_KK(Mc)\ge s-t\ge s-r.}
\end{equation}
\end{theorem}

\begin{proof}
Write \(x=(\Ein(\alpha_1),\ldots,\Ein(\alpha_n))^{\mathsf T}\).  Then
\(x=c+L\lambda\), so \(Mx=Mc+ML\lambda\).  The \(s\) independent rows of
\(M\) give \(s\) algebraically independent coordinates of \(Mx\) by
\cref{lem:linear-images}.  The vector \(ML\lambda\) uses at most \(t\)
logarithmic coordinates.  Hence
\[
 s\le \trdeg_KK(Mc)+t,
\]
which is the assertion.
\end{proof}

There is also a sharp statement in the directions where both the Euler
coefficient and all logarithms cancel.

\begin{corollary}[Complete cancellation space]
\label{cor:complete-cancellation-space}
Let
\begin{equation}\label{eq:cancellation-space}
 U=\ker\begin{pmatrix}\mathbf1^{\mathsf T}\\L^{\mathsf T}\end{pmatrix}
 \subset\Qbar^n,
 \qquad
 d=\dim U=n-\operatorname{rank}[\mathbf1,L].
\end{equation}
For any basis \(u_1,\ldots,u_d\) of \(U\), the numbers
\begin{equation}\label{eq:cancellation-values}
 W_j=\sum_{i=1}^n u_{j,i}\Ein(\alpha_i)
     =\sum_{i=1}^n u_{j,i}E_1(\alpha_i)
 \qquad(1\le j\le d)
\end{equation}
are algebraically independent over \(K\).  Thus
\eqref{eq:cancellation-space} is not only the complete space of
gamma-free, log-free coefficient directions: every independent direction
produces a new algebraically independent period.
\end{corollary}

\begin{proof}
The two defining equations for \(U\), together with
\eqref{eq:Ein-connection}, prove the equality in
\eqref{eq:cancellation-values}.  The basis rows have rank \(d\); apply
\cref{lem:linear-images} to the algebraically independent \(\Ein\)-coordinates.
\end{proof}

For a directed graph \(G=(V,E)\), let \(\partial\) be its incidence matrix,
and attach a positive algebraic number \(\alpha_v\) to each vertex.  If
\(L\) is the corresponding exponent matrix, then
\cref{thm:linear-rank-deficit} gives the graph form
\begin{equation}\label{eq:graph-rank-deficit}
 \trdeg_KK\bigl(E_1(\alpha_{t(e)})-E_1(\alpha_{s(e)}):e\in E\bigr)
 \ge |V|-c(G)-\operatorname{rank}(\partial^{\mathsf T}L),
\end{equation}
where \(c(G)\) is the number of connected components.  This is a genuine
incidence-geometric extension of the one-dimensional prime-defect counts
used below.

\section{External realizations of the independence engine}
\label{sec:external-realizations}

We record three consequences that export the algebraic-point input beyond
the exponential-integral language.  They classify relations rather than
asserting new irrationality results for Euler's constant.

\subsection*{Sine-integral graphic matroids}

Let \(G=(V,E)\) be a finite loopless directed multigraph, and choose nonzero algebraic numbers
\(a_v\) such that \(a_u^2\ne a_v^2\) for \(u\ne v\).  Put
\[
 s_v=\operatorname{Si}(a_v)
 =\frac{\Ein(ia_v)-\Ein(-ia_v)}{2i},
 \qquad
 I_e=s_{t(e)}-s_{s(e)}
 =\int_{a_{s(e)}}^{a_{t(e)}}\frac{\sin z}{z}\,dz.
\]

\begin{theorem}[Complete graphic relation ideal]
\label{thm:sine-graphic-ideal}
Let \(\partial_G:K^E\to K^V\) be the incidence map
\(\partial_Ge=e_{t(e)}-e_{s(e)}\).  Then
\begin{equation}\label{eq:sine-graphic-ideal}
 \boxed{\displaystyle
 \ker\!\left(K[X_e:e\in E]\longrightarrow K[I_e:e\in E]\right)
 =
 \left\langle\sum_{e\in E}c_eX_e:
 c\in\ker\partial_G\right\rangle.}
\end{equation}
Consequently
\begin{equation}\label{eq:sine-graphic-trdeg}
 \trdeg_KK(I_e:e\in E)=|V|-c(G),
\end{equation}
and a set of edge integrals is algebraically independent exactly when its
edges form a forest.
\end{theorem}

\begin{proof}
The pairs \(\{ia_v,-ia_v\}\) are disjoint.  The full-rank linear-image
lemma applied to the corresponding algebraically independent \(\Ein\)
values shows that the \(s_v\) are algebraically independent.  The edge
vector is \(I=\partial_G^{\mathsf T}s\).  Hence its coordinate ring is
the symmetric algebra of the incidence row space, whose defining ideal is
generated by the linear cycle relations.  The rank of the incidence matrix
is \(|V|-c(G)\), and a set of incidence columns is independent exactly when
the corresponding edges form a forest in the underlying multigraph.
\end{proof}

This polynomial-relation statement strengthens the disjoint-endpoint
linear-independence results for sine-integral values in
\cite{Delaygue2022} by using the stronger input from the companion paper.

\subsection*{Generalized Dickman log-Laplace values}

For the generalized Dickman distribution of
\cite{GrahovacEtAl2024}, the Laplace transform is
\begin{equation}\label{eq:Dickman-Laplace}
 \psi_{\theta,a}(s)=\exp\{-\theta\Ein(as)\}.
\end{equation}
Take positive algebraic \(\theta_j,a_j,s_j\), set \(x_j=a_js_j\), and
\[
 \Lambda_j=-\log\psi_{\theta_j,a_j}(s_j)=\theta_j\Ein(x_j).
\]

\begin{theorem}[Dickman log-Laplace classification]
\label{thm:Dickman-classification}
For every finite \(J\),
\begin{equation}\label{eq:Dickman-trdeg}
 \boxed{\displaystyle
 \trdeg_KK(\Lambda_j:j\in J)=\#\{x_j:j\in J\}.}
\end{equation}
For each distinct \(x\), choose a representative \(j_x\).  The complete
polynomial relation ideal is
\begin{equation}\label{eq:Dickman-relation-ideal}
 \ker\!\left(K[X_j:j\in J]\to K[\Lambda_j:j\in J]\right)
 =
 \left\langle
 \theta_{j_x}X_j-\theta_jX_{j_x}:
 x_j=x,\ j\ne j_x
 \right\rangle.
\end{equation}
Moreover,
for integers \(m_j\),
\begin{equation}\label{eq:Dickman-multiplicative}
 \boxed{\displaystyle
 \prod_{j\in J}\psi_{\theta_j,a_j}(s_j)^{m_j}=1
 \Longleftrightarrow
 \sum_{j:x_j=x}m_j\theta_j=0
 \quad\text{for every distinct }x.}
\end{equation}
\end{theorem}

\begin{proof}
Choose one representative for each distinct \(x_j\).  The corresponding
\(\Ein(x_j)\) are algebraically independent; every \(\Lambda_j\) is a nonzero
algebraic multiple of its block representative.  This proves the
transcendence degree and the linear defining ideal.  The Laplace values are
positive real numbers, so a multiplicative relation equal to one is
equivalent, after the real logarithm, to a linear relation among the
distinct \(\Ein(x)\); algebraic independence forces the displayed block
conditions.
\end{proof}

The theorem classifies the logarithmic values and the multiplicative
relations equal to one.  It does not claim to classify every polynomial
relation among the Laplace values themselves.

\subsection*{Regularized incomplete-Gamma jets}

For \(\alpha>0\) algebraic, let
\[
 \Gamma_{<}(s,\alpha)=\int_0^\alpha e^{-t}t^{s-1}\,dt,
 \qquad
 \Gamma_{>}(s,\alpha)=\int_\alpha^\infty e^{-t}t^{s-1}\,dt,
\]
initially in their half-planes of convergence and then by continuation.
At \(s=0\), define
\begin{equation}\label{eq:incomplete-gamma-jets}
 R_{\alpha,k}=[s^k]\{s^{-1}-\Gamma_{<}(s,\alpha)\},
 \qquad
 U_{\alpha,k}=[s^k]\Gamma_{>}(s,\alpha)
 \quad(k\ge0).
\end{equation}
This boundary-Laurent and multipoint regime is distinct from the
positive-parameter incomplete-Gamma special values studied, for example,
in \cite{MurtySaha2015}; no priority claim beyond the precise statements
proved below is intended.

\begin{theorem}[Cofinite independence of regularized lower jets]
\label{thm:lower-gamma-jets}
Let \(\Gamma\subset\Qbar_{>0}^{\times}\) have multiplicative rank \(r\),
and let \(S\subset\Gamma\times\N_0\) be finite.  Then
\begin{equation}\label{eq:lower-gamma-jet-defect}
 \boxed{\displaystyle
 \trdeg_KK(R_{\alpha,k}:(\alpha,k)\in S)\ge |S|-r.}
\end{equation}
Consequently, after deleting at most \(r\) members, the entire two-index
family \(\{R_{\alpha,k}:\alpha\in\Gamma,\ k\ge0\}\) is algebraically
independent over \(K\).  In particular at most \(r\) of these values are
algebraic over \(K\).

More explicitly, with \(L_\alpha=\log\alpha\),
\begin{equation}\label{eq:lower-gamma-jet-explicit}
 R_{\alpha,k}=-\frac{L_\alpha^{k+1}}{(k+1)!}
 +\sum_{j=0}^k\frac{(-1)^jL_\alpha^{k-j}}{(k-j)!}
 \Ein_{j+1}(\alpha).
\end{equation}
Thus \(R_{\alpha,0}=c_\alpha\), while
\begin{equation}\label{eq:lower-gamma-at-one}
 R_{1,k}=(-1)^k\Ein_{k+1}(1)
\end{equation}
and the latter all-order family is algebraically independent over \(K\)
without exceptions.
\end{theorem}

\begin{proof}
Splitting \(e^{-t}=1+(e^{-t}-1)\), expanding \(\alpha^s\), and then
expanding \((n+s)^{-1}\) gives \eqref{eq:lower-gamma-jet-explicit}.
For the finite set \(S\), collect the finitely many coordinates
\(\Ein_j(\alpha)\) that occur on the right.  The rows defining the
\(R_{\alpha,k}\) have rank \(|S|\): within each fixed \(\alpha\), the row
of order \(k\) has the new highest coordinate \(\Ein_{k+1}(\alpha)\) with
nonzero coefficient, and different \(\alpha\)-blocks are disjoint.

Let \(X\) be the resulting vector of \(M\) \(\Ein\)-coordinates and let
\(L=K(\lambda_1,\ldots,\lambda_r)\), where the \(\lambda_i\) span the real
logarithms of \(\Gamma\).  Over \(L\), complete the \(|S|\) displayed rows
to an invertible linear coordinate system \((R,Z)\) on \(X\), with
\(|Z|=M-|S|\).  By \eqref{eq:resonant-input},
\[
 M=\trdeg_KK(X)
 \le\trdeg_KL(R,Z)
 \le r+\trdeg_KK(R)+M-|S|.
\]
This proves \eqref{eq:lower-gamma-jet-defect}.  The maximality argument of
\cref{thm:cofinite-independence} then gives the cofinite assertion.
Setting \(L_1=0\) proves \eqref{eq:lower-gamma-at-one}.
\end{proof}

Write \(\Gamma(s)=s^{-1}+\sum_{k\ge0}g_ks^k\).  Since
\(\Gamma=\Gamma_<+\Gamma_>\), one has
\begin{equation}\label{eq:upper-lower-jet-split}
 U_{\alpha,k}=g_k+R_{\alpha,k}.
\end{equation}

\begin{corollary}[Upper-jet defect]
\label{cor:upper-gamma-jets}
For finite \(S\subset\Gamma\times\N_0\), let
\(d(S)=\#\{k:(\alpha,k)\in S\}\).  Then
\begin{equation}\label{eq:upper-gamma-jet-defect}
 \boxed{\displaystyle
 \trdeg_KK(U_{\alpha,k}:(\alpha,k)\in S)
 \ge |S|-r-d(S).}
\end{equation}
For every fixed \(k\), all but at most \(r+1\) members of
\(\{U_{\alpha,k}:\alpha\in\Gamma\}\) form an algebraically independent
family over \(K\).  At \(k=0\), \(U_{\alpha,0}=E_1(\alpha)\), so this
recovers \cref{thm:E1-cofinite}.
\end{corollary}

\begin{proof}
By \eqref{eq:upper-lower-jet-split}, adjoining the \(d(S)\) Gamma
coefficients \(g_k\) to the upper jets generates all lower jets in \(S\).
Combine \eqref{eq:lower-gamma-jet-defect} with the transcendence-degree
inequality.  For fixed \(k\), apply the same maximality argument as before.
\end{proof}

The finite-order shift ledger can now be replaced by a complete invariant
ring.  This is a structural classification, not an irrationality theorem.

\begin{proposition}[Complete finite Gamma-jet quotient]
\label{prop:finite-gamma-jet-quotient}
Fix \(N\ge0\), let
\[
 A_N(s)=1+\sum_{k=0}^Ng_ks^{k+1}\pmod {s^{N+2}},
 \qquad T_k=U_{1,k},\qquad d_k=g_k-T_k,
\]
and put \(\kappa_j=[s^j]\log A_N(s)\).  On the polynomial ring in the
\(g_k,T_k\), define a \(\mathbb G_a\)-action by
\begin{equation}\label{eq:gamma-jet-shift}
 A_N(s)\longmapsto e^{-ts}A_N(s)\pmod {s^{N+2}},
 \qquad d_k\longmapsto d_k.
\end{equation}
Then
\begin{equation}\label{eq:gamma-jet-invariant-ring}
 \boxed{\displaystyle
 \Q[g_0,\ldots,g_N,T_0,\ldots,T_N]^{\mathbb G_a}
 =\Q[\kappa_2,\ldots,\kappa_{N+1},d_0,\ldots,d_N].}
\end{equation}
The analogous equality holds for fraction fields.  At the actual Gamma
split,
\begin{equation}\label{eq:actual-gamma-jet-invariants}
 \kappa_1=-\gamma,qquad
 \kappa_j=\frac{(-1)^j}{j}\zeta(j)\ (j\ge2),qquad
 d_k=(-1)^{k+1}\Ein_{k+1}(1).
\end{equation}
Thus \(\gamma\) is the unique orbit coordinate at every finite order; the
zeta cumulants and resonant \(\Ein\)-coordinates generate all polynomial
and rational invariants.
\end{proposition}

\begin{proof}
The change of variables
\[
 (g_0,\ldots,g_N,T_0,\ldots,T_N)
 \longleftrightarrow
 (\kappa_1,\ldots,\kappa_{N+1},d_0,\ldots,d_N)
\]
is triangular and polynomial in both directions, by the truncated formal
identities \(\log A_N\) and \(A_N=\exp(\log A_N)\).  Under
\eqref{eq:gamma-jet-shift}, \(\kappa_1\mapsto\kappa_1-t\) and every other
displayed coordinate is fixed.  The invariant ring of translation in one
variable is generated by the remaining variables, proving
\eqref{eq:gamma-jet-invariant-ring}.  Finally
\(A_N(s)\) is the truncation of
\(s\Gamma(s)=\Gamma(1+s)\), so
\[
 \log\Gamma(1+s)=-\gamma s+
 \sum_{j\ge2}\frac{(-1)^j}{j}\zeta(j)s^j.
\]
Equation \eqref{eq:upper-lower-jet-split} at \(\alpha=1\), together with
\eqref{eq:lower-gamma-at-one}, gives the formula for \(d_k\).
\end{proof}

\begin{remark}[Correct reading of the all-order ledger]
The gauge orbit has dimension exactly one at every order.  Any additional
unknown excess in a truncated equation count is not another copy of the
Euler shift: the new odd-zeta cumulants are themselves invariant arithmetic
coordinates.  Consequently a growing unknown-minus-equation count is useful
bookkeeping, but is not by itself a no-go theorem.
\end{remark}

The corresponding export to arbitrary characteristic-zero algebraic
matroids uses an external linearization theorem and is therefore separated
from these direct applications; see \cref{app:Ein-matroids}.

\begin{remark}[A conditional density consequence of Powers]
\label{rem:Powers-conditional}
In the notation of Powers' later multisection manuscript
\cite[\S5]{Powers2025}, assume that
\[
 G_1(1),H_1(1),\ldots,G_N(1),H_N(1)
\]
are algebraically independent over \(\Qbar(e)\) for every \(N\).  This
assertion appears in the concluding multisection discussion rather than as
a proved numbered theorem in the cited version.  There
\begin{align*}
 G_n(1)&=\frac{\widetilde\delta^{(n)}-\delta^{(n)}}{2e},\\
 H_n(1)&=\gamma^{(n)}+
 \frac{\widetilde\delta^{(n)}+\delta^{(n)}}{2e}.
\end{align*}
Conditional on it, combine the two generalized Euler--Gompertz rows
\[
 \delta^{(1)},\ldots,\delta^{(N)},\qquad
 \widetilde\delta^{(1)},\ldots,\widetilde\delta^{(N)}.
\]
Their field satisfies
\[
 \trdeg_{\Qbar}\Qbar\bigl(
 \delta^{(1)},\ldots,\delta^{(N)},
 \widetilde\delta^{(1)},\ldots,\widetilde\delta^{(N)}\bigr)
 \ge\left\lfloor\frac{3N}{2}\right\rfloor-1.
\]
Indeed the Gamma-moment identities quoted by Powers give
\[
 \trdeg_{\Qbar}\Qbar(\gamma^{(1)},\ldots,\gamma^{(N)})
 \le N-\left\lfloor\frac N2\right\rfloor+1.
\]
The assumed \(2N\) independent coordinates, together with the
transcendence of \(e\), lie in the field generated by \(e\), the two
Euler--Gompertz rows, and these Gamma moments.  Subtracting transcendence
degrees yields the displayed bound.  Greedy basis selection in the
associated algebraic matroid, ordered as
\(\delta^{(1)},\widetilde\delta^{(1)},\delta^{(2)},\ldots\), then gives a
jointly algebraically independent subfamily of lower density at least
\(3/4\).  This entire paragraph is conditional, and no unconditional
result elsewhere in the paper depends on it.
\end{remark}

\section{Universal level-trace polynomials}
\label{sec:trace-polynomials}

For \(c\in\C^\times\), set
\[
 F_c(z)=\Ein(z)-c,
\]
and denote by \(Z_c\) its zero multiset, with multiplicities.  Since \(F_c\)
has order one, its inverse-power sums converge absolutely from order two
onward.

\begin{theorem}[Universal level-trace polynomial]
\label{thm:universal-trace}
For \(m\ge2\), define
\begin{equation}\label{eq:Pm-def}
 P_m(X):=-m[z^m]\log\bigl(1-X\Ein(z)\bigr).
\end{equation}
Then \(P_m\in\Q[X]\), its constant term is zero, and it is monic of degree
\(m\).  Moreover
\begin{equation}\label{eq:universal-trace}
 \boxed{\displaystyle
 T_m(c):=\sum_{\rho\in Z_c}\rho^{-m}=P_m(c^{-1}).}
\end{equation}
In particular
\begin{align}
 P_2(X)&=X^2-\frac X2,\label{eq:P2}\\
 P_3(X)&=X^3-\frac34X^2+\frac16X,\label{eq:P3}\\
 P_4(X)&=X^4-X^3+\frac{25}{72}X^2-\frac1{24}X.\label{eq:P4}
\end{align}
\end{theorem}

\begin{proof}
A genus-one Hadamard factorization has the form
\[
 F_c(z)=-c\,e^{a_cz}
 \prod_{\rho\in Z_c}
 \left(1-\frac z\rho\right)e^{z/\rho}.
\]
Taking the logarithm at the origin gives, for every \(m\ge2\),
\[
 [z^m]\log\frac{F_c(z)}{F_c(0)}
 =-\frac1m\sum_{\rho\in Z_c}\rho^{-m}.
\]
On the other hand,
\[
 \frac{F_c(z)}{F_c(0)}=1-\frac{\Ein(z)}c.
\]
Comparison proves \eqref{eq:universal-trace}.  Formula
\eqref{eq:Pm-def} is a polynomial in \(X\) of degree at most \(m\).  Since
\(\Ein(z)=z+O(z^2)\), the contribution of
\(X^m\Ein(z)^m/m\) to \eqref{eq:Pm-def} is \(X^m\).  Thus \(P_m\) is monic
of degree \(m\).  The displayed formulas follow by expanding through degree
four.
\end{proof}

\begin{remark}[Multiplicity and genus]
No simplicity assertion is required: all zeros are counted with their
multiplicities.  The first inverse sum is not used.  At order one the
canonical genus-one exponential factor absorbs the non-absolutely convergent
sum, whereas \eqref{eq:universal-trace} is absolute for \(m\ge2\).
\end{remark}

\begin{corollary}[Transversal trace independence]
\label{cor:transversal}
Let \(F\subset\C\) be a field of characteristic zero, and let
\(\{c_j:j\in J\}\subset\C^\times\) be algebraically independent over \(F\).
Choose one integer \(m_j\ge2\) for each \(j\).  Then
\[
 \{T_{m_j}(c_j):j\in J\}
\]
is algebraically independent over \(F\), in the finite-subset sense.
\end{corollary}

\begin{proof}
Fix a finite set \(I\subset J\), and write
\(x_j=c_j^{-1}\) and \(t_j=P_{m_j}(x_j)\).  Since
\[
 P_{m_j}(x_j)-t_j=0,
\]
the field \(F(x_j:j\in I)\) is algebraic over \(F(t_j:j\in I)\).  Hence
\[
 \trdeg_FF(t_j:j\in I)
 =\trdeg_FF(x_j:j\in I)=|I|.
\]
\end{proof}

\begin{theorem}[Exact even-trace recovery and ring closure]
\label{cor:fixed-level-closure}
For every \(c\ne0\),
\begin{equation}\label{eq:even-trace-recovery}
 \boxed{\displaystyle
 c^{-1}=144T_4(c)-144T_2(c)^2-14T_2(c).}
\end{equation}
Consequently, for every characteristic-zero field \(F\subset\C\),
\[
 \boxed{\displaystyle
 F[T_m(c):m\ge2]
 =F[T_2(c),T_4(c)]
 =F[c^{-1}]}
\]
and
\[
 \boxed{F(T_2(c),T_4(c))=F(c).}
\]
\end{theorem}

\begin{proof}
With \(x=c^{-1}\), \eqref{eq:P2} and \eqref{eq:P4} give
\[
 144P_4(x)-144P_2(x)^2-14P_2(x)=x.
\]
This proves \eqref{eq:even-trace-recovery}.  Every trace is a polynomial in
\(x\) by \cref{thm:universal-trace}, while the displayed identity recovers
\(x\) from \(T_2,T_4\).  Taking fraction fields and using \(F(x)=F(c)\)
finishes the proof.
\end{proof}

The preceding recovery identity also gives the exact, rather than merely
heuristic, form of the finite-stage shift obstruction.  It is important to
distinguish this gauge action, which fixes every connection coordinate,
from common translation of the connection coordinates themselves;
augmentation zero pertains to the latter action, not to the former.

\begin{theorem}[Exact finite-stage shift quotient]
\label{thm:exact-shift-quotient}
Let \(B\) be a characteristic-zero field, put
\[
 R=B[G,H_1,\ldots,H_n],\qquad c_i=G+H_i,
\]
and let \(\mathbb G_a\) act by
\begin{equation}\label{eq:gauge-action}
 G\longmapsto G+t,\qquad H_i\longmapsto H_i-t.
\end{equation}
Then
\begin{equation}\label{eq:shift-invariant-ring}
 \boxed{R^{\mathbb G_a}=B[c_1,\ldots,c_n],\qquad
 \operatorname{Frac}(R)^{\mathbb G_a}=B(c_1,\ldots,c_n).}
\end{equation}
If the universal traces are formed at these levels, then
\begin{align}
 B\bigl(T_m(c_i):1\le i\le n,\ m\ge2\bigr)
 &=B(c_1,\ldots,c_n),\label{eq:trace-invariant-field}\\
 B\bigl[T_m(c_i):1\le i\le n,\ m\ge2\bigr]
 &=B[c_1^{-1},\ldots,c_n^{-1}].\label{eq:trace-invariant-ring}
\end{align}
Thus the complete finite multi-level trace field is exactly the rational
quotient by the gauge direction.  It loses one orbit coordinate and no
others.
\end{theorem}

\begin{proof}
The polynomial change of variables \(H_i=c_i-G\) identifies
\(R=B[c_1,\ldots,c_n,G]\), and \eqref{eq:gauge-action} becomes translation
of \(G\) with every \(c_i\) fixed.  In characteristic zero the polynomial
and rational invariants of translation in one variable are respectively
\(B[c_1,\ldots,c_n]\) and \(B(c_1,\ldots,c_n)\), proving
\eqref{eq:shift-invariant-ring}.  Every trace is a polynomial in
\(c_i^{-1}\), while \eqref{eq:even-trace-recovery} recovers \(c_i^{-1}\)
from \(T_2(c_i),T_4(c_i)\).  This proves
\eqref{eq:trace-invariant-field}--\eqref{eq:trace-invariant-ring}.
\end{proof}

\begin{remark}[What the formal quotient does and does not prove]
The elimination ideal of the equations \(c_i=G+H_i\) in \(B[G]\) is zero:
the same finite invariant data admit a model with \(G=0\) and another with
\(G\) transcendental.  Hence invariant algebraic identities alone cannot
decide the arithmetic of the orbit coordinate.  This is a statement about
the formal axiom system, not about the actual complex number \(\gamma\).
Indeed the analytic completion supplies extra boundary data.  If
\(y_N\to g\) in a Hausdorff topological field and a continuous action fixes
every \(y_N\), it fixes \(g\).  Applied to \(Y_N\to\gamma\), and to the two
carrier limits \(\tau_N\to S_2\), \(\sigma_N\to S_4\), this is precisely
where the finite-stage gauge symmetry breaks; \(S_2,S_4\) recover
\(\gamma\) by \eqref{eq:gamma-even-recovery}.
\end{remark}

\begin{corollary}[Arithmetic classification of every Euler trace]
\label{cor:Euler-trace-arithmetic}
For \(m\ge2\), put \(S_m=T_m(\gamma)\).  Then
\[
 S_m\in\Qbar\quad\Longleftrightarrow\quad\gamma\in\Qbar,
 \qquad
 S_m\text{ is transcendental}
 \quad\Longleftrightarrow\quad\gamma\text{ is transcendental}.
\]
Moreover the two even traces recover Euler's constant without exceptions:
\begin{equation}\label{eq:gamma-even-recovery}
 \boxed{\displaystyle
 \gamma=\frac1{144S_4-144S_2^2-14S_2}.}
\end{equation}
In particular
\[
 \gamma\in\Q\quad\Longleftrightarrow\quad S_2,S_4\in\Q.
\]
\end{corollary}

\begin{proof}
If \(S_m=P_m(\gamma^{-1})\) is algebraic, then \(\gamma^{-1}\) is a root
of the nonconstant polynomial \(P_m(X)-S_m\) over \(\Qbar\), and hence is
algebraic.  The converse is immediate.  Formula
\eqref{eq:gamma-even-recovery} is \eqref{eq:even-trace-recovery} at
\(c=\gamma\); it also proves the rationality equivalence.
\end{proof}

\begin{theorem}[The unique two-trace level collision]
\label{thm:T2T3-collision}
For \(c,d\in\C^\times\),
\[
 (T_2(c),T_3(c))=(T_2(d),T_3(d))
\]
holds if and only if either \(c=d\), or
\begin{equation}\label{eq:T2T3-exceptional-pair}
 \boxed{\{c,d\}=\{6-2\sqrt3,6+2\sqrt3\}.}
\end{equation}
In the exceptional case the common pair is
\begin{equation}\label{eq:T2T3-exceptional-value}
 \boxed{(T_2,T_3)=\left(-\frac1{24},\frac1{96}\right).}
\end{equation}
\end{theorem}

\begin{proof}
Put \(x=c^{-1}\) and \(y=d^{-1}\).  From \eqref{eq:P2}, equality of the
second traces gives
\[
 (x-y)\left(x+y-\frac12\right)=0.
\]
If \(x=y\), then \(c=d\).  Otherwise \(x+y=1/2\), and equality of the
third traces, using \eqref{eq:P3}, reduces to
\[
 x^2+xy+y^2-\frac34(x+y)+\frac16=0,
\]
or \(xy=1/24\).  Thus \(x,y\) are the roots of
\[
 X^2-\frac12X+\frac1{24},
\]
and taking reciprocals gives \eqref{eq:T2T3-exceptional-pair}.  Direct
substitution in \eqref{eq:P2} and \eqref{eq:P3} gives
\eqref{eq:T2T3-exceptional-value}.
\end{proof}

\begin{corollary}[Exact vanishing traces with transcendental divisors]
\label{cor:vanishing-trace-levels}
For every \(m\ge2\), define
\[
 \mathcal V_m=
 \{c\in\C^\times:T_m(c)=0\}.
\]
Then
\begin{equation}\label{eq:vanishing-trace-levels}
 \boxed{\mathcal V_m=
 \{x^{-1}:x\ne0,\ P_m(x)=0\}\subset\Qbar^\times.}
\end{equation}
This set is nonempty and finite; its elements arise from exactly \(m-1\)
nonzero roots of \(P_m\), counted with multiplicity.  For every
\(c\in\mathcal V_m\), every zero \(\rho\) of \(\Ein(z)-c\) is
transcendental, although the absolutely convergent inverse trace satisfies
\[
 \sum_{\Ein(\rho)=c}\rho^{-m}=0.
\]
For example, \(\mathcal V_2=\{2\}\).
\end{corollary}

\begin{proof}
The trace identity \eqref{eq:universal-trace} gives
\eqref{eq:vanishing-trace-levels}.  The polynomial \(P_m\) has degree
\(m\), constant term zero, and linear coefficient
\[
 [X]P_m(X)=m[z^m]\Ein(z)=\frac{(-1)^{m-1}}{m!}\ne0.
\]
Hence zero is a simple root and the remaining \(m-1\) roots are nonzero,
counted with multiplicity.  They are algebraic because \(P_m\in\Q[X]\).

Now fix \(c\in\mathcal V_m\).  If a zero \(\rho\) of \(\Ein(z)-c\) were
algebraic, then \(\rho\ne0\) and
\(\Ein(\rho)=c\in\Qbar\), contradicting
\cref{thm:input-independence}.  Thus every point of the divisor is
transcendental.  Finally \(P_2(X)=X(X-1/2)\) gives the example.
\end{proof}

Thus even all inverse traces at one fixed level still contain only one
coordinate.  Moving the level is the only way to restore a transversal
algebraically independent family.

\section{Balanced dyadic levels and all-order moving traces}
\label{sec:dyadic}

For \(N\ge1\), define
\begin{equation}\label{eq:YN}
 Y_N=(N+1)\Ein(2^N)-N\Ein(2^{N+1}).
\end{equation}
The companion paper proved the first assertion below.  The remaining
assertions are the analytic bridge to the Euler level.

\begin{proposition}[Balanced dyadic approach]
\label{prop:dyadic-approach}
The family \(\{Y_N:N\ge1\}\) is countably algebraically independent over
\(K\).  Moreover
\begin{equation}\label{eq:epsilonN}
 Y_N=\gamma+\varepsilon_N,
 \qquad
 \varepsilon_N=(N+1)E_1(2^N)-NE_1(2^{N+1}),
\end{equation}
and
\begin{equation}\label{eq:epsilon-properties}
 \varepsilon_N>0,\qquad
 \varepsilon_N\downarrow0,\qquad
 \varepsilon_N\sim\frac{N+1}{2^N}e^{-2^N}.
\end{equation}
In particular \(Y_N\downarrow\gamma\).
\end{proposition}

\begin{proof}
The algebraic independence is \cite[Lem.~11.2]{KEIO2026I}; it also follows directly from
\cref{thm:input-independence,lem:linear-images}, because the finite row family
\(\bigl((N+1)e_N-Ne_{N+1}\bigr)\) has full row rank.

Substitution of \eqref{eq:Ein-connection} in \eqref{eq:YN} cancels both the
logarithmic and the \(N\log2\) terms and gives \eqref{eq:epsilonN}.  Positivity
follows from
\[
 \varepsilon_N
 =E_1(2^N)+N\{E_1(2^N)-E_1(2^{N+1})\}>0.
\]
For strict decrease put \(g(t)=E_1(2^t)\).  Since
\[
 g'(t)=-(\log2)e^{-2^t},\qquad
 g''(t)=(\log2)^2 2^t e^{-2^t}>0,
\]
strict convexity gives
\[
 \varepsilon_N-\varepsilon_{N+1}
 =(N+1)\{g(N)-2g(N+1)+g(N+2)\}>0.
\]
Finally, the standard positive-axis expansion
\begin{equation}\label{eq:E1-asymptotic}
 E_1(x)=\frac{e^{-x}}x\left(1+O(x^{-1})\right)
 \qquad(x\to+\infty)
\end{equation}
gives the asymptotic in \eqref{eq:epsilon-properties}; see
\cite[\S6.12\(i\)]{DLMF}.
\end{proof}

For \(m\ge2\), let
\begin{equation}\label{eq:tau-mN}
 \tau_{m,N}:=\sum_{\Ein(\rho)=Y_N}\rho^{-m},
\end{equation}
where the solutions are counted with multiplicity.

\begin{theorem}[All-order transversal tower]
\label{thm:all-order-tower}
For every \(m\ge2\) and \(N\ge1\),
\begin{equation}\label{eq:tau-Pm}
 \tau_{m,N}=P_m(Y_N^{-1}).
\end{equation}
For any selection \(m_N\ge2\), the family
\[
 \{\tau_{m_N,N}:N\ge1\}
\]
is countably algebraically independent over \(K\).  For each fixed
\(m\ge2\),
\begin{equation}\label{eq:tau-limit}
 \tau_{m,N}\longrightarrow S_m:=P_m(\gamma^{-1}).
\end{equation}
\end{theorem}

\begin{proof}
Equation \eqref{eq:tau-Pm} is \cref{thm:universal-trace}.  Apply
\cref{cor:transversal} to the independent levels \(Y_N\).  The limit follows
from \cref{prop:dyadic-approach} and continuity of \(P_m(1/x)\) at
\(x=\gamma\).
\end{proof}

\begin{remark}[Why transversal is sharp]
For a fixed \(N\), the even traces already recover the level exactly:
\begin{equation}\label{eq:YN-even-recovery}
 \boxed{\displaystyle
 Y_N^{-1}=144\tau_{4,N}-144\tau_{2,N}^2-14\tau_{2,N}.}
\end{equation}
Thus
\[
 K[\tau_{m,N}:m\ge2]
 =K[\tau_{2,N},\tau_{4,N}]
 =K[Y_N^{-1}].
\]
The full double family is therefore not algebraically independent: the
theorem gives exactly one independent spectral coordinate per moving level.
\end{remark}

The order-two row has additional positivity.  Put
\begin{equation}\label{eq:f-def}
 f(x)=\frac1{x^2}-\frac1{2x}=\frac{2-x}{2x^2},
 \qquad \tau_N:=\tau_{2,N}=f(Y_N).
\end{equation}

\begin{theorem}[Exact positive convergence at order two]
\label{thm:tau2-convergence}
One has
\[
 0<\tau_1<\tau_2<\cdots<S_2,
 \qquad \tau_N\longrightarrow S_2.
\]
More precisely,
\begin{equation}\label{eq:S2-tail-exact}
 \boxed{\displaystyle
 S_2-\tau_N=
 \frac{\varepsilon_N
 \{\gamma(4-\gamma)+(2-\gamma)\varepsilon_N\}}
 {2\gamma^2(\gamma+\varepsilon_N)^2}.}
\end{equation}
Consequently
\begin{equation}\label{eq:S2-tail-asymptotic}
 S_2-\tau_N\sim
 \frac{4-\gamma}{2\gamma^3}
 \frac{N+1}{2^N}e^{-2^N}.
\end{equation}
\end{theorem}

\begin{proof}
Since \(E_1(2)<e^{-2}/2\) and \(\gamma<1\),
\(Y_1=\gamma+2E_1(2)-E_1(4)<\gamma+e^{-2}<2\), so all \(Y_N\) lie in
\((0,2)\).  But
\[
 f'(x)=\frac{x-4}{2x^3}<0\qquad(0<x<2).
\]
Thus \(Y_N\downarrow\gamma\) implies \(\tau_N\uparrow f(\gamma)=S_2\), and
positivity follows from \(f(x)>0\) on \((0,2)\).  Direct subtraction with
\(Y_N=\gamma+\varepsilon_N\) gives \eqref{eq:S2-tail-exact}.  Combine it
with \eqref{eq:epsilon-properties} to obtain
\eqref{eq:S2-tail-asymptotic}.
\end{proof}

\begin{theorem}[Non-P-recursiveness of the dyadic moving sequences]
\label{thm:non-p-recursive}
Neither \((Y_N)_{N\ge1}\) nor \((\tau_N)_{N\ge1}\) is P-recursive over
\(\C\).  Equivalently, neither sequence satisfies a nonzero recurrence
\[
 \sum_{j=0}^r p_j(N)u_{N+j}=0
 \qquad\bigl(p_j\in\C[N]\bigr)
\]
for all sufficiently large \(N\).
\end{theorem}

\begin{proof}
Suppose first that such a recurrence holds for \(Y_N=\gamma+\varepsilon_N\).
Since every polynomial multiple of \(\varepsilon_{N+j}\) tends to zero, we
obtain
\[
 \gamma\sum_{j=0}^rp_j(N)\longrightarrow0.
\]
The sum is a polynomial, hence it vanishes identically.  Let \(j_0\) be the
least index for which \(p_{j_0}\ne0\).  Dividing the remaining relation by
\(p_{j_0}(N)\varepsilon_{N+j_0}\) gives
\[
 1+\sum_{j>j_0}
 \frac{p_j(N)}{p_{j_0}(N)}
 \frac{\varepsilon_{N+j}}{\varepsilon_{N+j_0}}=0.
\]
But \eqref{eq:epsilon-properties} implies, for every fixed \(j>j_0\),
\[
 \frac{\varepsilon_{N+j}}{\varepsilon_{N+j_0}}
 =\exp\{-\bigl(2^j-2^{j_0}\bigr)2^N+O(N)\},
\]
which dominates every rational function of \(N\).  The displayed equality
therefore tends to \(1=0\), a contradiction.

For \(\tau_N=f(Y_N)\), \cref{thm:tau2-convergence} gives
\[
 \tau_N=S_2+f'(\gamma)\varepsilon_N+O(\varepsilon_N^2),
 \qquad f'(\gamma)\ne0.
\]
The same least-index argument applies verbatim.
\end{proof}

\begin{remark}[Scope of the non-P-recursiveness result]
The theorem does not say that no holonomic construction can involve the
finite-stage values.  It says that the natural moving index itself cannot be
packaged as a fixed P-recursive sequence.  Any application of arithmetic
holonomy must therefore replace the moving tower by a genuinely different,
fixed differential or recurrence system rather than merely reindexing
\(Y_N\).
\end{remark}

For orientation, the first values are shown in \cref{tab:dyadic-values}.
The table is illustrative only; no numerical value is used in a proof.

\begin{table}[ht]
\centering
\caption{The dyadic level and its order-two zero trace.}
\label{tab:dyadic-values}
\begin{tabular}{@{}rcc@{}}
\toprule
\(N\)&\(Y_N\)&\(\tau_N\)\\
\midrule
1&0.671237333907806&1.474569389580200\\
2&0.588478390885392&2.037963940792663\\
3&0.577366307471447&2.133831985200994\\
4&0.577215698103968&2.135171685377284\\
5&0.577215664901535&2.135171980841625\\
\midrule
limit&0.577215664901533&2.135171980841646\\
\bottomrule
\end{tabular}
\end{table}

\section{Positive spectral carriers for
\texorpdfstring{\(S_2\) and \(S_4\)}{S2 and S4}}
\label{sec:spectral}

Define
\begin{equation}\label{eq:dN-def}
 d_0=\tau_1,
 \qquad
 d_N=\tau_{N+1}-\tau_N\quad(N\ge1).
\end{equation}

\begin{theorem}[Positive algebraically independent increments]
\label{thm:increments}
The sequence \(\{d_N:N\ge0\}\) is positive and countably algebraically
independent over \(K\).  It satisfies
\begin{equation}\label{eq:S2-sum}
 \boxed{\displaystyle S_2=\sum_{N=0}^\infty d_N}
\end{equation}
and, for \(N\ge1\),
\begin{equation}\label{eq:dN-asymptotic}
 d_N\sim
 \frac{4-\gamma}{2\gamma^3}
 \frac{N+1}{2^N}e^{-2^N}.
\end{equation}
\end{theorem}

\begin{proof}
Positivity follows from \cref{thm:tau2-convergence}.  The finite
transformation
\[
 (\tau_1,\ldots,\tau_{M+1})
 \longleftrightarrow(d_0,\ldots,d_M)
\]
is an invertible triangular \(\Q\)-linear transformation.  The
\(\tau_N\) are algebraically independent by
\cref{thm:all-order-tower}; hence so are the \(d_N\).  The telescoping
identity
\[
 \sum_{N=0}^Md_N=\tau_{M+1}
\]
and \cref{thm:tau2-convergence} give \eqref{eq:S2-sum}.  Finally the ratio
of two successive tails in \eqref{eq:S2-tail-asymptotic} tends to zero, so
\[
 d_N=(S_2-\tau_N)-(S_2-\tau_{N+1})\sim S_2-\tau_N,
\]
which proves \eqref{eq:dN-asymptotic}.
\end{proof}

Let \(\{e_N:N\ge0\}\) be the standard basis of \(\ell^2(\N_0)\).

\begin{corollary}[A positive operator carrier]
\label{cor:operator}
The diagonal operator
\[
 \mathsf A e_N=d_Ne_N
\]
is positive, self-adjoint and trace class, with
\[
 \Tr\mathsf A=S_2.
\]
More strongly,
\[
 \sum_{N\ge0}d_N^p<\infty\qquad(p>0).
\]
Thus \(\mathsf A\in\mathcal S_p\) for every \(p\ge1\), and satisfies the
corresponding quasi-Schatten condition for \(0<p<1\).  Its nonzero
eigenvalues are simple and countably algebraically independent over \(K\).
\end{corollary}

\begin{proof}
All assertions except simplicity follow immediately from
\cref{thm:increments}.  The double-exponential asymptotic
\eqref{eq:dN-asymptotic} gives summability of every positive power.  If
\(d_j=d_k\) for \(j\ne k\), the algebraically independent family would
satisfy a nonzero linear relation; hence the eigenvalues are distinct.
\end{proof}

The Fredholm determinant of the operator gives a scalar entire realization.

\begin{theorem}[Order-zero Laguerre--P\'olya realization]
\label{thm:Delta2}
The product
\begin{equation}\label{eq:Delta2}
 \Delta_2(z)=
 \det\!\left(I-\frac{z^2}{2}\mathsf A\right)
 =\prod_{N=0}^\infty\left(1-\frac{d_N}{2}z^2\right)
\end{equation}
converges locally uniformly and defines an order-zero entire function in the
Laguerre--P\'olya class.  Its zeros are simple and equal to
\[
 \pm r_N,\qquad r_N=\sqrt{\frac2{d_N}}.
\]
The positive half \(\{r_N:N\ge0\}\) is countably algebraically independent
over \(K\), and
\begin{equation}\label{eq:Delta2-trace}
 \boxed{\displaystyle
 \sum_{\rho\in Z(\Delta_2)}\rho^{-2}=S_2.}
\end{equation}
\end{theorem}

\begin{proof}
Since \(\sum d_N<\infty\), the product in \eqref{eq:Delta2} converges
normally on compact sets.  Every finite partial product has only real zeros,
so their locally uniform limit belongs to the Laguerre--P\'olya class; see
\cite[Chapter~VIII]{Levin1996}.  For every \(\eta>0\),
\[
 \sum_{N\ge0}r_N^{-\eta}
 =2^{-\eta/2}\sum_{N\ge0}d_N^{\eta/2}<\infty.
\]
Thus the zero exponent of convergence, and hence the order of the canonical
product, is zero.  Simplicity follows from the distinctness of the \(d_N\).

For a finite index set \(I\), the relations \(d_N=2/r_N^2\) imply
\[
 \trdeg_KK(r_N:N\in I)
 \ge \trdeg_KK(d_N:N\in I)=|I|,
\]
which proves algebraic independence of the positive zeros.  Finally,
absolute convergence and \eqref{eq:S2-sum} give
\[
 \sum_{\rho\in Z(\Delta_2)}\rho^{-2}
 =\sum_{N\ge0}\left(r_N^{-2}+(-r_N)^{-2}\right)
 =\sum_{N\ge0}d_N=S_2.
\]
\end{proof}

\begin{theorem}[A second positive carrier and even-trace recovery]
\label{thm:positive-S4-carrier}
Put
\[
 \sigma_N=T_4(Y_N)=P_4(Y_N^{-1})\quad(N\ge1),
\]
and define
\[
 d_0^{(4)}=\sigma_1,\qquad
 d_N^{(4)}=\sigma_{N+1}-\sigma_N\quad(N\ge1).
\]
Then every \(d_N^{(4)}\) is positive, the family
\(\{d_N^{(4)}:N\ge0\}\) is algebraically independent over \(K\), and
\begin{equation}\label{eq:S4-positive-sum}
 S_4=\sum_{N\ge0}d_N^{(4)}.
\end{equation}
Consequently the diagonal operator
\(\mathsf A_4=\diag(d_N^{(4)}:N\ge0)\) belongs to every Schatten class and
\(\Tr\mathsf A_4=S_4\).  Together with the order-two carrier
\(\mathsf A_2:=\mathsf A\) of \cref{cor:operator}, it gives
\begin{equation}\label{eq:gamma-two-positive-carriers}
 \boxed{\displaystyle
 \gamma=
 \left\{144\Tr\mathsf A_4
 -144(\Tr\mathsf A_2)^2-14\Tr\mathsf A_2\right\}^{-1}.}
\end{equation}
\end{theorem}

\begin{proof}
The earlier bound \(Y_N<2\) and the monotonicity \(Y_N\downarrow\gamma\)
show that \(x_N=Y_N^{-1}\) increases and satisfies \(x_N>1/2\).  From
\eqref{eq:P4},
\[
 P_4(1/2)=\frac1{288},\qquad P_4'(1/2)=\frac1{18},
\]
and
\[
 P_4''(x)=12x^2-6x+\frac{25}{36}>0\qquad(x\ge1/2).
\]
Thus \(P_4\) is positive and strictly increasing on this range, so all
increments are positive.  The \(\sigma_N\) are algebraically
independent by \cref{thm:all-order-tower}; the triangular change from
\((\sigma_1,\ldots,\sigma_{M+1})\) to
\((d_0^{(4)},\ldots,d_M^{(4)})\) preserves this property.  Telescoping and
\(T_4(Y_N)\to T_4(\gamma)=S_4\) prove
\eqref{eq:S4-positive-sum}.  The doubly exponential convergence of \(Y_N\)
gives summability of every positive power of \(d_N^{(4)}\).  Finally apply
\eqref{eq:gamma-even-recovery}.
\end{proof}

\begin{remark}[Three distinct meanings of ``gamma-free'']
Each finite datum \(Y_N,\tau_N,d_N\) is defined by algebraic-point
\(\Ein\)-values without displaying \(\gamma\).  The limit and sum
\[
 S_2=\lim_N\tau_N=\sum_Nd_N
\]
are countable constructions.  They do not prove that \(S_2\) belongs to the
finite field generated by algebraic-point \(\Ein\)-values.  Indeed
\([z^2]\Delta_2(z)=-S_2/2\), so its Taylor coefficients are already beyond
what the finite-stage algebraic-independence theorem controls.  No claim is
made that \(\Delta_2\) is an \(E\)-function, a \(D\)-finite function, or an
arithmetic entire function.
\end{remark}

\begin{remark}[No arithmetic conclusion from the infinite sum]
Algebraic independence controls finite polynomial relations.  It says
nothing by itself about the arithmetic nature of an infinite sum.  Thus
\eqref{eq:S2-sum} and \eqref{eq:Delta2-trace} do not prove that \(S_2\), or
equivalently \(\gamma\), is transcendental.
\end{remark}

\section{Further moving-level towers}
\label{sec:other-towers}

The dyadic tower is not isolated.  This section gives a fixed-coefficient
power-orbit family and a logarithmic lattice.  They have different decay
rates and different independence defects.

\subsection{Fixed-coefficient power orbits}

Fix an integer \(n\ge2\) and a positive algebraic number \(q>1\).  For
\(k\ge0\), put
\begin{equation}\label{eq:Bnk}
 B_{n,k}(q)=
 \frac{n\Ein(q^k)-\Ein(q^{nk})}{n-1}.
\end{equation}

\begin{theorem}[Power-orbit moving tower]
\label{thm:power-orbit}
For fixed \(n\) and \(q\), the family
\(\{B_{n,k}(q):k\ge0\}\) is countably algebraically independent over \(K\).
Moreover
\begin{equation}\label{eq:eta-nk}
 B_{n,k}(q)=\gamma+\eta_{n,k}(q),\qquad
 \eta_{n,k}(q)=
 \frac{nE_1(q^k)-E_1(q^{nk})}{n-1},
\end{equation}
and
\begin{equation}\label{eq:eta-properties}
 \eta_{n,k}>0,\qquad
 \eta_{n,k}\downarrow0,\qquad
 \eta_{n,k}\sim
 \frac n{n-1}\frac{e^{-q^k}}{q^k}.
\end{equation}
For every choice \(m_k\ge2\), the exact traces
\[
 T_{m_k}(B_{n,k}(q))
\]
are countably algebraically independent over \(K\).  For fixed \(m\), they
converge to \(S_m\).  At \(m=2\) they increase to \(S_2\) and yield, by
successive differences, another positive algebraically independent spectral
realization of \(S_2\).
\end{theorem}

\begin{proof}
Let \(X_j=\Ein(q^j)\).  The distinct algebraic points \(q^j\) give an
algebraically independent coordinate family by
\cref{thm:input-independence}.  The row vectors
\[
 v_k=ne_k-e_{nk}\qquad(k\ge0)
\]
are linearly independent: in a finite relation, the coordinate \(nk_0\)
associated with the largest occurring \(k_0\) appears as the terminal
coordinate of that row and as the initial coordinate of no row in the
finite set.  Thus its coefficient vanishes, and descending induction
finishes the proof.  If the finite set consists only of \(k_0=0\), then
\(v_0=(n-1)e_0\ne0\), which is the same conclusion.  Apply
\cref{lem:linear-images}.

The connection formula \eqref{eq:Ein-connection} cancels the logarithmic
terms and gives \eqref{eq:eta-nk}.  For \(x>1\),
\[
 \frac d{dx}\left(
 \frac{nE_1(x)-E_1(x^n)}{n-1}\right)
 =\frac n{n-1}\frac{e^{-x^n}-e^{-x}}x<0.
\]
This proves strict decrease; positivity follows from
\(E_1(x)>E_1(x^n)>0\), and the asymptotic follows from
\eqref{eq:E1-asymptotic}.  The trace assertions follow from
\cref{cor:transversal,thm:universal-trace}.  The order-two construction is
the proof of \cref{thm:increments,thm:Delta2} applied to this decreasing
level sequence: indeed \(B_{n,0}(q)=\Ein(1)<1<2\), since
\((1-e^{-t})/t<1\) on \((0,1]\).
\end{proof}

This tower samples the rank-one witness of the companion paper exactly:
\begin{equation}\label{eq:B-Gn}
 B_{n,k}(q)-\gamma=\frac{G_n(q^k)}{n-1}.
\end{equation}

\begin{remark}[Fixing \(n\) is essential]
The full two-parameter family is not algebraically independent.  For
\(r\ge1\),
\begin{equation}\label{eq:cross-n-relation}
 B_{n^r,k}(q)=\frac{n-1}{n^r-1}
 \sum_{j=0}^{r-1}n^{r-1-j}B_{n,n^jk}(q).
\end{equation}
For example \(B_{4,k}=(2B_{2,k}+B_{2,2k})/3\).
\end{remark}

\subsection{Algebraic-level inverse transcendence}

The first difference between two members of the \(n=2\) power orbit gives
the gamma-free function
\begin{equation}\label{eq:scriptR-def}
 \mathscr R(q)=E_1(1)-2E_1(q)+E_1(q^2)\qquad(q>1).
\end{equation}
It has an unexpectedly rigid inverse on algebraic levels.

\begin{theorem}[Transcendental inverse levels]
\label{thm:inverse-levels}
The map \(\mathscr R:(1,\infty)\to(0,E_1(1))\) is a strictly increasing
bijection.  If \(a\in\Qbar\cap(0,E_1(1))\), then the unique number
\(r_a>1\) satisfying \(\mathscr R(r_a)=a\) is transcendental.

In particular, for every integer \(m\ge6\), the unique root \(r_m>1\) of
\begin{equation}\label{eq:rm-level}
 E_1(1)-2E_1(r_m)+E_1(r_m^2)=\frac1m
\end{equation}
is transcendental, and
\begin{equation}\label{eq:rm-asymptotic}
 r_m=1+\sqrt{\frac em}+\frac e{2m}+O(m^{-3/2}).
\end{equation}
\end{theorem}

\begin{proof}
Direct differentiation gives
\[
 \mathscr R'(q)=\frac2q(e^{-q}-e^{-q^2})>0.
\]
Moreover \(\mathscr R(1)=0\) and
\(\mathscr R(q)\to E_1(1)\) as \(q\to\infty\), proving the bijection.
If \(r_a\) were algebraic, the connection formula would turn
\(\mathscr R(r_a)=a\) into
\[
 \Ein(1)-2\Ein(r_a)+\Ein(r_a^2)=a.
\]
The three arguments are distinct nonzero algebraic numbers, so this relation
contradicts \cref{thm:input-independence}.

Indeed, after writing \(t=1+u\), Jensen's inequality for the probability
density \(e^{-u}\,du\) on \([0,\infty)\) gives
\[
 E_1(1)=e^{-1}\int_0^\infty\frac{e^{-u}}{1+u}\,du
 >\frac1{2e}>\frac16.
\]
Thus \(a=1/m\) is admissible for \(m\ge6\).  Finally, writing \(q=1+h\),
direct differentiation at one gives
\[
 \mathscr R(1+h)=e^{-1}(h^2-h^3+O(h^4)).
\]
Series inversion at \(\mathscr R=1/m\) proves
\eqref{eq:rm-asymptotic}.
\end{proof}

\subsection{A logarithmic lattice and density-one transcendence}

For \(r\ge1\), put
\begin{equation}\label{eq:cr}
 c_r=\Cin(r)-\log r=\gamma-\Ci(r).
\end{equation}
These levels approach \(\gamma\) only polynomially, but the loss caused by
the logarithms can be counted by primes.

\begin{lemma}\label{lem:cr-positive}
One has \(c_r>0\) for every integer \(r\ge1\).
\end{lemma}

\begin{proof}
For \(x>0\), set \(c(x)=\gamma-\Ci(x)\).  Since
\(c'(x)=-\cos x/x\), its local minima occur at
\(x_k=\pi/2+2\pi k\).  At the first minimum,
\[
 \Cin(\pi/2)
 \ge \int_0^{\pi/2}\left(\frac t2-\frac{t^3}{24}\right)dt
 =\frac{\pi^2}{16}-\frac{\pi^4}{1536}
 >\log(\pi/2),
\]
using \(1-\cos t\ge t^2/2-t^4/24\) on this interval.  The last strict
inequality follows, for example, from
\(333/106<\pi<355/113\) and
\(\log(1+u)\le u-u^2/2+u^3/3\) for \(0<u<1\).
Thus
\(c(x_0)>0\).  For \(k\ge1\), integration by parts gives
\(\Ci(x_k)\le2/x_k<1/2<\gamma\).  Here we used
\(x_k\ge5\pi/2\) and the elementary bound \(\gamma>1/2\).
Hence all local minima are positive, and so is
\(c(x)\).  In particular \(c_r>0\).
\end{proof}

For each \(r\), let
\begin{equation}\label{eq:Ur}
 U_r=T_2(c_r)=\frac1{c_r^2}-\frac1{2c_r}.
\end{equation}

\begin{theorem}[Prime-defect lattice theorem]
\label{thm:prime-defect}
Let \(\pi(M)\) be the number of primes not exceeding \(M\).  Then
\begin{equation}\label{eq:prime-defect-trdeg}
 \trdeg_KK(c_1,\ldots,c_M)
 =\trdeg_KK(U_1,\ldots,U_M)
 \ge M-\pi(M).
\end{equation}
Consequently at most \(\pi(M)\) of the first \(M\) members of each family
\(\{c_r\}\) and \(\{U_r\}\) are algebraic over \(K\).  Both families are
\(K\)-transcendental on a set of natural density one.
\end{theorem}

\begin{proof}
The companion result \cite[Cor.~5.5]{KEIO2026I} makes
\(\Cin(1),\ldots,\Cin(M)\) algebraically independent over \(K\).
Equivalently, this follows from
\(\Cin(r)=\{\Ein(ir)+\Ein(-ir)\}/2\): the pairs
\(\{ir,-ir\}\) are disjoint, so these averaging rows have full rank.
Since
\[
 \log r=\sum_{p\le M}v_p(r)\log p,
\]
we have
\[
 K(\Cin(1),\ldots,\Cin(M))
 \subset K(c_1,\ldots,c_M,\log p:p\le M).
\]
Therefore
\[
 M\le \trdeg_KK(c_1,\ldots,c_M)+\pi(M),
\]
which gives the lower bound for the \(c_r\).  Equation
\[
 2U_rc_r^2+c_r-2=0
\]
shows that \(K(c_1,\ldots,c_M)\) is algebraic over
\(K(U_1,\ldots,U_M)\), while the reverse field inclusion is rational.
Thus the two transcendence degrees agree.  If \(a\) of the first \(M\)
coordinates were algebraic over \(K\), the transcendence degree would be at
most \(M-a\); hence \(a\le\pi(M)\).  The prime number theorem completes the
density statement.
\end{proof}

The trace \(U_r\) also converges to \(S_2\).  Direct subtraction gives
\begin{equation}\label{eq:Ur-S2-exact}
 U_r-S_2=
 \frac{\Ci(r)\{\gamma(4-\gamma)-(2-\gamma)\Ci(r)\}}
 {2\gamma^2(\gamma-\Ci(r))^2},
\end{equation}
and the standard expansion of \(\Ci\) yields
\begin{align}
 U_r-S_2
 &=\frac{4-\gamma}{2\gamma^3}\frac{\sin r}{r}
 -\frac{4-\gamma}{2\gamma^3}\frac{\cos r}{r^2}\notag\\
 &\quad+
 \frac{6-\gamma}{2\gamma^4}\frac{\sin^2r}{r^2}
 +O(r^{-3}).\label{eq:Ur-S2-asymptotic}
\end{align}
Thus this lattice is generally only \(O(r^{-1})\) from the Euler trace.
The theorem does not decide any specified \(c_r,U_r\) with \(r\ge2\);
\(c_1=\Cin(1)\) is already \(K\)-transcendental.

\section{Normalized certificates and cross-index rigidity}
\label{sec:certificates}

The moving towers concern limits.  We now return to finite witnesses and
show that all normalized linear certificates are unique.

Let
\[
 V=\bigoplus_{\alpha\in\Qbar^\times}K^{\mathrm{alg}}[\alpha]
\]
be the free \(K^{\mathrm{alg}}\)-vector space with finite support on the
nonzero algebraic points, and define
\[
 \Phi:V\longrightarrow\C,
 \qquad \Phi([\alpha])=\Ein(\alpha).
\]

\begin{lemma}[The evaluation intersection]
\label{lem:evaluation-intersection}
The map \(\Phi\) is injective and
\begin{equation}\label{eq:Phi-intersection}
 \Phi(V)\cap K^{\mathrm{alg}}=\{0\}.
\end{equation}
\end{lemma}

\begin{proof}
By \cref{thm:input-independence}, every finite collection of the displayed
\(\Ein\)-values is algebraically independent over \(K^{\mathrm{alg}}\), and
hence linearly independent.  Moreover any nonzero linear form in
algebraically independent variables is transcendental over the base field:
extend it to an invertible linear coordinate system.  This proves both
claims.
\end{proof}

\begin{definition}
A finite \(K^{\mathrm{alg}}\)-certificate for \(\gamma\) is a triple
\((v,c,a)\in V\times(K^{\mathrm{alg}})^\times\times K^{\mathrm{alg}}\)
satisfying
\begin{equation}\label{eq:certificate}
 \Phi(v)=c\gamma+a.
\end{equation}
Its normalized data are \((v/c,a/c)\).
\end{definition}

\begin{theorem}[Normalized certificate uniqueness]
\label{thm:certificate-uniqueness}
If \((v_1,c_1,a_1)\) and \((v_2,c_2,a_2)\) are two certificates, then
\begin{equation}\label{eq:normalized-certificate-unique}
 \boxed{\displaystyle
 \frac{v_1}{c_1}=\frac{v_2}{c_2},\qquad
 \frac{a_1}{c_1}=\frac{a_2}{c_2}.}
\end{equation}
\end{theorem}

\begin{proof}
Subtracting the two normalized identities gives
\[
 \Phi\!\left(\frac{v_1}{c_1}-\frac{v_2}{c_2}\right)
 =\frac{a_1}{c_1}-\frac{a_2}{c_2}\in K^{\mathrm{alg}}.
\]
Apply \eqref{eq:Phi-intersection}.
\end{proof}

This is an exclusion theorem, not an existence theorem.  The certificate
set may be empty.  Its force is that whenever two apparently unrelated
algebraicity events would produce certificates, their normalized supports
must coincide exactly.

We apply it to the real Euler witnesses recalled in the introduction.  The
existence, location and within-index algebraic-zero rigidity of \(q_*\),
\(\beta_n\) and \(\lambda_n\) were established in
\cite[Thm.~14.2, Prop.~14.3, and Thm.~14.4]{KEIO2026I}.

\begin{theorem}[Cross-index witness rigidity]
\label{thm:cross-index}
At most one member of
\begin{equation}\label{eq:witness-set}
 \mathcal W=
 \{q_*\}\cup\{\beta_n:n\ge2\}
 \cup\{\lambda_n:n\ge3,\ n\text{ odd}\}
\end{equation}
is algebraic.  Thus precisely one of the following alternatives holds:
\begin{enumerate}[label=\textup{(\roman*)}]
 \item every member of \(\mathcal W\) is transcendental;
 \item exactly one member is algebraic, and in that event \(\gamma\) and
 \(\delta\) are algebraically independent over \(K\).
\end{enumerate}
\end{theorem}

\begin{proof}
If \(q_*\) is algebraic, its normalized certificate vector is
\[
 v_*=[q_*],\qquad \Phi(v_*)=\gamma.
\]
If an algebraic number \(\alpha\) is a zero of \(G_n\), its normalized
certificate vector is
\begin{equation}\label{eq:G-certificate-vector}
 v_{n,\alpha}
 =\frac n{n-1}[\alpha]-\frac1{n-1}[\alpha^n],
 \qquad \Phi(v_{n,\alpha})=\gamma.
\end{equation}
The real-zero classification places \(\alpha\) in \((0,1)\) or, for odd
\(n\), in \((-\infty,-1)\).  Hence \(\alpha\ne\alpha^n\), and
\(v_{n,\alpha}\) has exactly one positive and one negative support
coefficient.  It cannot equal the one-point vector \(v_*\).

If \(v_{n,\alpha}=v_{m,\beta}\), equality of the positive and negative
support coordinates gives
\[
 \frac n{n-1}=\frac m{m-1},
\]
so \(n=m\), and then \(\alpha=\beta\).  Certificate uniqueness therefore
excludes two distinct algebraic members of \eqref{eq:witness-set}.  The
algebraic independence of \(\gamma,\delta\) in the exceptional case is
\cite[Thm.~14.2 and Thm.~14.4]{KEIO2026I}.
\end{proof}

\subsection{A defect-one family of critical cancellation coefficients}

Fix an algebraic number \(q>1\).  For \(N\ge1\), put
\begin{align}
 D_N(q)&=E_1(q^N)-2E_1(q^{N+1})+E_1(q^{N+2}),\label{eq:DNq}\\
 A_N(q)&=(N+2)E_1(q^{N+1})-(N+1)E_1(q^{N+2}),\label{eq:ANq}\\
 u_N(q)&=\frac{A_N(q)}{D_N(q)}.\label{eq:uNq}
\end{align}
Both numerator and denominator are positive: \(D_N(q)>0\) is the strict
second-difference inequality for the strictly convex function
\(t\mapsto E_1(q^t)\), while monotonicity gives
\(A_N(q)>E_1(q^{N+1})>0\).  The number \(u_N(q)\) is the
unique positive coefficient for which the three-point augmentation-one tail
vanishes:
\begin{equation}\label{eq:uN-tail-zero}
 -u_NE_1(q^N)+(N+2+2u_N)E_1(q^{N+1})
 -(N+1+u_N)E_1(q^{N+2})=0.
\end{equation}

\begin{theorem}[One-defect cancellation family]
\label{thm:uN-defect}
For every finite set \(I\subset\N_{\ge1}\),
\begin{equation}\label{eq:uN-trdeg}
 \boxed{\displaystyle
 \trdeg_K K(u_N(q):N\in I)\ge |I|-1.}
\end{equation}
Consequently at most one member of \(\{u_N(q):N\ge1\}\) is algebraic over
\(K\).  If one member is algebraic over \(K\), all the remaining members are
countably algebraically independent over \(K\), and \(\gamma,\delta\) are
algebraically independent over \(K\).  If \(\gamma\) is algebraic over
\(K\), then the whole family \(\{u_N(q):N\ge1\}\) is countably
algebraically independent over \(K\).

Finally,
\begin{equation}\label{eq:uN-asymptotic}
 u_N(q)\sim\frac{N+2}{q}
 \exp\!\bigl(-(q-1)q^N\bigr).
\end{equation}
\end{theorem}

\begin{proof}
Let \(X_j=\Ein(q^j)\).  The connection formula and
\eqref{eq:uN-tail-zero} give
\begin{equation}\label{eq:uN-certificate}
 -u_NX_N+(N+2+2u_N)X_{N+1}
 -(N+1+u_N)X_{N+2}=\gamma.
\end{equation}
For finite \(I\), let
\(U=\bigcup_{N\in I}\{N,N+1,N+2\}\).  Since \(u_N>0\), the relations
\eqref{eq:uN-certificate}, used in decreasing order of \(N\), solve every
\(X_N\) with \(N\in I\) in terms of \(u_N\), \(\gamma\), and the
coordinates \(X_j\) with \(j\in U\setminus I\).  Hence
\[
 K(X_j:j\in U)
 \subset
 K(u_N:N\in I,\gamma,X_j:j\in U\setminus I).
\]
The left-hand coordinates are algebraically independent over \(K\), so
\[
 |U|\le
 \trdeg_KK(u_N:N\in I)+1+|U|-|I|,
\]
which proves \eqref{eq:uN-trdeg}.  If one \(u_{N_0}\) is algebraic over
\(K\), apply the bound to \(I\cup\{N_0\}\) to obtain full independence of
every finite set with \(N_0\) omitted.

In that exceptional case \eqref{eq:uN-certificate} expresses \(\gamma\) as
a nonconstant \(K^{\rm alg}\)-linear form in three algebraically independent
\(\Ein\)-coordinates.  Together with
\(\delta/e=\Ein(1)-\gamma\), where \(1\) is a fourth distinct algebraic
point, two linearly independent coordinate forms are obtained.  Thus
\(\gamma,\delta\) are algebraically independent over \(K\).  If instead
\(\gamma\in K^{\rm alg}\), the extra ``\(+1\)'' in the preceding
transcendence-degree estimate disappears, proving full independence of the
\(u_N\).

The asymptotic follows from \eqref{eq:E1-asymptotic}: the first terms in
\eqref{eq:DNq} and \eqref{eq:ANq} dominate and
\[
 \frac{E_1(q^{N+1})}{E_1(q^N)}
 \sim q^{-1}\exp\!\bigl(-(q-1)q^N\bigr).
\]
\end{proof}

\section{Relative exclusion and finite trace descent}
\label{sec:relative}

Put
\[
 A=\Ein(1),\qquad B=\Ein(-1),\qquad
 C=\Cin(1)=\frac{\Ein(i)+\Ein(-i)}2.
\]
The numbers (A,B,C) are algebraically independent over \(K\), and
\begin{align}
 E_1(1)&=A-\gamma,&
 \Ei(1)&=\gamma-B,\label{eq:six-identities}\\
 \Ci(1)&=\gamma-C,&
 \Chi(1)&=\gamma-\frac{A+B}{2},&
 \delta&=e(A-\gamma).\notag
\end{align}

\begin{theorem}[Five relative certificate classes]
\label{thm:five-classes}
Among the following five assertions, at most one can hold:
\begin{align*}
 \mathcal C_0&:\ \gamma\in K^{\mathrm{alg}},\\
 \mathcal C_1&:\ E_1(1)\in K^{\mathrm{alg}}
 \quad\Longleftrightarrow\quad \delta\in K^{\mathrm{alg}},\\
 \mathcal C_2&:\ \Ei(1)\in K^{\mathrm{alg}},\\
 \mathcal C_3&:\ \Ci(1)\in K^{\mathrm{alg}},\\
 \mathcal C_4&:\ \Chi(1)\in K^{\mathrm{alg}}.
\end{align*}
Consequently at least four members of
\[
 \gamma,\quad\delta,\quad E_1(1),\quad\Ei(1),\quad\Ci(1),\quad\Chi(1)
\]
are transcendental over \(K\).  If the possible exceptional class is not
\(\mathcal C_1\), then at least five are transcendental over \(K\).
\end{theorem}

\begin{proof}
The equivalence in \(\mathcal C_1\) follows from
\(\delta=eE_1(1)\) and \(e\in K^\times\).  The five assertions produce the
following normalized certificate vectors, respectively:
\[
 0,\qquad [1],\qquad[-1],\qquad
 \frac{[i]+[-i]}2,\qquad
 \frac{[1]+[-1]}2.
\]
For example, if \(E_1(1)\in K^{\mathrm{alg}}\), then
\(\Phi([1])=\gamma+E_1(1)\); if \(\Ci(1)\in K^{\mathrm{alg}}\), then
\(\Phi(([i]+[-i])/2)=\gamma-\Ci(1)\).  These five vectors are distinct, so
\cref{thm:certificate-uniqueness} permits at most one assertion.  Counting
the sizes (1,2,1,1,1) of the five classes gives the final statement.
\end{proof}

\begin{remark}[The necessary correction]
The theorem does \emph{not} say that at most one of the six constants is
algebraic over \(K\).  The pair \(E_1(1),\delta\) forms one relative class:
they are simultaneously algebraic or transcendental over \(K\).  Over
\(\Qbar\), the companion paper excludes their simultaneous algebraicity by
the additional identity \(e=\delta/E_1(1)\) and the transcendence of \(e\).
\end{remark}

\begin{corollary}[Witnesses exclude every relative class]
\label{cor:witness-all-K-trans}
If one member of the witness set \(\mathcal W\) in
\eqref{eq:witness-set} is algebraic, then all six constants displayed in
\cref{thm:five-classes} are transcendental over \(K\).
\end{corollary}

\begin{proof}
The normalized vector of an algebraic witness is either \([q_*]\) or a
two-point vector with one positive and one negative coefficient as in
\eqref{eq:G-certificate-vector}.  Its support and sign pattern differ from
each of the five vectors in the proof of \cref{thm:five-classes}.  Certificate
uniqueness therefore excludes all five relative assertions.
\end{proof}

We next separate finite field membership from countable analytic limits.
Let
\begin{equation}\label{eq:G-field}
 \calG=K\bigl(\Ein(\alpha):\alpha\in\Qbar^\times\bigr).
\end{equation}
Every element of \(\calG\) uses only finitely many algebraic-point values,
although the generating family is countable.

\begin{proposition}[Relative trace descent]
\label{prop:relative-trace-descent}
For \(m\ge2\),
\[
 S_m=P_m(\gamma^{-1}),\qquad
 \calG(S_m)\subset\calG(\gamma),\qquad
 [\calG(\gamma):\calG(S_m)]\le m.
\]
Consequently
\begin{equation}\label{eq:relative-trdeg}
 \trdeg_{\calG}\calG(S_m)=\trdeg_{\calG}\calG(\gamma).
\end{equation}
Moreover
\begin{equation}\label{eq:G-S2-S3}
 \boxed{\calG(S_2,S_3)=\calG(S_2,S_4)=\calG(\gamma)},
\end{equation}
and hence
\begin{equation}\label{eq:joint-gamma-free}
 S_2,S_3\in\calG
 \quad\Longleftrightarrow\quad
 S_2,S_4\in\calG
 \quad\Longleftrightarrow\quad
 \gamma\in\calG.
\end{equation}
Individually,
\[
 S_2\in\calG\Longrightarrow[\calG(\gamma):\calG]\le2,
 \qquad
 S_3\in\calG\Longrightarrow[\calG(\gamma):\calG]\le3,
 \qquad
 S_4\in\calG\Longrightarrow[\calG(\gamma):\calG]\le4.
\]
\end{proposition}

\begin{proof}
The polynomial identity comes from \cref{thm:universal-trace} at the Euler
level.  Thus \(\gamma^{-1}\) is algebraic of degree at most \(m\) over
\(\calG(S_m)\), proving the first assertions and
\eqref{eq:relative-trdeg}.  Equations \eqref{eq:gamma-recovery} and
\eqref{eq:gamma-even-recovery} prove \eqref{eq:G-S2-S3}; the remaining
statements follow.
\end{proof}

The reverse of the first individual implication for \(S_2\) is false without
an additional norm--trace condition.

\begin{theorem}[Exact quadratic descent criterion]
\label{thm:quadratic-descent}
One has \(S_2\in\calG\) if and only if either \(\gamma\in\calG\), or
\[
 [\calG(\gamma):\calG]=2
 \quad\text{and}\quad
 \Norm_{\calG(\gamma)/\calG}(\gamma)
 =2\Tr_{\calG(\gamma)/\calG}(\gamma).
\]
In the quadratic case the conjugate is
\begin{equation}\label{eq:gamma-conjugate}
 \gamma'=-\frac{2\gamma}{2-\gamma}.
\end{equation}
\end{theorem}

\begin{proof}
Let \(f(x)=x^{-2}-(2x)^{-1}\).  If \(\gamma\notin\calG\) and
\(f(\gamma)=S_2\in\calG\), then \cref{prop:relative-trace-descent} gives a
quadratic extension.  If \(\gamma'\) is its nontrivial conjugate, then
\[
 f(\gamma)-f(\gamma')
 =\frac{(\gamma'-\gamma)
 \{2(\gamma+\gamma')-\gamma\gamma'\}}
 {2\gamma^2{\gamma'}^2}.
\]
Thus invariance of \(f(\gamma)\) is equivalent to
\(\gamma\gamma'=2(\gamma+\gamma')\), which is exactly the displayed
norm--trace condition and rearranges to \eqref{eq:gamma-conjugate}.
Conversely that condition makes \(f(\gamma)\) fixed by the quadratic Galois
involution and hence an element of \(\calG\).
\end{proof}

For instance, degree two alone is insufficient: with base field \(\Q\) and
\(x=\sqrt2\), one has
\(f(x)=1/2-\sqrt2/4\notin\Q\).

\section{An explicit solution of a 2026 Euler recurrence}
\label{sec:ramanujan-recurrence}

The Ramanujan Challenge for AI \cite{RamanujanChallenge2026} asks, in its
Euler-constant problem, for a proof that the quotient of two distinguished
solutions of the following four-term recurrence tends to \(\gamma\).  We
give a self-contained solution.  The finite sums below belong to the
multiple-Laguerre framework developed by Wolfs--Van Assche
\cite{WolfsVanAssche2025}.  What is proved here is the adjacent
specialization, explicit all-index recurrence certificates, an elementary
positive concentration proof, and the flat matrix field.

Put
\begin{align*}
 A_n&=8n^3+51n^2+105n+68,\\
 B_n&=24n^5+337n^4+1833n^3+4818n^2+6092n+2928,\\
 C_n&=(n+2)(n+3)
 (24n^5+273n^4+1150n^3+2154n^2+1635n+268),\\
 D_n&=(n+1)(n+2)^4(n+3)
 (8n^3+75n^2+231n+232),
\end{align*}
and
\begin{equation}\label{eq:challenge-recurrence}
 A_nu_n=B_nu_{n-1}-C_nu_{n-2}+D_nu_{n-3}
 \qquad(n\ge0).
\end{equation}

\begin{theorem}[Closed solution of the Euler recurrence]
\label{thm:challenge-solution}
For \(n\ge-3\), set \(a=n+3\), \(b=n+4=a+1\), and
\begin{equation}\label{eq:Wab}
 W_{a,b}(k)=
 \frac{(a!b!)^2}{(k!)^3(a-k)!(b-k)!}
 \qquad(0\le k\le a).
\end{equation}
Then the two sequences determined by
\[
 (p_{-3},p_{-2},p_{-1})=(0,7,179),
 \qquad
 (q_{-3},q_{-2},q_{-1})=(1,12,306)
\]
and \eqref{eq:challenge-recurrence} have the exact formulas
\begin{align}
 q_n&=b\sum_{k=0}^{a}W_{a,b}(k),
 \label{eq:qn-closed}\\
 p_n&=b\sum_{k=0}^{a}W_{a,b}(k)
 \bigl(3H_k-H_{a-k}-H_{b-k}\bigr)+a!.
 \label{eq:pn-closed}
\end{align}
In particular,
\begin{equation}\label{eq:challenge-limit}
 \boxed{\displaystyle\lim_{n\to\infty}\frac{p_n}{q_n}=\gamma.}
\end{equation}
\end{theorem}

\begin{proof}
Substitution of \((a,b)=(0,1),(1,2),(2,3)\) gives the six initial
values.  The two symbolic telescoping identities in
\cref{app:telescopers} prove \eqref{eq:challenge-recurrence} for both
closed forms at every \(n\ge0\).  We prove the limit directly.

After cancelling factors independent of \(k\), define a probability law by
\begin{equation}\label{eq:challenge-pmf}
 \Pr(K_n=k)=
 \frac{\binom ak\binom bk/k!}
 {\sum_{j=0}^{a}\binom aj\binom bj/j!}.
\end{equation}
Then
\begin{equation}\label{eq:challenge-expectation}
 \frac{p_n}{q_n}
 =\mathbb E\bigl(3H_{K_n}-H_{a-K_n}-H_{b-K_n}\bigr)
 +\frac{a!}{q_n}.
\end{equation}
The ratio of consecutive unnormalized masses is
\begin{equation}\label{eq:challenge-ratio}
 r_k=\frac{(a-k)(b-k)}{(k+1)^3},
\end{equation}
which decreases strictly in \(k\).  Put \(x_n=(ab)^{1/3}\).  For each
fixed \(\eta>0\), the ratios outside
\([(1-\eta)x_n,(1+\eta)x_n]\) are uniformly separated from one once
\(n\) is large.  Multiplying the ratios on either side of the mode gives
\begin{equation}\label{eq:challenge-concentration}
 \Pr\bigl(|K_n/x_n-1|>\eta\bigr)
 =O\bigl(e^{-c_\eta x_n}\bigr).
\end{equation}
Hence \(K_n/x_n\to1\) in probability.

On the event \(|K_n/x_n-1|\le\eta\), the estimate
\(H_m=\log m+\gamma+O(m^{-1})\) yields
\begin{align*}
 3H_k-H_{a-k}-H_{b-k}
 &=\gamma+\log\frac{k^3}{(a-k)(b-k)}+o(1)\\
 &=\gamma+O(\eta)+o(1).
\end{align*}
The expression is \(O(\log n)\) on the full support, while the exceptional
probability in \eqref{eq:challenge-concentration} is exponentially small in
\(x_n\).  First let \(n\to\infty\), and then let \(\eta\downarrow0\); its
expectation therefore tends to \(\gamma\).  Finally the
\(k=0\) summand in \eqref{eq:qn-closed} gives
\[
 0\le\frac{a!}{q_n}\le\frac1{b\,b!}\longrightarrow0.
\]
Equation \eqref{eq:challenge-expectation} proves
\eqref{eq:challenge-limit}.
\end{proof}

\begin{remark}[What has and has not been solved]
The theorem gives an independent proof of the publicly stated limit problem.
The official proof and code are now public \cite{RamanujanOfficial2026}; the
sentence in Version 2.0 saying otherwise is superseded.  No priority claim is
made for the official Meijer--\(G\) trajectory, Miller selection, or the broad
multiple-Laguerre family.  The comparison with the terminating state is
established in \cref{sec:meijer-duality}.
The convergence \(p_n/q_n\to\gamma\) does not by itself imply the
irrationality of \(\gamma\).
\end{remark}

The same recurrence has a second nontrivial Ap\'ery limit.  Let \(w_n\) be
the solution with
\begin{equation}\label{eq:w-initial}
 (w_{-3},w_{-2},w_{-1})=(1,18,440),
\end{equation}
and let \(\mathfrak h\) be the positive constant in
\eqref{eq:h-constant} below.

\begin{theorem}[Complete initial-value limit classification]
\label{thm:ram-full-solution}
One has
\begin{equation}\label{eq:w-limit}
 \boxed{\displaystyle\frac{w_n}{q_n}\longrightarrow\mathfrak h
 =1.4322057346532244148\ldots.}
\end{equation}
For arbitrary rational initial data
\((x,y,z)=(u_{-3},u_{-2},u_{-1})\), the corresponding solution satisfies
\begin{align}
 \lim_{n\to\infty}\frac{u_n}{q_n}
 =\frac1{136}\bigl[&(-6x+179y-7z)\mathfrak h\notag\\
 &+(-228x-134y+6z)\gamma
 +(142x-179y+7z)\bigr].
 \label{eq:all-initial-limits}
\end{align}
Consequently the set of all rational-initial-data limits is exactly
\begin{equation}\label{eq:attainable-limit-span}
 \boxed{\Q+\Q\gamma+\Q\mathfrak h.}
\end{equation}
No assertion that this span has dimension three is made.
\end{theorem}

\begin{proof}
The official cyclic gauge at zero has rows
\begin{equation}\label{eq:U0-rows}
 \mathsf U(0)=
 \begin{pmatrix}1&18&440\\0&7&179\\0&1&45\end{pmatrix}.
\end{equation}
They generate solutions \(w,p,s\), respectively, and
\(q=w-p+s\).  The official Miller-selection theorem, transported through
the exact gauge of \cref{thm:ram-gauge}, gives
\begin{equation}\label{eq:Miller-limit-vector}
 \frac1{q_n}(w_n,p_n,s_n)
 \longrightarrow
 (\mathfrak h,\gamma,1+\gamma-\mathfrak h);
\end{equation}
see \cite[\S5.2]{RamanujanOfficial2026} and
\cref{thm:meijer-probability}.  The identity \(q=w-p+s\) normalizes the
right side to one.  Solving the initial vector
\((x,y,z)\) in the basis \((q,p,w)\) gives the three coefficients in
\eqref{eq:all-initial-limits}.  Conversely \(q,p,w\) realize
\(1,\gamma,\mathfrak h\), proving \eqref{eq:attainable-limit-span}.
\end{proof}

\begin{theorem}[Complete integer initial lattice]
\label{thm:ram-integer-lattice}
For a solution of \eqref{eq:challenge-recurrence},
\begin{equation}\label{eq:integer-lattice}
 \boxed{\displaystyle
 u_n\in\Z\ \text{for every }n\ge-3
 \Longleftrightarrow
 x,y,z\in\Z\ \text{and }13x+6y+z\equiv0\pmod{17}.}
\end{equation}
\end{theorem}

\begin{proof}
The first new term is
\begin{equation}\label{eq:u0-congruence}
 u_0=\frac{2784x-402y+732z}{17},
\end{equation}
so integrality forces the stated congruence.  Conversely its kernel in
\(\Z^3\) has basis
\[
 (1,0,4),\qquad(0,1,11),\qquad(0,0,17).
\]
The finite companion-gauge certificate in
\cref{app:integer-lattice-certificate} proves that the three corresponding
solutions are integral at every index.  Every integral vector satisfying the
congruence is their integral combination.
\end{proof}

\section{A flat conservative matrix field and exterior arithmetic}
\label{sec:cmf}

For integers \(a,b\ge0\), define
\begin{equation}\label{eq:Fab}
 F_{a,b}(z)={}_{2}F_{2}
 \!\left(\begin{matrix}-a,-b\\1,1\end{matrix};z\right),
 \qquad
 \theta=z\frac d{dz},
 \qquad
 \mathbf Y_{a,b}=
 \begin{pmatrix}F_{a,b}\\ \theta F_{a,b}\\ \theta^2F_{a,b}\end{pmatrix}.
\end{equation}
It satisfies
\begin{equation}\label{eq:Fab-ode}
 \theta^3F_{a,b}=z(\theta-a)(\theta-b)F_{a,b}.
\end{equation}
At \(z=1\), its terminating series gives the exact link
\begin{equation}\label{eq:q-Fab-link}
 F_{a,b}(1)=\sum_{k=0}^{a}\binom ak\binom bk\frac1{k!},
 \qquad q_n=b\,a!b!\,F_{a,b}(1)\quad(a=n+3,\ b=n+4).
\end{equation}

\begin{theorem}[Flat \(SL_3\) conservative matrix field]
\label{thm:flat-cmf}
For \(a,b>0\), put
\begin{align}
 M_a(a,b;z)&=\frac1a
 \begin{pmatrix}
 a&-1&0\\
 0&a&-1\\
 -zab&z(a+b)&a-z
 \end{pmatrix},\label{eq:Ma}\\
 M_b(a,b;z)&=\frac1b
 \begin{pmatrix}
 b&-1&0\\
 0&b&-1\\
 -zab&z(a+b)&b-z
 \end{pmatrix}.\label{eq:Mb}
\end{align}
Then
\[
 \mathbf Y_{a-1,b}=M_a(a,b;z)\mathbf Y_{a,b},
 \qquad
 \mathbf Y_{a,b-1}=M_b(a,b;z)\mathbf Y_{a,b},
\]
and
\begin{align}
 \det M_a&=\det M_b=1,\label{eq:cmf-det}\\
 M_b(a-1,b;z)M_a(a,b;z)
 &=M_a(a,b-1;z)M_b(a,b;z).
 \label{eq:cmf-flatness}
\end{align}
Thus the denominator system of \cref{thm:challenge-solution} is the
adjacent-diagonal path \((a,b)=(n+3,n+4)\) in an explicit flat \(SL_3\)
conservative matrix field.
\end{theorem}

\begin{proof}
The two contiguous relations follow by comparing the terminating
hypergeometric coefficients and using \eqref{eq:Fab-ode} for the third row.
Expansion of the determinants gives \(1\).  Direct matrix multiplication
gives the polynomial identity \eqref{eq:cmf-flatness}; an entrywise symbolic
certificate is also included with the source package.
\end{proof}

\begin{remark}[Path rigidity]
Flatness says that changing a lattice path while keeping the same endpoint
\((a,b)\) leaves \(F_{a,b}\), and hence the endpoint-defined denominator,
unchanged.  The numerator and quotient in \cref{thm:challenge-solution} are
also endpoint-defined by their closed formulas.  Rerouting a path therefore
does not change these particular values.  Changing the endpoint curve is a
different problem and is not excluded by this observation.
\end{remark}

The numerator is the first Frobenius jet of the same flat connection.  Let
\begin{equation}\label{eq:Phi-epsilon}
 \Phi_{a,b}(\epsilon,z)=
 \frac{a!b!}{(1+\epsilon)_a(1+\epsilon)_b}
 {}_3F_3\!\left(
 \begin{matrix}-a-\epsilon,-b-\epsilon,1\\
 1-\epsilon,1-\epsilon,1-\epsilon
 \end{matrix};z\right).
\end{equation}

\begin{theorem}[Flat first-jet \(SL_6\) lift]
\label{thm:ram-sl6}
Fix positive integers \(a,b\), expand the analytic \(\epsilon\)-germ at
zero modulo \(\epsilon^2\), and put
\(\mathbf Y_{a,b}(\epsilon)
=(\Phi_{a,b},\theta\Phi_{a,b},\theta^2\Phi_{a,b})^{\mathsf T}\).
Then
\begin{equation}\label{eq:epsilon-contiguous}
 \mathbf Y_{a-1,b}(\epsilon)=\mathscr M_a(\epsilon)\mathbf Y_{a,b}(\epsilon),
 \qquad
 \mathbf Y_{a,b-1}(\epsilon)=\mathscr M_b(\epsilon)\mathbf Y_{a,b}(\epsilon),
\end{equation}
where
\begin{equation}\label{eq:epsilon-Ma}
 \mathscr M_a(\epsilon)=\frac1a
 \begin{pmatrix}
 a+\epsilon&-1&0\\
 0&a+\epsilon&-1\\
 -zab-\epsilon z(a+b)&z(a+b)+2\epsilon z&a-z-2\epsilon
 \end{pmatrix},
\end{equation}
and
\begin{equation}\label{eq:epsilon-Mb}
 \mathscr M_b(\epsilon)=\frac1b
 \begin{pmatrix}
 b+\epsilon&-1&0\\
 0&b+\epsilon&-1\\
 -zab-\epsilon z(a+b)&z(a+b)+2\epsilon z&b-z-2\epsilon
 \end{pmatrix}.
\end{equation}
In
\(\Q(a,b,z)[\epsilon]/(\epsilon^2)\), these matrices have determinant one
and satisfy the same zero-curvature identity as
\eqref{eq:cmf-flatness}.

Writing
\(\mathscr M_a=\mathscr M_{a,0}+\epsilon\mathscr M_{a,1}\), the block matrix
\begin{equation}\label{eq:SL6-block}
 \begin{pmatrix}
 \mathscr M_{a,0}&0\\
 \mathscr M_{a,1}&\mathscr M_{a,0}
 \end{pmatrix}\in SL_6
\end{equation}
transports
\(\binom{\mathbf Y_{a,b}(0)}
{\partial_\epsilon\mathbf Y_{a,b}(0)}\), and similarly in the
\(b\)-direction.  Thus denominator and numerator path-independence are the
zero- and first-jet parts of one flat connection.

At \(b=a+1,z=1\),
\begin{equation}\label{eq:Phi-pq-jets}
 b\,a!b!\Phi_{a,b}(0,1)=q_{a-3},\qquad
 b\,a!b!\partial_\epsilon\Phi_{a,b}(0,1)=p_{a-3}.
\end{equation}
The boundary \(a=0,b=1\) in \eqref{eq:Phi-pq-jets} is verified directly.
\end{theorem}

\begin{proof}
Modulo \(\epsilon^2\), the Frobenius deformation satisfies
\begin{equation}\label{eq:Phi-dual-ODE}
 (\theta-\epsilon)^3\Phi_{a,b}
 =z(\theta-a-\epsilon)(\theta-b-\epsilon)\Phi_{a,b}.
\end{equation}
Coefficient comparison gives the first two rows of
\eqref{eq:epsilon-contiguous}, while \eqref{eq:Phi-dual-ODE} gives the third.
Direct determinant expansion and multiplication give determinant one and
flatness in the dual-number ring.  Expanding a
dual-number matrix action gives the block form \eqref{eq:SL6-block}.
At \(\epsilon=0\), \(\Phi_{a,b}=F_{a,b}\).  The derivative of the terminating
part gives the harmonic weight in \eqref{eq:pn-closed}; the regularized term
at \(k=b=a+1\) is \(1/(bb!)\) before multiplication by \(b\,a!b!\), and
hence gives exactly the endpoint \(+a!\).
\end{proof}

\begin{proposition}[First-jet exactness barrier]
\label{prop:first-jet-barrier}
Let \(\mathscr K_a\) be the rational carrier defined in
\eqref{eq:Kas-def}.  For \(b=a+1,z=1\),
\begin{equation}\label{eq:Phi-K-contact}
 \Phi_{a,b}(\epsilon,1)-\mathscr K_a(\epsilon)=O(\epsilon^2).
\end{equation}
Return here to the full analytic \(\epsilon\)-germ.  For every \(r\ge1\),
the term with \(k=b+r\) contributes
\begin{equation}\label{eq:E-tail-second-jet}
 [\epsilon^2]\,\Phi_{a,b}(\epsilon,1)\big|_{k=b+r}
 =-\frac{a!b!\,r!(r-1)!}{(b+r)!^3}\ne0.
\end{equation}
The total second-order difference cannot cancel:
\begin{equation}\label{eq:second-jet-total-positive}
 [\epsilon^2]\{\Phi_{a,b}(\epsilon,1)-\mathscr K_a(\epsilon)\}
 =
 \frac1{b^2a!}\left\{
 3H_b-\sum_{r\ge1}\frac{r!(r-1)!}{(b+1)_r^3}\right\}>0.
\end{equation}
Indeed \((b+1)_r\ge(r+1)!\), so the sum is less than one, while
\(3H_b\ge3\).
Thus the first derivative is exactly terminating after its resonant endpoint,
whereas the second derivative develops a genuine infinite hypergeometric
tail.  This is a barrier for this Frobenius deformation, not a universal
no-go for all higher-jet constructions.
\end{proposition}

\begin{proof}
For \(0\le k\le a\), the hypergeometric summand is identically the
corresponding summand of \(\mathscr J_a(\epsilon)\).  At the resonant index
\(k=b\), its analytic continuation is
\[
 \frac{\epsilon}{b^2a!}
 \prod_{j=1}^b\left(1-\frac{\epsilon}{j}\right)^{-3};
\]
its linear coefficient is \(1/(bb!)\), exactly the correction in
\(\mathscr K_a\), and its quadratic coefficient is
\(3H_b/(b^2a!)\).  For \(k=b+r\), \(r\ge1\), direct expansion gives
\eqref{eq:E-tail-second-jet}.  Summing those coefficients gives
\eqref{eq:second-jet-total-positive}; the displayed comparison proves its
strict positivity.  Hence the zero and first jets agree and the second
jets do not.
\end{proof}

\begin{proposition}[Forced divisor and exact Casoratian]
\label{prop:challenge-arithmetic}
Let \(d_b=\lcm(1,\ldots,b)\).  For every \(n\ge-3\),
\begin{equation}\label{eq:challenge-gcd}
 \boxed{\displaystyle
 \frac{b!}{d_b}\gcd\!\left(b,\frac{d_b}{b}\right)
 \mid\gcd(p_n,q_n).}
\end{equation}
If \(u^{(1)},u^{(2)},u^{(3)}\) are three solutions of
\eqref{eq:challenge-recurrence}, put, for \(n\ge-1\),
\[
 \mathcal W_n=\det\bigl(u^{(i)}_{n-j}\bigr)_{1\le i\le3,\,0\le j\le2},
\]
so in particular
\(\mathcal W_{-1}=\det(u^{(i)}_{-1-j})_{1\le i\le3,\,0\le j\le2}\).
Then, for every \(n\ge0\),
\begin{equation}\label{eq:challenge-casoratian}
 \boxed{\displaystyle
 \mathcal W_n=\mathcal W_{-1}
 \frac{A_{n+1}}{136}
 (n+1)!(n+2)!^4(n+3)!.}
\end{equation}
\end{proposition}

\begin{proof}
Write
\[
 W_{a,b}(k)=a!b!\binom ak\binom bk\frac1{k!}.
\]
Put \(g_b=\gcd(b,d_b/b)\).  After division by \(b!/d_b\), each summand of
\(q_n\) is \(b\) times an integer, because \(a!/k!\in\Z\).  The harmonic
part of \(p_n\) is likewise \(b\) times an integer, since
\(d_b(3H_k-H_{a-k}-H_{b-k})\in\Z\).  Its endpoint becomes \(d_b/b\).
Hence every term of both quotients is divisible by \(g_b\), proving
\eqref{eq:challenge-gcd}.

The companion determinant of \eqref{eq:challenge-recurrence} is
\(D_n/A_n\).  Since
\[
 D_n=(n+1)(n+2)^4(n+3)A_{n+1},
\]
iteration from \(-1\) to \(n\) telescopes and gives
\eqref{eq:challenge-casoratian}.
\end{proof}

\begin{remark}[Exterior-volume barrier]
The flat base field has determinant one, whereas scalar elimination together
with its factorial gauge produces the factorial-six Casoratian
\eqref{eq:challenge-casoratian}.  The forced divisor
\eqref{eq:challenge-gcd} is a lower bound for the common divisor, not an
upper bound.  These formulas record the factorial cost but do not rule out
further common-divisor cancellation; no small primitive integer linear form
is obtained here.
\end{remark}

\section{Exact Meijer--\texorpdfstring{\(G\)}{G}/hypergeometric duality}
\label{sec:meijer-duality}

The official solution of the Challenge is now public
\cite{RamanujanOfficial2026}.  It selects a recessive Meijer--\(G\) state by
Miller's algorithm, whereas \cref{sec:ramanujan-recurrence,sec:cmf} uses a
terminating \({}_2F_2\) state and a positive saddle.  The terminating module
is rationally gauge-equivalent to the contragredient of the normalized
official rank-three module.

Write
\[
 F_a(z)=F_{a,a+1}(z),\qquad
 \mathbf Y_a=
 \begin{pmatrix}F_a(1)\\ \theta F_a(1)\\ \theta^2F_a(1)\end{pmatrix}.
\]
The adjacent transition obtained by composing the two contiguous matrices
of \cref{thm:flat-cmf} is
\begin{equation}\label{eq:Delta-adjacent}
 \mathbf Y_a=\mathsf D_a\mathbf Y_{a+1},
\end{equation}
where
\begin{equation}\label{eq:D-adjacent}
\mathsf D_a=
\begin{pmatrix}
1&-\dfrac{2a+3}{(a+1)(a+2)}&\dfrac1{(a+1)(a+2)}\\[2mm]
1&\dfrac{a^2+a-1}{(a+1)(a+2)}&-\dfrac2{a+2}\\[2mm]
-(2a+1)&\dfrac{5a^2+11a+5}{(a+1)(a+2)}&\dfrac{a-2}{a+2}
\end{pmatrix}.
\end{equation}
For \(a\ge1\) this is the composite contiguous transfer.  The boundary
identity \(\mathbf Y_0=\mathsf D_0\mathbf Y_1\) is checked directly from the
terminating series.

Define
\begin{equation}\label{eq:official-Meijer-function}
 \mathcal H_a^G(z)=
 G_{2,3}^{3,2}\!\left(
 z\,\middle|\begin{matrix}1,1\\a+1,a+1,a+1\end{matrix}\right)
\end{equation}
with the Mellin--Barnes normalization of
\cite[\S5.2]{RamanujanOfficial2026}, and let
\[
 \mathbf S_a^G=
 \bigl(\mathcal H_a^G(1),\theta\mathcal H_a^G(1),
       \theta^2\mathcal H_a^G(1)\bigr)
\]
be the official row state, with the normalization of
\cite[\S5.2]{RamanujanOfficial2026}.  With \(b=a+1\), its transfer matrix is
\begin{equation}\label{eq:official-transfer}
 \mathbf S_{a+1}^G=\mathbf S_a^G\mathsf M_a^G,\qquad
 \mathsf M_a^G=
 \begin{pmatrix}
 0&b^3&b^3(3a+4)\\
 0&-3b^2&-b^2(8a+11)\\
 1&3b&6a^2+12a+5
 \end{pmatrix}.
\end{equation}
Put
\begin{equation}\label{eq:official-normalized-gauge}
 \mathsf P_a^G=b^{-2}\mathsf M_a^G,\qquad
 \mathsf B_a=
 \begin{pmatrix}
 b^2&-(2a+1)&1\\
 0&a/b&-1/b\\
 0&1/b&0
 \end{pmatrix}.
\end{equation}

\begin{theorem}[Exact rational \(SL_3\) coboundary]
\label{thm:ram-gauge}
For every \(a\ge0\),
\begin{equation}\label{eq:ram-determinants}
 \det\mathsf D_a=\det\mathsf B_a=\det\mathsf P_a^G=1,
 \qquad \det\mathsf M_a^G=(a+1)^6,
\end{equation}
and the exact gauge identity is
\begin{equation}\label{eq:ram-coboundary}
 \boxed{\displaystyle
 \mathsf D_a^{-\mathsf T}\mathsf B_{a+1}
 =\mathsf B_a(\mathsf P_a^G)^{-\mathsf T}.}
\end{equation}

If \(\mathbf V_a=\mathbf Y_a^{\mathsf T}\mathsf B_a\) and
\(\Pi_a^G=\mathsf M_0^G\cdots\mathsf M_{a-1}^G\), then
\begin{align}
 \mathbf V_{a+1}&=\mathbf V_a(\mathsf P_a^G)^{-\mathsf T},
 \label{eq:dual-trajectory}\\
 \boxed{\displaystyle
 \mathbf Y_a^{\mathsf T}\mathsf B_a
 =(a!)^2(1,-1,1)(\Pi_a^G)^{-\mathsf T}.}
 \label{eq:dual-trajectory-product}
\end{align}
The denominator coordinate has the precise factorial gauge
\begin{equation}\label{eq:q-dual-coordinate}
 \boxed{\displaystyle
 q_{a-3}=(a!)^2\mathbf V_ae_1.}
\end{equation}
Equivalently, \(\widehat{\mathbf V}_a=(a!)^2\mathbf V_a\) has first
coordinate \(q_{a-3}\) and evolves by the denominator-cleared dual matrix
\begin{equation}\label{eq:denominator-cleared-dual}
 \widehat{\mathbf V}_{a+1}
 =\widehat{\mathbf V}_a\widehat{\mathsf M}_a^G,\qquad
 \widehat{\mathsf M}_a^G=b^4(\mathsf M_a^G)^{-\mathsf T},\qquad
 \det\widehat{\mathsf M}_a^G=b^6.
\end{equation}
\end{theorem}

\begin{proof}
For \(a\ge1\), equation \eqref{eq:Delta-adjacent} is the product
\(M_b(a,a+2;1)M_a(a+1,a+2;1)\) from
\cref{thm:flat-cmf}; multiplication gives \eqref{eq:D-adjacent}.  At
\(a=0\), direct evaluation of the two terminating polynomials gives the
same identity.
The determinants in \eqref{eq:ram-determinants} are direct expansions.
Substitution of \cref{eq:D-adjacent,eq:official-normalized-gauge} into
\eqref{eq:ram-coboundary} reduces its nine entries to rational identities in
\(a\).

Rearranging \eqref{eq:ram-coboundary} gives
\(\mathsf D_a^{\mathsf T}\mathsf B_a
=\mathsf B_{a+1}(\mathsf P_a^G)^{\mathsf T}\), which proves
\eqref{eq:dual-trajectory}.  Since
\(\mathbf Y_0=(1,0,0)^{\mathsf T}\) and
\(\mathbf Y_0^{\mathsf T}\mathsf B_0=(1,-1,1)\), iteration gives
\eqref{eq:dual-trajectory-product}.  Finally,
\(\mathbf V_ae_1=b^2F_a(1)\), whereas
\eqref{eq:q-Fab-link} gives
\(q_{a-3}=(a!)^2b^2F_a(1)\).  This proves
\eqref{eq:q-dual-coordinate}; the last assertions follow by rescaling.
\end{proof}

The gauge fixes the finite connection pairing, not only its projective
direction.

\begin{theorem}[Finite Meijer--\(G\)/\({}_2F_2\) bilinear identity]
\label{thm:ram-bilinear}
For every \(a\ge0\),
\begin{equation}\label{eq:ram-bilinear-matrix}
 \boxed{\displaystyle
 \mathbf S_a^G\mathsf B_a^{\mathsf T}\mathbf Y_a=(a!)^2.}
\end{equation}
Writing \(b=a+1\), \(F=F_a(1)\), and
\(H=\mathcal H_a^G(1)\), this is
\begin{align}
 b^2HF
 &+(\theta H)\left(-(2a+1)F+\frac ab\theta F
                         +\frac1b\theta^2F\right)\notag\\
 &+(\theta^2H)\left(F-\frac1b\theta F\right)=(a!)^2.
 \label{eq:ram-bilinear-expanded}
\end{align}
Thus the normalized conserved primal--dual pairing is exactly one.
\end{theorem}

\begin{proof}
The official normalization gives
\(\mathbf S_0^G(1,-1,1)^{\mathsf T}=1\).  Use
\(\mathbf S_a^G=\mathbf S_0^G\Pi_a^G\) and
\eqref{eq:dual-trajectory-product}; the factors \(\Pi_a^G\) cancel.
Expanding \(\mathsf B_a^{\mathsf T}\mathbf Y_a\) gives
\eqref{eq:ram-bilinear-expanded}.
\end{proof}

The determinant-six phenomenon in \cref{prop:challenge-arithmetic} can now
be located exactly.  If \(\mathsf U(a)\) is the cyclic gauge of the official
solution, then
\begin{equation}\label{eq:cyclic-gauge-det}
 \det\mathsf U(a)=(a+1)^5(a+2)A_a,\qquad \det\mathsf U(0)=136.
\end{equation}
The coboundary transformation
\begin{equation}\label{eq:companion-gauge-definition}
 \mathsf C^{\mathrm{comp}}(a)=
 \mathsf U(a)^{-1}\widehat{\mathsf M}_a^G\mathsf U(a+1)
\end{equation}
gives the scalar companion transfer.  Consequently
\begin{align}
 \det\bigl(\mathsf C^{\mathrm{comp}}(0)\cdots
                 \mathsf C^{\mathrm{comp}}(n)\bigr)
 &=((n+1)!)^6\frac{\det\mathsf U(n+1)}{136}\notag\\
 &=\frac{A_{n+1}}{136}(n+1)!(n+2)!^4(n+3)!,
 \label{eq:casoratian-gauge-origin}
\end{align}
which is the multiplier of \(\mathcal W_{-1}\) in
\eqref{eq:challenge-casoratian}.  The sixth factorial power comes from the
denominator-cleared dual determinant; the shifted factorials and
\(A_{n+1}/136\) are the endpoint cyclic-gauge determinant.

\section{Euler--Beta--log transmutation and the offset ladder}
\label{sec:beta-log}

The adjacent sums are an exact integral transform of the cubic Pilehrood
polynomials, not merely an analogous family.  Define
\begin{align}
 \mathcal Q_m(t)&=\sum_{j=0}^m\binom mj^3j!\,t^j,
 \label{eq:cubic-Q}\\
 \mathcal P_m(t)&=\sum_{j=0}^m\binom mj^3j!
   \bigl(3H_{m-j}-2H_j\bigr)t^j.
 \label{eq:cubic-P}
\end{align}
These are the cubic polynomials occurring in the Euler-constant construction
of Hessami Pilehrood--Hessami Pilehrood
\cite{Pilehrood2013,PilehroodBell2012}.  For \(d\ge1\), put
\begin{equation}\label{eq:Beta-operator}
 \mathcal B_d[f]=\frac1{(d-1)!}\int_0^1(1-t)^{d-1}f(t)\,dt
\end{equation}
and
\begin{equation}\label{eq:transmutation-operator}
 \mathcal T_{m,d}[f](z)=
 \frac{(m+d)!}{m!^2}z^m\mathcal B_d[f(t/z)].
\end{equation}

\begin{theorem}[Exact Beta--log transmutation]
\label{thm:beta-log-transform}
For every \(m\ge0\) and \(d\ge1\), the first identity below holds for
\(z\ne0\) and extends polynomially to \(z=0\); at \(z=1\), the second
identity also holds:
\begin{align}
 F_{m,m+d}(z)&=\mathcal T_{m,d}[\mathcal Q_m](z),
 \label{eq:offset-F-transform}\\
 \mathcal G_{m,d}(1)&=
 \mathcal T_{m,d}[\mathcal P_m+\mathcal Q_m\log t](1),
 \label{eq:offset-G-transform}
\end{align}
where
\begin{equation}\label{eq:Gmd-def}
 \mathcal G_{m,d}(z)=
 \sum_{k=0}^m\binom mk\binom{m+d}k\frac{z^k}{k!}
 \bigl(3H_k-H_{m-k}-H_{m+d-k}\bigr).
\end{equation}
Moreover the transform intertwines Euler derivations:
\begin{equation}\label{eq:Beta-intertwining}
 \theta_z\mathcal T_{m,d}
 =\mathcal T_{m,d}(m-\theta_t).
\end{equation}

For the Challenge specialization \(m=a\), \(d=1\), \(z=1\),
\begin{align}
 q_{a-3}&=(a+1)^3a!\int_0^1\mathcal Q_a(t)\,dt,
 \label{eq:q-Beta-transform}\\
 p_{a-3}&=(a+1)^3a!
 \int_0^1\bigl(\mathcal P_a(t)+\mathcal Q_a(t)\log t\bigr)\,dt+a!.
 \label{eq:p-Beta-transform}
\end{align}
\end{theorem}

\begin{proof}
Put \(j=m-k\).  The beta integrals
\begin{align*}
 \mathcal B_d[t^j]&=\frac{j!}{(j+d)!},\\
 \mathcal B_d[t^j\log t]&=\frac{j!}{(j+d)!}(H_j-H_{j+d})
\end{align*}
turn the coefficient of \(z^{m-j}=z^k\) on the right of
\eqref{eq:offset-F-transform} into
\[
 \frac{m!(m+d)!}{k!^3(m-k)!(m+d-k)!},
\]
which is the coefficient of \(F_{m,m+d}\).  Adding
\(3H_{m-j}-2H_j\) and \(H_j-H_{j+d}\) gives exactly
\(3H_k-H_{m-k}-H_{m+d-k}\), proving
\eqref{eq:offset-G-transform} at \(z=1\).  Differentiating
\(z^mf(t/z)\) gives
\eqref{eq:Beta-intertwining}.  Finally
\cref{eq:q-Beta-transform,eq:p-Beta-transform} follow from
\cref{eq:qn-closed,eq:pn-closed}.
\end{proof}

\begin{corollary}[Fixed-offset Euler ladder]
\label{cor:offset-ladder}
For every fixed \(d\ge1\),
\begin{equation}\label{eq:offset-gamma-limit}
 \boxed{\displaystyle
 \frac{\mathcal G_{m,d}(1)}{F_{m,m+d}(1)}\longrightarrow\gamma.}
\end{equation}
\end{corollary}

\begin{proof}
Use the probability weights
\[
 \Prob(K=k)\propto\binom mk\binom{m+d}k\frac1{k!}.
\]
Their consecutive ratio is
\((m-k)(m+d-k)/(k+1)^3\), so, with
\(s_m=(m(m+d))^{1/3}\), the same ratio argument as in
\cref{thm:challenge-solution} gives \(K/s_m\to1\) exponentially in
probability.  On the window \(|K/s_m-1|\le\eta\),
\[
 3H_K-H_{m-K}-H_{m+d-K}
 =\gamma+\log\frac{K^3}{(m-K)(m+d-K)}+o(1)
 =\gamma+O(\eta)+o(1).
\]
The logarithmic growth on the complement is absorbed by the exponential
tail.  First let \(m\to\infty\), and then let \(\eta\downarrow0\).
Taking expectations proves \eqref{eq:offset-gamma-limit}.
\end{proof}

The multiple-Laguerre background and general approximation families are
already present in the cited literature and in
\cite{WolfsVanAssche2025}.  The assertion here is the exact transmutation
between that cubic system and the 2026 Challenge normalization, together
with its full fixed-offset form.

\section{Rational Gamma jets and contact polynomials}
\label{sec:gamma-jets}

The numerator and denominator are the first two jets of one rational
deformation.  Continue to write \(b=a+1\), and define
\begin{align}
 \mathscr J_a(s)&=a!b!\sum_{k=0}^a
 \frac1{(1-s)_k^3(1+s)_{a-k}(1+s)_{b-k}},
 \label{eq:Jas-def}\\
 \mathscr K_a(s)&=\mathscr J_a(s)+\frac{s}{b\,b!}.
 \label{eq:Kas-def}
\end{align}
The small final term is essential: it records the resonant endpoint
correction \(+a!\) in \eqref{eq:pn-closed}.

\begin{theorem}[Master Gamma-jet transform]
\label{thm:gamma-jet}
For every \(a\ge0\),
\begin{equation}\label{eq:zero-first-jets}
 \boxed{\displaystyle
 q_{a-3}=b\,a!b!\,\mathscr K_a(0),\qquad
 p_{a-3}=b\,a!b!\,\mathscr K_a'(0),}
\end{equation}
and hence
\begin{equation}\label{eq:gamma-first-log-jet}
 \frac{p_{a-3}}{q_{a-3}}=(\log\mathscr K_a)'(0).
\end{equation}
As \(a\to\infty\), locally uniformly on the strip \(|\Re s|<1\),
\begin{equation}\label{eq:Gamma-master-limit}
 \boxed{\displaystyle
 \frac{\mathscr J_a(s)}{\mathscr J_a(0)},\quad
 \frac{\mathscr K_a(s)}{\mathscr K_a(0)}
 \longrightarrow
 \Gamma(1-s)^3\Gamma(1+s)^2.}
\end{equation}
Consequently, for every \(m\ge2\),
\begin{equation}\label{eq:zeta-log-jets}
 \boxed{\displaystyle
 \frac{\left.\dfrac{d^m}{ds^m}\log\mathscr K_a(s)\right|_{s=0}}
 {(3+2(-1)^m)(m-1)!}\longrightarrow\zeta(m),}
\end{equation}
while the first logarithmic jet tends to \(\gamma\).
Every finite-stage logarithmic jet in \eqref{eq:zeta-log-jets} is rational.
\end{theorem}

\begin{proof}
Differentiating a reciprocal Pochhammer symbol at zero gives
\begin{align*}
 \left.\frac d{ds}\log(1-s)_k^{-3}\right|_{s=0}&=3H_k,\\
 \left.\frac d{ds}\log(1+s)_r^{-1}\right|_{s=0}&=-H_r.
\end{align*}
Equations \cref{eq:qn-closed,eq:pn-closed} now give
\eqref{eq:zero-first-jets}; the derivative of \(s/(bb!)\) produces exactly
the endpoint \(a!\).

Normalize the summands of \(\mathscr J_a(0)\) to the probability law
\eqref{eq:challenge-pmf}.  Uniformly for \(s\) in a compact subset of the
strip \(|\Re s|<1\), the ratio of a deformed summand to its value at zero is
 \begin{align}
 &\left(\frac{k!}{(1-s)_k}\right)^3
 \frac{(a-k)!}{(1+s)_{a-k}}
 \frac{(b-k)!}{(1+s)_{b-k}}\notag\\
 &\quad=\Gamma(1-s)^3\Gamma(1+s)^2
 k^{3s}(a-k)^{-s}(b-k)^{-s}(1+o(1)).
 \label{eq:uniform-gamma-ratio}
 \end{align}
Choose a shrinking central window
\(|k/(ab)^{1/3}-1|\le\eta_a\), where
\(\eta_a\downarrow0\) and \(\eta_a^2(ab)^{1/3}\to\infty\).
The concentration estimate \eqref{eq:challenge-concentration} makes its
complement exponentially small.  On the window,
\(k=(ab)^{1/3}(1+o(1))\), \(a-k=a(1+o(1))\), and
\(b-k=b(1+o(1))\); hence the power factor in
\eqref{eq:uniform-gamma-ratio} tends uniformly to one.  Uniform Gamma-ratio
bounds on each compact substrip give at most a fixed polynomial in
\(k+1,a-k+1,b-k+1\), which the exponential tail absorbs.  This proves the
locally uniform limit for \(\mathscr J_a\).  The extra term in
\(\mathscr K_a\) is \(O((bb!)^{-1})\), whereas
\(\mathscr J_a(0)\ge1\), so the same limit holds for \(\mathscr K_a\).

Finally,
\begin{equation}\label{eq:Gamma-log-series}
 \log\{\Gamma(1-s)^3\Gamma(1+s)^2\}
 =\gamma s+\sum_{m\ge2}\frac{3+2(-1)^m}{m}\zeta(m)s^m.
\end{equation}
Local uniformity permits termwise differentiation of logarithms near zero
and proves \eqref{eq:zeta-log-jets}.  Rationality at finite \(a\) follows
because \(\mathscr K_a\in\Q(s)\) and \(\mathscr K_a(0)\ne0\).
\end{proof}

For example, as \(a\to\infty\), the correctly normalized first cases are
\begin{equation}\label{eq:first-zeta-jets}
 \frac{(\log\mathscr K_a)''(0)}5\to\zeta(2),\qquad
 \frac{(\log\mathscr K_a)^{(3)}(0)}2\to\zeta(3),\qquad
 \frac{(\log\mathscr K_a)^{(4)}(0)}{30}\to\zeta(4).
\end{equation}
The raw derivative \(\mathscr K_a''(0)\) cannot replace the second
logarithmic derivative.

\begin{corollary}[Exact parity separation]
\label{cor:gamma-jet-parity}
As \(a\to\infty\), locally uniformly on \(|\Re s|<1\),
\begin{align}
 \frac{\mathscr K_a(s)\mathscr K_a(-s)}{\mathscr K_a(0)^2}
 &\longrightarrow
 \left(\frac{\pi s}{\sin\pi s}\right)^5,
 \label{eq:even-jet-limit}\\
 \frac{\mathscr K_a(s)}{\mathscr K_a(-s)}
 &\longrightarrow\frac{\Gamma(1-s)}{\Gamma(1+s)}.
 \label{eq:odd-jet-limit}
\end{align}
Thus the logarithm of the first quotient generates only even zeta values,
while the logarithm of the second generates \(\gamma\) and the odd zeta
values.
\end{corollary}

\begin{proof}
Multiply and divide the two limits in \eqref{eq:Gamma-master-limit}, and use
\(\Gamma(1-s)\Gamma(1+s)=\pi s/\sin\pi s\).
\end{proof}

The Bell-polynomial zeta tower itself is classical
\cite{PilehroodBell2012}.  The point of
\cref{thm:gamma-jet,cor:gamma-jet-parity} is its exact finite connection to
the individual 2026 Challenge approximants, including their endpoint term.

The same deformation has a finite integer spectral avatar.  Put
\begin{equation}\label{eq:contact-denominator}
 \mathscr D_a(s)=
 \prod_{j=1}^a(s-j)^3(s+j)^2(s+a+1).
\end{equation}

\begin{theorem}[Monic integer contact polynomial]
\label{thm:contact-polynomial}
There is a unique polynomial \(\mathscr N_a\in\Z[s]\) such that
\begin{equation}\label{eq:contact-rational-form}
 \mathscr K_a(s)=
 \frac{\mathscr N_a(s)}{a!(a+1)^2\mathscr D_a(s)}.
\end{equation}
It is monic of degree \(5a+2\), and
\begin{equation}\label{eq:contact-degree}
 \deg\{\mathscr N_a-s\mathscr D_a\}=3a.
\end{equation}
Consequently \(\mathscr N_a\) and \(s\mathscr D_a\) have the same first
\(2a+1\) Newton power sums.  Explicitly, if the zeros of
\(\mathscr N_a\) are counted with multiplicity, then for
\(1\le m\le2a+1\),
\begin{equation}\label{eq:contact-Newton}
 \sum_{\mathscr N_a(\nu)=0}\nu^m
 =3\sum_{j=1}^aj^m+2\sum_{j=1}^a(-j)^m+(-a-1)^m.
\end{equation}
At the opposite end of the polynomial,
\begin{align}
 \mathscr N_a(0)
 &=(-1)^a(a+1)(a!)^4q_{a-3},
 \label{eq:contact-N0}\\
 \mathscr N_a'(0)
 &=(-1)^a(a+1)(a!)^4
 \left[p_{a-3}-q_{a-3}\left(H_a-\frac1{a+1}\right)\right].
 \label{eq:contact-N1}
\end{align}
\end{theorem}

\begin{proof}
Multiplication of \eqref{eq:Kas-def} by
\(a!(a+1)^2\mathscr D_a(s)\) clears every Pochhammer denominator.  Each
resulting quotient is in \(\Z[s]\); the endpoint term becomes exactly
\(s\mathscr D_a(s)\).  The latter is monic of degree \(5a+2\), while the
\(k\)-th summand of the remaining part has degree \(3a-k\).  The \(k=0\)
term has nonzero leading coefficient, proving monicity and
\eqref{eq:contact-degree}.

Two monic degree-\(5a+2\) polynomials whose difference has degree \(3a\)
share their first \(2a+1\) coefficients below the leading term.  Newton's
identities give \eqref{eq:contact-Newton}; the roots of
\(s\mathscr D_a\) are \(0\), three copies of \(1,\ldots,a\), two copies of
\(-1,\ldots,-a\), and \(-a-1\).  Finally evaluate
\eqref{eq:contact-rational-form} and its logarithmic derivative at zero,
use \eqref{eq:zero-first-jets}, and compute
\[
 \frac{\mathscr D_a'(0)}{\mathscr D_a(0)}
 =-H_a+\frac1{a+1}.
\]
This gives \cref{eq:contact-N0,eq:contact-N1}.
\end{proof}

\section{Probabilistic realizations and zero geometry}
\label{sec:ram-probability}

The positive law in \eqref{eq:challenge-pmf} has two complementary
realizations: exponential maxima at finite index and independent Bernoulli
variables through the zeros of its generating polynomial.

For \(m\ge1\), let \(M_m\) be the maximum of \(m\) independent
\(\operatorname{Exp}(1)\) variables, put \(M_0=0\), and use independent
copies whenever superscripts are present.  Then
\begin{equation}\label{eq:maxima-mgf}
 \mathbb Ee^{tM_m}=\prod_{j=1}^m\frac j{j-t}
 =\frac{m!}{(1-t)_m}\qquad(\Re t<1).
\end{equation}

\begin{theorem}[Exact finite maxima model]
\label{thm:ram-probability}
For \(a\ge0\), let \(J_a\) have the law
\begin{equation}\label{eq:Ja-law}
 \Prob(J_a=k)=
 \frac{\binom ak\binom{a+1}k/k!}{F_{a,a+1}(1)}
 \qquad(0\le k\le a),
\end{equation}
and, conditionally on \(J_a\), define
\begin{equation}\label{eq:Za0}
 Z_a^{(0)}=
 M_{J_a}^{(1)}+M_{J_a}^{(2)}+M_{J_a}^{(3)}
 -M_{a-J_a}^{(4)}-M_{a+1-J_a}^{(5)}.
\end{equation}
Then
\begin{equation}\label{eq:exact-maxima-mgf}
 \boxed{\displaystyle
 \mathbb Ee^{tZ_a^{(0)}}
 =\frac{\mathscr J_a(t)}{\mathscr J_a(0)}}
 \qquad(|\Re t|<1).
\end{equation}
If
\begin{equation}\label{eq:Za-correction}
 c_a=\frac{a!}{q_{a-3}},\qquad Z_a=Z_a^{(0)}+c_a,
\end{equation}
then
\begin{equation}\label{eq:finite-mean-approximant}
 \boxed{\displaystyle\mathbb EZ_a=\frac{p_{a-3}}{q_{a-3}}.}
\end{equation}
Moreover, as \(a\to\infty\), with independent standard Gumbel variables
\(G_1,\ldots,G_5\),
\begin{equation}\label{eq:Za-all-moments}
 Z_a\longrightarrow X_0:=G_1+G_2+G_3-G_4-G_5
\end{equation}
in distribution and with convergence of every moment.  Equivalently,
locally uniformly for \(|\Re t|<1\),
\begin{equation}\label{eq:Za-mgf-limit}
 \mathbb Ee^{tZ_a}\longrightarrow
 \Gamma(1-t)^3\Gamma(1+t)^2.
\end{equation}
Thus
\begin{align}
 \kappa_1(Z_a)&\longrightarrow\gamma,\notag\\
 \kappa_m(Z_a)&\longrightarrow
 (m-1)!\{3+2(-1)^m\}\zeta(m)\qquad(m\ge2).
 \label{eq:Za-cumulants}
\end{align}
\end{theorem}

\begin{proof}
Multiply the weight in \eqref{eq:Ja-law} by the five conditional moment
generating functions from \eqref{eq:maxima-mgf}.  All ordinary factorials
cancel, leaving exactly the summand of \(\mathscr J_a(t)\); normalization at
zero proves \eqref{eq:exact-maxima-mgf}.  Taking the first derivative and
using \eqref{eq:zero-first-jets} gives
\eqref{eq:finite-mean-approximant}.  Since \(c_a\to0\),
\cref{thm:gamma-jet} gives \eqref{eq:Za-mgf-limit} on a neighborhood of the
origin.  Uniform exponential moments on every smaller strip imply
distributional and all-moment convergence.  Finally differentiate the
logarithm and use \eqref{eq:Gamma-log-series}.
\end{proof}

The distinction between the two deformations is important:
\begin{equation}\label{eq:mgf-not-K}
 \mathbb Ee^{tZ_a}
 =e^{c_at}\frac{\mathscr J_a(t)}{\mathscr J_a(0)}
 \ne\frac{\mathscr K_a(t)}{\mathscr K_a(0)}
\end{equation}
at finite \(a\).  The function \(\mathscr K_a\) is the exact rational
zero/first-jet carrier; \(\mathscr J_a\) is the exact probability carrier.
They have the same Gamma limit.

We next identify the full official Meijer state.  Let
\begin{equation}\label{eq:Xn-def}
 A_{n,i}\sim\Gamma(n+1,1)\quad(1\le i\le3),\qquad
 E_4,E_5\sim\operatorname{Exp}(1)
\end{equation}
be independent, and put
\begin{equation}\label{eq:Xn-ratio}
 X_n=\log\frac{E_4E_5}{A_{n,1}A_{n,2}A_{n,3}}.
\end{equation}

\begin{theorem}[Stop-loss interpretation of the Meijer state]
\label{thm:meijer-probability}
For the official normalization in \cref{sec:meijer-duality},
\begin{equation}\label{eq:Meijer-stop-loss}
 \boxed{\displaystyle
 \mathcal H_n^G(e^{-x})=(n!)^3\mathbb E(X_n-x)_+}
 \qquad(x\in\R).
\end{equation}
Consequently
\begin{equation}\label{eq:Meijer-state-probability}
 \boxed{\displaystyle
 \frac{\bigl(\mathcal H_n^G(1),\theta\mathcal H_n^G(1),
 \theta^2\mathcal H_n^G(1)\bigr)}{(n!)^3}
 =\bigl(\mathbb E(X_n)_+,\Prob(X_n>0),f_{X_n}(0)\bigr).}
\end{equation}
The functional shift is
\begin{equation}\label{eq:Meijer-functional-shift}
 \mathcal H_{n+1}^G(z)=z\theta^2\mathcal H_n^G(z),
\end{equation}
and hence
\begin{equation}\label{eq:probability-conservation}
 f_{X_n}(x)=e^x(n+1)^3\mathbb E(X_{n+1}-x)_+.
\end{equation}

At \(n=0\),
\begin{equation}\label{eq:h-constant}
 \mathfrak h:=\mathbb E(X_0)_+
 =\int_0^\infty e^{-x}\frac{\log x}{x-1}\,dx
 =1.4322057346532244148\ldots,
\end{equation}
whereas
\begin{equation}\label{eq:Powers-gamma}
 \mathbb EX_0=\Prob(X_0>0)=\gamma.
\end{equation}
In particular \(\mathfrak h\ne\gamma\).
\end{theorem}

\begin{proof}
The Mellin transform of \(X_n\) is
\begin{equation}\label{eq:Xn-mgf}
 \mathbb Ee^{sX_n}
 =\frac{\Gamma(1+s)^2\Gamma(n+1-s)^3}{(n!)^3},
 \qquad -1<\Re s<n+1.
\end{equation}
Perron's inverse Mellin formula on a line \(0<c<n+1\) gives
\eqref{eq:Meijer-stop-loss}; differentiation in \(x\), using
\(\theta=-d/dx\) at \(z=e^{-x}\), gives
\eqref{eq:Meijer-state-probability}.  Shifting the three gamma factors in
the Mellin--Barnes integral proves \eqref{eq:Meijer-functional-shift}, and
comparison with \eqref{eq:Meijer-stop-loss} gives
\eqref{eq:probability-conservation}.  At \(n=0\), substitute the five
exponential densities into \(\mathbb E(X_0)_+\), integrate four variables
first, and apply Frullani's identity; the remaining integral is
\[
 \int_0^\infty e^{-x}\frac{\log x}{x-1}\,dx,
\]
which proves \eqref{eq:h-constant}.
The probability identity in \eqref{eq:Powers-gamma} is due to Powers
\cite{Powers2023}; the expectation identity is immediate from
\(\mathbb EG_j=\gamma\).
\end{proof}

The Bernoulli realization comes from real-rootedness.  For integers
\(b\ge a\ge1\), write the zeros of \(F_{a,b}\), with multiplicities, as
\(-\rho_{a,b,1},\ldots,-\rho_{a,b,a}\).

\begin{theorem}[Poisson--binomial law and zero-density limit]
\label{thm:poisson-binomial}
All zeros of \(F_{a,b}\) are negative real numbers, and
\begin{equation}\label{eq:Poisson-binomial-factor}
 \frac{F_{a,b}(z)}{F_{a,b}(1)}
 =\prod_{j=1}^a(1-p_j+p_jz),
 \qquad p_j=\frac1{1+\rho_{a,b,j}}.
\end{equation}
Thus the law with weights proportional to
\(\binom ak\binom bk/k!\) is exactly a sum of independent Bernoulli
variables.  Its mass sequence is strictly log-concave and ultra-log-concave,
and it is monotone in each of \(a,b\) in the monotone-likelihood-ratio order.

Fix \(d\ge0\), put \(b=a+d\) and \(s_a=(ab)^{1/3}\).  As \(a\to\infty\),
if \(J_a\) has these weights, then
\begin{align}
 \mathbb EJ_a&\sim s_a,&
 \Var(J_a)&\sim\frac{s_a}{3},
 \label{eq:Ja-mean-var}\\
 \frac{J_a-\mathbb EJ_a}{\sqrt{\Var(J_a)}}&\Longrightarrow N(0,1).
 \label{eq:Ja-CLT}
\end{align}
The corresponding local central limit theorem holds, and \(J_a/s_a\)
satisfies a large-deviation principle with speed \(s_a\) and rate
\begin{equation}\label{eq:Ja-rate}
 I(x)=
 \begin{cases}
 3(x\log x-x+1),&x\ge0,\\
 +\infty,&x<0,
 \end{cases}
 \qquad(0\log0:=0).
\end{equation}

Finally, the locally finite zero measures converge vaguely on
\([0,\infty)\):
\begin{equation}\label{eq:zero-density-limit}
 \boxed{\displaystyle
 \frac1{s_a}\sum_{j=1}^a\delta_{\rho_{a,b,j}}
 \Longrightarrow
 \frac{\sqrt3}{2\pi}\rho^{-2/3}\,d\rho.}
\end{equation}
\end{theorem}

\begin{proof}
The polynomial \(L_a(-z)=\sum_k\binom akz^k/k!\) has only negative real
zeros.  The coefficient sequence \(\{\binom bk\}_{k\ge0}\) is a classical
P\'olya--Schur multiplier sequence because its exponential generating
function is \(L_b(-z)\); see \cite{PolyaSchur1914}.  Applying the multiplier
theorem to \(L_a(-z)\)
proves real-rootedness, and positive coefficients make every root negative.
Factorization and normalization at \(z=1\) give
\eqref{eq:Poisson-binomial-factor}.

The consecutive mass ratio
\[
 \frac{w_{k+1}}{w_k}=\frac{(a-k)(b-k)}{(k+1)^3}
\]
is strictly decreasing, proving strict log-concavity.  Ultra-log-concavity
follows because \(w_k/\binom ak=\binom bk/k!\) has decreasing consecutive
ratio.  Finally
\[
 \frac{w_{a+1,b;k}}{w_{a,b;k}}=\frac{a+1}{a+1-k}
\]
is increasing in \(k\), and similarly for \(b\), proving the likelihood
ratio assertion.

A Taylor expansion about the unique real saddle
\(\kappa_a\), determined by
\[
 (a-\kappa_a)(b-\kappa_a)=(\kappa_a+1)^3,
\]
shows that the discrete law has one mode or two adjacent modes neighboring
\(\kappa_a\).  Its logarithmic curvature is
\[
 \frac1{a-\kappa_a}+\frac1{b-\kappa_a}
 +\frac3{\kappa_a+1}\sim\frac3{s_a},
\]
which gives \eqref{eq:Ja-mean-var}.  The Bernoulli representation,
\(\Var(J_a)\to\infty\), and the standard Poisson--binomial local limit
theorem \cite[Chapter~VII]{Petrov1975} give the central and local limit
assertions.  Stirling's
formula at \(k=xs_a\), uniformly on compact \(x\)-intervals, gives
\eqref{eq:Ja-rate} and exponential tightness gives the full large-deviation
principle.

For \(z>0\), tilt the weights by \(z^k\).  Their ratio has the extra factor
\(z\), so \(J_{a,z}/s_a\to z^{1/3}\).  Therefore
\begin{equation}\label{eq:Stieltjes-transform-limit}
 \frac1{s_a}\frac{F'_{a,b}(z)}{F_{a,b}(z)}
 =\frac{\mathbb E_zJ_{a,z}}{s_az}\longrightarrow z^{-2/3}.
\end{equation}
The left side is the Stieltjes transform of the measure on the left of
\eqref{eq:zero-density-limit}, whereas
\[
 \int_0^\infty\frac1{\rho+z}
 \frac{\sqrt3}{2\pi}\rho^{-2/3}\,d\rho=z^{-2/3}.
\]
Uniqueness and the standard continuity theorem for Stieltjes transforms
\cite[Chapter~I]{Akhiezer1965} give the vague limit.
\end{proof}

No simplicity or strict-interlacing assertion is needed above.  Those
stronger statements require a separate strict multiplier or
multiple-orthogonality input and are not claimed here.

\begin{remark}[Priority boundary for the Ramanujan package]
The official solution already proves Question~2 through a rank-three
Meijer--\(G\) trajectory, a Mellin--Barnes recessive estimate, a rational
companion gauge, and inverse-transpose Miller selection
\cite{RamanujanOfficial2026}.  The cubic sums and their Bell--zeta
approximants occur in \cite{PilehroodBell2012,Pilehrood2013}, while the
five-exponential ratio law, its signed-Gumbel logarithm, and the associated
Gamma-product transform occur in \cite{Powers2023}.  The general
multiple-orthogonal and multiple-Laguerre background is developed in
\cite{VanAsscheWolfs2023,WolfsVanAssche2025}.  The contribution of the
present sections is the explicit finite-level transmutation among the
Challenge recurrence, its terminating \({}_2F_2\) dual, the official
Meijer--\(G\) trajectory, and the cubic Bell family, together with the exact
finite probability carrier and the all-order logarithmic-jet limits in the
Challenge normalization.
\end{remark}

\section{The regularized Gamma determinant}
\label{sec:gamma-determinant}

The positive trace-class operator of \cref{sec:spectral} encodes the single
Euler trace \(S_2\).  A different, signed Hilbert--Schmidt operator encodes
the complete limiting Gamma jet.  On
\(\ell^2(\N)^{\oplus5}\), define
\begin{equation}\label{eq:A-Gamma}
 \mathsf A_\Gamma=\diag\left(
 \frac1n,\frac1n,\frac1n,-\frac1n,-\frac1n:n\ge1\right).
\end{equation}
Let \(\mathsf A_{\Gamma,N}\) be its compression to the first \(N\)
coordinates in each of the five summands.  We use the standard
Hilbert--Schmidt determinant
\[
 \det{}_2(I-T)=\det\bigl((I-T)e^T\bigr).
\]

\begin{theorem}[Regularized determinant and zeta traces]
\label{thm:gamma-det2}
One has
\(\mathsf A_\Gamma\in\mathcal S_2\setminus\mathcal S_1\), and
\begin{equation}\label{eq:FP-trace-gamma}
 \boxed{\displaystyle
 \FPTr\mathsf A_\Gamma
 :=\lim_{N\to\infty}
 \left(\Tr\mathsf A_{\Gamma,N}-\log N\right)=\gamma.}
\end{equation}
For every \(m\ge2\),
\begin{equation}\label{eq:Gamma-operator-traces}
 \boxed{\displaystyle
 \Tr\mathsf A_\Gamma^m
 =\{3+2(-1)^m\}\zeta(m).}
\end{equation}
Equivalently,
\begin{equation}\label{eq:Gamma-det2}
 \boxed{\displaystyle
 \Gamma(1-t)^3\Gamma(1+t)^2
 =e^{\gamma t}\det{}_2(I-t\mathsf A_\Gamma)^{-1}.}
\end{equation}
In particular
\begin{equation}\label{eq:Gamma-even-odd-traces}
 \Tr\mathsf A_\Gamma^{2r}=5\zeta(2r),\qquad
 \Tr\mathsf A_\Gamma^{2r+1}=\zeta(2r+1)\qquad(r\ge1).
\end{equation}
\end{theorem}

\begin{proof}
Square summability and failure of absolute summability are immediate from
the spectrum.  The truncated trace is \(H_N\), so
\eqref{eq:FP-trace-gamma} is Euler's definition.  Taking powers of the five
eigenvalue strings gives \eqref{eq:Gamma-operator-traces}.  Finally, for a
Hilbert--Schmidt operator,
\[
 \log\det{}_2(I-t\mathsf A_\Gamma)
 =-\sum_{m\ge2}\frac{t^m}{m}\Tr\mathsf A_\Gamma^m.
\]
Compare with \eqref{eq:Gamma-log-series}; this first proves
\eqref{eq:Gamma-det2} for \(|t|<1\), and meromorphic continuation gives the
identity on \(\C\).
\end{proof}

This regularized determinant is an explicit diagonal spectral realization of the
Gamma product, but it does not furnish a new irrationality result for odd
zeta values.  It must also not be confused with the ordinary Fredholm
determinant \(\Delta_2\) of the positive trace-class operator in
\cref{thm:Delta2}.

\section{Cross-place Euler--Pad\'e convergence}
\label{sec:cross-place-pade}

The Pad\'e system for Euler's factorial series gives a separate arithmetic
transfer.  Put
\begin{equation}\label{eq:Euler-Pade-Q}
 Q_n(T)=\sum_{h=0}^nh!\binom nh^2(-T)^h,
\end{equation}
and let \(P_n(T)\), of degree at most \(n-1\), be the Pad\'e numerator
characterized by
\begin{equation}\label{eq:Euler-Pade-condition}
 Q_n(T)\sum_{m\ge0}m!T^m-P_n(T)=O(T^{2n})
\end{equation}
as a formal series.  For a prime \(p\), the series
\begin{equation}\label{eq:p-adic-Euler-series}
 E_p(t)=\sum_{m\ge0}m!t^m
\end{equation}
converges for \(|t|_p\le1\).  At a real negative argument put
\begin{equation}\label{eq:real-Euler-function}
 E_\infty(t)=\int_0^\infty\frac{e^{-x}}{1-tx}\,dx.
\end{equation}
The standard Pad\'e remainder gives, for \(t\in\Z_p\),
\begin{equation}\label{eq:p-adic-remainder-bound}
 v_p\!\left(E_p(t)-\frac{P_n(t)}{Q_n(t)}\right)
 \ge2v_p(n!)-v_p(Q_n(t)),
\end{equation}
whenever \(Q_n(t)\ne0\); see
\cite{ErnvallMatalaSeppala2022}.

\begin{theorem}[Unit-boundary residue criterion]
\label{thm:p-adic-unit-criterion}
Let \(p\) be prime and \(t\in\Z_p^\times\).  For
\(0\le r<p\), define
\begin{equation}\label{eq:q-rp}
 q_{r,p}(X)=\sum_{h=0}^rh!\binom rh^2(-X)^h\in\mathbb F_p[X].
\end{equation}
If
\begin{equation}\label{eq:good-residue-condition}
 q_{r,p}(\bar t)\ne0\qquad(0\le r<p),
\end{equation}
then every \(Q_n(t)\) is a \(p\)-adic unit and
\begin{equation}\label{eq:full-p-adic-convergence}
 \boxed{\displaystyle
 \frac{P_n(t)}{Q_n(t)}\longrightarrow E_p(t).}
\end{equation}
The good nonzero residue classes for the first odd primes are
\begin{equation}\label{eq:good-residue-table}
 \begin{array}{c|c}
 p&t\pmod p\\ \hline
 3&2\\
 5&2,3,4\\
 7&4,5\\
 13&6,9,11,12
 \end{array}.
\end{equation}
In particular, at \(t=-1\), the full convergent sequence tends to
\(E_p(-1)\) for \(p=2,3,5,13\).
\end{theorem}

\begin{proof}
Modulo \(p\), all terms with \(h\ge p\) vanish and, for \(h<p\),
\(\binom nh\bmod p\) depends only on \(n\bmod p\).  Hence
\begin{equation}\label{eq:Q-period-p}
 Q_n(t)\equiv q_{r,p}(\bar t)\pmod p
 \qquad(n\equiv r\pmod p).
\end{equation}
Condition \eqref{eq:good-residue-condition} makes every denominator a unit;
then \eqref{eq:p-adic-remainder-bound} tends to infinity.  The table is a
finite calculation in the indicated residue fields.

For \(p=2,t=-1\), an elementary calculation modulo four gives
\begin{equation}\label{eq:Q-v2}
 v_2(Q_n(-1))=
 \begin{cases}0,&n\text{ even},\\1,&n\text{ odd}.
 \end{cases}
\end{equation}
The same remainder estimate therefore proves full convergence also at two.
For \(p=3,5,13\), the assertion is contained in the table.
\end{proof}

The congruence behind a simultaneous subsequence is stronger than periodicity.

\begin{theorem}[A common real and all-prime subsequence]
\label{thm:common-placewise-subsequence}
For all integers \(n\ge0\) and \(k\ge1\),
\begin{equation}\label{eq:Carlitz-congruence}
 Q_{n+k}(T)\equiv Q_n(T)\pmod{k\Z[T]}.
\end{equation}
Let \(N_j=\lcm(1,\ldots,j)\).  For every fixed negative integer \(t\), the
single rational sequence
\begin{equation}\label{eq:common-rational-sequence}
 x_j(t)=\frac{P_{N_j}(t)}{Q_{N_j}(t)}
\end{equation}
satisfies
\begin{equation}\label{eq:common-placewise-limits}
 \boxed{\displaystyle
 x_j(t)\longrightarrow E_\infty(t)\quad\text{in }\R,
 \qquad
 x_j(t)\longrightarrow E_p(t)\quad\text{in }\Q_p
 \text{ for every fixed }p.}
\end{equation}
At \(t=-1\), the real limit is the Euler--Gompertz constant
\(\delta=eE_1(1)\), while the \(p\)-adic limit is \(E_p(-1)\).
\end{theorem}

\begin{proof}
For each \(h\), the polynomial
\(h!\binom Xh=X(X-1)\cdots(X-h+1)\) belongs to \(\Z[X]\).  Hence
\[
 h!\left\{\binom{n+k}h^2-\binom nh^2\right\}
 =\left\{h!\binom{n+k}h-h!\binom nh\right\}
  \left\{\binom{n+k}h+\binom nh\right\}
\]
is divisible by \(k\), coefficient by coefficient in
\eqref{eq:Euler-Pade-Q}.  This proves \eqref{eq:Carlitz-congruence}.
Taking \(n=0,k=N_j\) gives
\begin{equation}\label{eq:QN-unit-all-p}
 Q_{N_j}(T)\equiv1\pmod{N_j\Z[T]}.
\end{equation}
Thus for every fixed \(p\) and all \(j\ge p\), \(Q_{N_j}(t)\) is a
\(p\)-adic unit, and \eqref{eq:p-adic-remainder-bound} proves the \(p\)-adic
limit.  The real Stieltjes--Pad\'e convergence to
\eqref{eq:real-Euler-function} holds for every \(t<0\), and therefore along
the same subsequence.
\end{proof}

\begin{proposition}[\(p\)-adic denominator interpolation]
\label{prop:p-adic-denominator-interpolation}
For every prime \(p\), the series
\begin{equation}\label{eq:q-p-interpolation}
 q_p(x)=\sum_{h\ge0}h!\binom xh^2\qquad(x\in\Z_p)
\end{equation}
converges uniformly and defines a \(1\)-Lipschitz function satisfying
\(q_p(n)=Q_n(-1)\) for \(n\in\N_0\).  Moreover
\begin{equation}\label{eq:q-p-root-criterion}
 \boxed{\displaystyle
 \sup_{n\ge0}v_p(Q_n(-1))=\infty
 \Longleftrightarrow q_p\text{ has a zero in }\Z_p.}
\end{equation}
In particular, if \(q_p\) is rootless then the full Euler convergent
sequence at \(t=-1\) tends to \(E_p(-1)\).
\end{proposition}

\begin{proof}
Uniform convergence follows from \(v_p(h!)\to\infty\) and
\(\binom xh\in\Z_p\).  For each summand,
\[
 h!\left\{\binom xh^2-\binom yh^2\right\}
 =\left(x^{\underline h}-y^{\underline h}\right)
  \left\{\binom xh+\binom yh\right\},
\]
which is divisible by \(x-y\); passing to the uniform limit gives the
Lipschitz bound.  The interpolation identity follows because
\(\binom nh=0\) for \(h>n\).  If the valuations are unbounded, compactness
of \(\Z_p\) supplies a convergent subsequence of the corresponding integers,
and continuity gives a zero.  Conversely, approximate a zero by
nonnegative integers and use continuity.  If there is no zero, compactness
bounds \(v_p(q_p(x))\), and \eqref{eq:p-adic-remainder-bound} proves the
last assertion.
\end{proof}

The word ``simultaneous'' in \cref{thm:common-placewise-subsequence} is
placewise: no convergence in the restricted adelic product topology is
asserted.  Nor does the theorem prove irrationality of any \(E_p(-1)\).
At the residue \(n=p-1\),
\begin{equation}\label{eq:Kurepa-contact}
 Q_{p-1}(-1)\equiv\sum_{h=0}^{p-1}h!\pmod p,
\end{equation}
so Kurepa's left-factorial problem appears as one boundary congruence, but is
not solved here.  Full convergence at \(t=-1\) for every prime remains open;
the observed growth bound \(v_p(Q_n(-1))=O(\log n)\) is not presently a
theorem.  The existence of a zero of \(q_p\) means only that denominator
valuations are unbounded; it is not equivalent to failure of convergence,
which would require comparison with \(v_p(n!)\).

\section{The critical moving-point barrier}
\label{sec:barrier}

We first separate the coefficient of \(\gamma\) from the gamma-free tail.
Fix \(q>1\), put
\[
 f_q(t)=E_1(q^t),
\]
and consider a finitely supported coefficient vector \((c_n)\).  The two
linear statistics
\[
 \omega(c)=\sum_nc_n,
 \qquad
 \mu(c)=\sum_nnc_n
\]
give the exact decomposition
\begin{equation}\label{eq:augmentation-decomposition}
 \sum_nc_n\Ein(q^n)
 =\omega(c)\gamma+\mu(c)\log q+\sum_nc_nf_q(n).
\end{equation}
We call \(\omega=1,\mu=0\) the augmentation-one, log-balanced stratum.
The dyadic level belongs to this stratum.  The next theorem completely solves
its natural cumulative-positive cone.

Direct differentiation gives
\begin{equation}\label{eq:fq-convex}
 f_q'(t)=-(\log q)e^{-q^t}<0,
 \qquad
 f_q''(t)=(\log q)^2q^te^{-q^t}>0.
\end{equation}

\begin{theorem}[Sharp cumulative-positive extremality]
\label{thm:augmentation-extremality}
Let \(M\ge2\), and suppose that real coefficients \(c_0,\ldots,c_M\)
satisfy
\[
 \sum_{n=0}^Mc_n=1,
 \qquad
 \sum_{n=0}^Mnc_n=0.
\]
Put \(C_k=\sum_{n=0}^kc_n\), and assume \(C_k\ge0\) for \(0\le k<M\).
Then the complete sharp range is
\begin{equation}\label{eq:augmentation-range}
 \begin{split}
 M f_q(M-1)-(M-1)f_q(M)
 &\le \sum_{n=0}^Mc_nf_q(n)\\
 &\le f_q(M)+M\{f_q(0)-f_q(1)\}.
 \end{split}
\end{equation}
The lower endpoint is attained uniquely by
\begin{equation}\label{eq:augmentation-minimizer}
 c_{M-1}=M,\qquad c_M=-(M-1),
 \qquad c_0=\cdots=c_{M-2}=0.
\end{equation}
The upper endpoint is attained uniquely by
\(c_0=M,c_1=-M,c_M=1\), with the other coefficients zero.
Consequently the moving level \(Y_{M-1}(q)-\gamma\) is the unique minimum
tail in this cone, and
\[
 \min\sum_nc_nf_q(n)
 \sim \frac{M}{q^{M-1}}e^{-q^{M-1}}.
\]
\end{theorem}

\begin{proof}
Abel summation and the moment normalization give
\begin{align*}
 \sum_{n=0}^Mc_nf_q(n)
 &=f_q(M)+\sum_{k=0}^{M-1}C_k\{f_q(k)-f_q(k+1)\},\\
 \sum_{k=0}^{M-1}C_k&=M.
\end{align*}
By strict convexity, the positive differences
\(f_q(k)-f_q(k+1)\) decrease strictly with \(k\).  The displayed sum is
therefore minimized by putting all mass \(M\) at \(k=M-1\), and maximized
by putting it at \(k=0\).  Unwinding the partial sums gives the two unique
coefficient vectors.  The feasible \(C\)-vectors form the simplex
\(\{C_k\ge0:\sum C_k=M\}\); its continuous linear image is the whole
interval between these two endpoints, which proves the asserted complete
range.  The asymptotic is \eqref{eq:E1-asymptotic}.
\end{proof}

The same doubly exponential scale survives when the two levels use
multiplicatively independent bases.

\begin{proposition}[Two-base moving-level obstruction]
\label{prop:two-base-obstruction}
For \(q>1\), write
\[
 Y_N(q)=(N+1)\Ein(q^N)-N\Ein(q^{N+1})
       =\gamma+\varepsilon_N(q).
\]
Let
\[
 M_N=\left\lfloor\frac{N\log2}{\log3}\right\rfloor+1,
 \qquad W_N=Y_N(2)-Y_{M_N}(3).
\]
Then \(W_N>0\) for all sufficiently large \(N\), and
\begin{equation}\label{eq:two-base-asymptotic}
 \boxed{\displaystyle
 -\log W_N
 =2^N+N\log2-\log(N+1)+O(1).}
\end{equation}
In particular \(-\log W_N=2^N+O(N)\).  Thus changing from one
multiplicative orbit to two independent ones does not remove the
moving-point scale of the fixed-form obstruction.
\end{proposition}

\begin{proof}
Put \(x=2^N\) and \(y=3^{M_N}\).  Since
\(N\log2/\log3\) is not an integer, \(y>x\); both are integers, so
\(y\ge x+1\).  The standard bounds
\[
 \frac{e^{-u}}{u+1}<E_1(u)<\frac{e^{-u}}u
\]
show, uniformly here, that
\[
 \varepsilon_N(2)
 =\frac{N+1}{x}e^{-x}\{1+O(x^{-1})\}.
\]
Moreover \(M_N/N\to\log2/\log3<1\), and the same bounds give
\[
 \frac{\varepsilon_{M_N}(3)}{\varepsilon_N(2)}
 \le
 \left(\frac{\log2}{\log3}+o(1)\right)e^{-1}<1.
\]
Hence \(W_N\) is positive and differs from \(\varepsilon_N(2)\) by a
factor bounded above and below away from zero.  Taking logarithms yields
\eqref{eq:two-base-asymptotic}.
\end{proof}

Outside this cone, augmentation one alone has no sign content.

\begin{proposition}[Three-point flexibility]
\label{prop:three-point-flexibility}
Let \(0<a<b<c\) be positive integers.  Put
\[
 c^{(0)}=\left(\frac b{b-a},-\frac a{b-a},0\right),
 \qquad
 v=(c-b,a-c,b-a).
\]
For \(t\in\R\), the vector \(c(t)=c^{(0)}+tv\), supported at \(a,b,c\),
is augmentation one and log balanced.  Its tail is \(A+tD\), where
\[
 A>0,
 \qquad
 D=(b-a)(c-b)(c-a)f_q[a,b,c]>0.
\]
Hence rational admissible coefficients give a dense set of tail values in
\(\R\); in particular both signs and arbitrarily small nonzero values occur.
For example, at \(q=2\) and exponents \(1,2,3\),
\[
 -E_1(2)+5E_1(4)-3E_1(8)<0.
\]
\end{proposition}

\begin{proof}
The two normalizations are immediate.  The slope \(D\) is a positive second
divided difference by \eqref{eq:fq-convex}.  Since \(A+D\Q\) is dense in
\(\R\), all qualitative assertions follow.  For the explicit example use
\[
 E_1(x)>\frac{e^{-x}}{x+1},
 \qquad
 E_1(x)<\frac{e^{-x}}x:
\]
the displayed value is less than
\(-e^{-2}/3+5e^{-4}/4<0\).
\end{proof}

The next estimate identifies the coefficient height at which such
cancellation can begin.  Write a rational coefficient vector primitively as
\(c_j=a_j/D\), with \(D\ge1\), and put
\(H(c)=\max(D,|a_1|,\ldots,|a_r|)\).

\begin{proposition}[First-tail dominance]
\label{prop:first-tail-dominance}
Let \(n_1<\cdots<n_r\) be integers with \(a_1\ne0\).  If
\begin{equation}\label{eq:height-subcritical}
 (r-1)H(c)\frac{1+q^{-n_1}}q
 e^{-(q-1)q^{n_1}}\le\frac12,
\end{equation}
then the tail has the sign of \(a_1\) and
\begin{equation}\label{eq:first-tail-lower}
 \boxed{\displaystyle
 \left|\sum_{j=1}^rc_jE_1(q^{n_j})\right|
 \ge\frac{e^{-q^{n_1}}}{2H(c)(q^{n_1}+1)}.}
\end{equation}
Thus, for fixed \(q,r\), cancellation of the first active tail requires
\[
 \log H(c)\ge(q-1)q^{n_1}+O(1).
\]
\end{proposition}

\begin{proof}
The reverse triangle inequality gives
\[
 \left|\sum_jc_jE_1(q^{n_j})\right|
 \ge\frac{E_1(q^{n_1})-(r-1)H(c)E_1(q^{n_1+1})}{H(c)}.
\]
The elementary bounds used in the preceding proof imply
\[
 \frac{E_1(q^{n+1})}{E_1(q^n)}
 \le\frac{1+q^{-n}}q e^{-(q-1)q^n}.
\]
Now apply \eqref{eq:height-subcritical}.
\end{proof}

For consecutive triples the transition is exact.  Let
\[
 F_{j,N}=E_1(q^{N+j}),
 \quad
 A_N=(N+2)F_{1,N}-(N+1)F_{2,N},
 \quad
 B_N=F_{0,N}-2F_{1,N}+F_{2,N},
\]
and \(\rho_N=B_N/A_N\).  Both \(A_N,B_N\) are positive.

\begin{theorem}[Exact denominator phase transition]
\label{thm:denominator-transition}
Every rational augmentation-one, log-balanced coefficient vector on
\(N,N+1,N+2\) has the form
\[
 c(s)=(s,N+2-2s,s-N-1).
\]
For \(s=a/D\) with \(D\ge1\), its tail \(\varepsilon_N(s)\) satisfies
\begin{equation}\label{eq:denominator-exact}
 \boxed{\displaystyle
 |\varepsilon_N(a/D)|
 =\frac{B_N}{D}\left|a+\frac D{\rho_N}\right|,}
\end{equation}
and hence
\[
 \min_{a\in\Z}|\varepsilon_N(a/D)|
 =\frac{B_N}{D}\left\|\frac D{\rho_N}\right\|.
\]
In particular, \(D\le\rho_N/2\) implies
\(|\varepsilon_N(a/D)|\ge A_N\) for every integer \(a\).  Conversely,
whenever \(\rho_N\ge\tfrac12\) (hence for all sufficiently large \(N\)),
choosing a positive integer \(D\) nearest to \(\rho_N\) and taking
\(a=-1\) gives either an exact zero or
\begin{equation}\label{eq:denominator-upper}
 0<|\varepsilon_N(-1/D)|\le\frac{A_N}{2D}.
\end{equation}
Furthermore, with \(X=q^N\),
\begin{equation}\label{eq:rho-asymptotic}
 \rho_N\sim\frac q{N+2}e^{(q-1)X}.
\end{equation}
Thus the denominator threshold is exactly
\(e^{(q-1)q^N}/N\) up to a constant factor.
\end{theorem}

\begin{proof}
Substitution gives \(\varepsilon_N(s)=A_N+B_Ns\), proving
\eqref{eq:denominator-exact}.  If \(D/\rho_N\le1/2\), its nearest integer
is zero.  If \(|D-\rho_N|\le1/2\), then
\(|1-D/\rho_N|\le1/(2\rho_N)\), which gives
\eqref{eq:denominator-upper}.  Finally \eqref{eq:E1-asymptotic} yields
\[
 B_N\sim E_1(X),
 \qquad
 A_N\sim(N+2)E_1(qX),
\]
and hence \eqref{eq:rho-asymptotic}.
\end{proof}

When \(q\) is algebraic, exact zero can occur for at most one \(N\); by
\cref{thm:uN-defect}, any such zero would already force the algebraic
independence of \(\gamma,\delta\) over \(K\).  Discarding that possible one
index, the nearest-denominator construction gives normalized coefficient
height
\[
 H_N=q e^{(q-1)X}(1+o(1))
\]
and
\begin{equation}\label{eq:critical-height-upper}
 0<|\varepsilon_N|
 \le\frac{(N+2)^2}{2q^2X}
 e^{-(2q-1)X}(1+o(1)).
\end{equation}
Consequently a uniform lower bound, on normalized nonzero rational
augmentation-one triples, of the form
\[
 |\varepsilon|\ge C e^{-\beta X}H^{-\kappa}X^{-\lambda}
\]
requires
\begin{equation}\label{eq:height-frontier}
 \boxed{\beta+\kappa(q-1)\ge2q-1.}
\end{equation}
At equality it also requires \(\lambda\ge1\).  In particular, at the natural
point cost \(\beta=1\), one must have \(\kappa\ge2\).  This exponent refers
to the affine normalization \(\sum c_j=1\); clearing denominators multiplies
the value and changes the corresponding homogeneous height exponent.

The preceding sections construct exact finite spectral data converging very
rapidly to Euler-level traces.  To extract the arithmetic nature of the
limit, one would need a lower bound uniform in a moving algebraic sample
point.  We first show that sublinear point cost is impossible.

Let
\[
 F_1(z)=\Ein(z),\qquad
 F_2(z)=\Ein(2z),\qquad
 F_3(z)=\Ein(4z),
\]
and let \(P_0=X_1-2X_2+X_3\).

\begin{theorem}[A fixed-form critical-scale no-go]
\label{thm:sublinear-no-go}
For \(x\ge1\),
\begin{equation}\label{eq:R-def}
 R(x):=P_0(F_1(x),F_2(x),F_3(x))
 =E_1(x)-2E_1(2x)+E_1(4x)>0,
\end{equation}
and
\begin{equation}\label{eq:R-asymptotic}
 R(x)=\frac{e^{-x}}x
 \left(1+O(x^{-1})+O(e^{-x})\right).
\end{equation}
Consequently, for every \(A>0\) and every \(\vartheta<1\),
\begin{equation}\label{eq:R-sublinear}
 0<R(n)<e^{-An^\vartheta}
\end{equation}
for all sufficiently large integers \(n\).

More sharply, for every \(C_0>0\), a lower bound
\begin{equation}\label{eq:critical-polynomial-template}
 R(n)\ge C_0e^{-\sigma n}n^{-\lambda}
\end{equation}
cannot hold for all large \(n\) if either \(\sigma<1\), or if
\(\sigma=1\) and \(\lambda<1\).

In particular, no lower bound
\begin{equation}\label{eq:uniform-template}
 |P(F(\xi))|\ge
 \exp\!\left[-C(1+|\xi|)^\vartheta
 \Phi(\deg P,\log H(P))\right]
\end{equation}
can hold uniformly for all positive integral \(\xi\), all nonzero values and
all fixed \(E\)-function systems if \(\vartheta<1\), when \(C\) and \(\Phi\)
are independent of \(\xi\).
\end{theorem}

\begin{proof}
The connection formula \eqref{eq:Ein-connection} cancels \(\gamma\),
\(\log x\), and the scale logarithms because
\[
 1-2+1=0,
 \qquad
 \log x-2\log(2x)+\log(4x)=0.
\]
This proves the identity in \eqref{eq:R-def}.  Using
\(E_1(x)=\int_1^\infty e^{-xt}\,dt/t\), its integrand numerator is, with
\(u=e^{-xt}\),
\[
 u-2u^2+u^4=u(1-u)(1-u-u^2)>0,
\]
because \(0<u\le e^{-1}<(\sqrt5-1)/2\).  This proves positivity.
Asymptotic formula \eqref{eq:R-asymptotic} follows from
\eqref{eq:E1-asymptotic}; the terms at \(2x\) and \(4x\) are exponentially
smaller.  Hence
\[
 -\log R(n)=n+\log n+o(1).
\]
Since \(n^\vartheta=o(n)\), \eqref{eq:R-sublinear} follows.  In
\eqref{eq:uniform-template}, \(P_0\) has fixed degree one and fixed height
two, so its right-hand exponent is \(O(n^\vartheta)\), contradicting
\eqref{eq:R-sublinear}.
Finally, \eqref{eq:R-asymptotic} gives
\[
 \frac{R(n)}{C_0e^{-\sigma n}n^{-\lambda}}
 =C_0^{-1}e^{-(1-\sigma)n}n^{\lambda-1}(1+o(1)),
\]
which tends to zero in each of the two stated cases and contradicts
\eqref{eq:critical-polynomial-template}.
\end{proof}

\begin{remark}[Exact scope of the no-go]
The theorem gives the necessary condition \(\vartheta\ge1\) under the
normalization \eqref{eq:uniform-template}, and at the critical exponential
scale it forces at least the polynomial loss \(n^{-1}\).  It does not prove
the existence or sufficiency of such a bound.  If arbitrary point-dependent
constants are allowed, or if the function system itself is changed with the
sample point, the exponent \(\vartheta\) alone is not an invariant.
\end{remark}

\begin{remark}[Augmentation zero versus one]
The fixed form \(P_0=(1,-2,1)\) has \(\omega(P_0)=0\) and zero logarithmic
moment, so it lies in the gamma-free stratum of
\eqref{eq:augmentation-decomposition}.  It disproves bounds that treat all
coefficient directions uniformly; by itself it does not disprove a theorem
restricted to \(\omega=1\).
\Cref{thm:augmentation-extremality,thm:denominator-transition}
describe what survives, and what fails, after that restriction.
\end{remark}

The same fixed second difference is built into the dyadic tower.

\begin{proposition}[Increment identity]
\label{prop:increment-identity}
For every \(N\ge1\),
\begin{equation}\label{eq:YN-increment}
 \boxed{\displaystyle
 Y_N-Y_{N+1}=(N+1)R(2^N).}
\end{equation}
Consequently
\begin{equation}\label{eq:epsilon-tail-R}
 \varepsilon_N=\sum_{k\ge N}(k+1)R(2^k).
\end{equation}
\end{proposition}

\begin{proof}
Expanding \eqref{eq:YN} gives
\[
 Y_N-Y_{N+1}
 =(N+1)\{\Ein(2^N)-2\Ein(2^{N+1})+\Ein(2^{N+2})\},
\]
which is \eqref{eq:YN-increment}.  Telescoping and
\(Y_N\to\gamma\) give \eqref{eq:epsilon-tail-R}.
\end{proof}

Thus the obstruction to a sublinear measure is not merely analogous to the
tower: it is its exact level increment.

We now formulate the remaining arithmetic question with a coefficient that
is fixed across every level.  For \(s\in\C\), define
\begin{equation}\label{eq:LN-def}
 L_N(s)=2sY_N^2+Y_N-2.
\end{equation}

\begin{theorem}[Common-coefficient critical-slope criterion]
\label{thm:critical-slope}
For \(s\in\C\), the limit
\[
 \lambda(s)=\lim_{N\to\infty}
 \frac{-\log|L_N(s)|}{2^N}
\]
exists whenever the logarithm is defined for all sufficiently large \(N\),
and
\begin{equation}\label{eq:lambda-dichotomy}
 \lambda(s)=
 \begin{cases}
 1,&s=S_2,\\
 0,&s\ne S_2.
 \end{cases}
\end{equation}
For \(s=S_2\), more precisely,
\begin{align}
 L_N(S_2)
 &=\frac{4-\gamma}{\gamma}\varepsilon_N
 +2S_2\varepsilon_N^2>0,\label{eq:LN-exact}\\
 -\log L_N(S_2)
 &=2^N+N\log2-\log(N+1)+O(1).
 \label{eq:LN-log-asymptotic}
\end{align}
Hence
\begin{equation}\label{eq:S2-slope-equivalence}
 \boxed{\displaystyle
 S_2\in\Qbar
 \quad\Longleftrightarrow\quad
 \exists s\in\Qbar\text{ with }\lambda(s)=1.}
\end{equation}
If \(s\in\Qbar\), the family \(\{L_N(s):N\ge1\}\) is countably
algebraically independent over \(K\).
\end{theorem}

\begin{proof}
Since \(2S_2\gamma^2+\gamma-2=0\), substituting
\(Y_N=\gamma+\varepsilon_N\) gives \eqref{eq:LN-exact}; combine this with
\eqref{eq:epsilon-properties} to get \eqref{eq:LN-log-asymptotic}.  If
\(s\ne S_2\), then
\[
 L_N(s)\longrightarrow2\gamma^2(s-S_2)\ne0,
\]
so the normalized logarithm tends to zero.  This proves
\eqref{eq:lambda-dichotomy} and \eqref{eq:S2-slope-equivalence}.

Finally, for algebraic \(s\), one has \(s\in K\) and
\[
 2sY_N^2+Y_N-(L_N(s)+2)=0.
\]
Thus \(K(Y_N:N\in I)\) is algebraic over
\(K(L_N(s):N\in I)\) for every finite set \(I\); when \(s=0\), the relation
is linear.  Algebraic independence of the \(Y_N\) therefore transfers to the
\(L_N(s)\).
\end{proof}

There is an equivalent bounded-distortion form that uses only finite-stage
gamma-free data.

\begin{corollary}[Bounded-distortion criterion]
\label{cor:bounded-distortion}
One has
\begin{equation}\label{eq:bounded-distortion-limit}
 \frac{L_N(S_2)}{Y_N-Y_{N+1}}
 \longrightarrow\frac{4-\gamma}{\gamma},
\end{equation}
whereas for \(s\ne S_2\),
\[
 \left|\frac{L_N(s)}{Y_N-Y_{N+1}}\right|\longrightarrow\infty.
\]
Consequently
\begin{equation}\label{eq:bounded-equivalence}
 S_2\in\Qbar
 \quad\Longleftrightarrow\quad
 \exists s\in\Qbar:\quad
 \left\{\frac{L_N(s)}{Y_N-Y_{N+1}}\right\}_{N\gg1}
 \text{ is bounded}.
\end{equation}
\end{corollary}

\begin{proof}
By \cref{prop:increment-identity} and \eqref{eq:R-asymptotic},
\[
 Y_N-Y_{N+1}\sim\frac{N+1}{2^N}e^{-2^N}
 \sim\varepsilon_N.
\]
Now use \eqref{eq:LN-exact}; if \(s\ne S_2\), the numerator tends to a
nonzero constant while the denominator tends to zero.
\end{proof}

\begin{remark}[Why a generic critical measure is not enough]
The fixed-form obstruction has size \(R(x)\sim e^{-x}/x\).  Under the
algebraicity hypothesis \(s=S_2\), the target at \(x=2^N\) has size
\[
 L_N(S_2)\sim C_\gamma(\log x)\frac{e^{-x}}x.
\]
It is therefore larger, by a logarithmic factor, than an unconditional
fixed-form example already permits.  A lower bound constrained only by this
worst-case pointwise scale cannot distinguish the two.  A surviving route
must use additional structure, for example the repetition of the \emph{same}
algebraic coefficient \(s\) over all levels.
\end{remark}

Existing quantitative \(E\)-function theorems cited here are fixed-point
theorems: the constants in the algebraic-independence measure depend on the
chosen nonzero algebraic evaluation point; see
Adamczewski--Faverjon \cite{AdamczewskiFaverjon2025} and
Fischler--Rivoal \cite{FischlerRivoal2024}.  They do not directly provide
the moving uniformity required by \cref{thm:critical-slope}.  This statement
is only about the scope of those theorems, not an impossibility claim for all
future methods.

\section{An exact asymmetric Sondow integral}
\label{sec:sondow}

One of the discarded barrier heuristics attempted to assign a universal
entropy floor to an asymmetric two-variable Sondow integral.  The correct
replacement is an exact finite formula.  For integers \(A\ge0\) and
\(B\ge1\), put
\begin{equation}\label{eq:Sondow-AB}
 \mathcal I_{A,B}=
 \int_0^1\!\!\int_0^1
 \frac{x^A(1-x)^B y^A(1-y)^B}
 {(1-xy)(-\log(xy))}\,dx\,dy.
\end{equation}

\begin{theorem}[Exact asymmetric evaluation]
\label{thm:Sondow-asymmetric}
One has
\begin{align}
 \mathcal I_{A,B}
 ={}&\binom{2B}{B}\gamma\notag\\
 &+2\sum_{j=0}^B\binom Bj^2
 (H_j-H_{B-j})\log((A+j)!)\notag\\
 &-\sum_{j=0}^B\binom Bj^2H_{A+j}.
 \label{eq:Sondow-exact}
\end{align}
In particular,
\begin{equation}\label{eq:Sondow-example}
 \mathcal I_{1,2}=6\gamma+3\log6-\frac{53}{6}.
\end{equation}
\end{theorem}

\begin{proof}
Use Tonelli's theorem,
\[
 -\frac1{\log(xy)}=\int_0^\infty(xy)^u\,du,
 \qquad
 \frac1{1-xy}=\sum_{k\ge0}(xy)^k,
\]
and the beta integral.  With
\[
 R_B(t)=\frac{B!}{t(t+1)\cdots(t+B)},
\]
this gives
\[
 \mathcal I_{A,B}
 =\sum_{m=A+1}^{\infty}\int_m^\infty R_B(t)^2\,dt.
\]
The rational kernel has the partial fraction identity
\begin{equation}\label{eq:Sondow-partial-fraction}
 \left\{\frac{B!}{t(t+1)\cdots(t+B)}\right\}^{\!2}
 =\sum_{j=0}^B\binom Bj^2
 \left\{\frac1{(t+j)^2}
 +\frac{2(H_j-H_{B-j})}{t+j}\right\}.
\end{equation}
Put \(a_j=\binom Bj^2\) and \(h_j=H_j-H_{B-j}\).  Integrating and summing
only up to \(m=M\) gives the finite identity
\begin{align*}
 \mathcal I_{A,B}^{(M)}
 ={}&\sum_{j=0}^Ba_j(H_{M+j}-H_{A+j})\\
 &-2\sum_{j=0}^Ba_jh_j
   \log\frac{(M+j)!}{(A+j)!}.
\end{align*}
The potentially divergent pieces cancel by
\[
 \sum_ja_j=\binom{2B}{B},\qquad
 \sum_ja_jh_j=0,
\]
together with the derivative of Vandermonde's identity
\[
 2\sum_jja_jh_j=\binom{2B}{B}.
\]
Indeed \(H_{M+j}=\log M+\gamma+o(1)\), while Stirling's formula gives
\[
 \log(M+j)!
 =M(\log M-1)+(j+\tfrac12)\log M
  +\tfrac12\log(2\pi)+o(1).
\]
Letting \(M\to\infty\) in the finite identity gives
\[
 \binom{2B}{B}\gamma
 -\sum_ja_jH_{A+j}
 +2\sum_ja_jh_j\log((A+j)!),
\]
which is \eqref{eq:Sondow-exact}.  Direct substitution gives
\eqref{eq:Sondow-example}.
\end{proof}

Formula \eqref{eq:Sondow-exact} replaces, rather than supports, any fixed
numerical entropy floor for this asymmetric family.  The remaining
Diophantine bottleneck is the transverse height of the prime-exponent vector
in the logarithmic kernel.  We make no novelty claim beyond this exact
evaluation without a separate priority comparison with the broader Sondow
literature \cite{Sondow2003}.

\section{The Pr\'evost grid and residual-prime saturation}
\label{sec:prevost-grid}

For \(0\le m\le n\), Pr\'evost's family \cite{Prevost2008} can be written
\begin{align}
 J_{n,m}&=\gamma+L_{n,m}-2H_n,\label{eq:prevost-form}\\
 L_{n,m}&=\sum_{k=0}^n(-1)^{n+k}
 \binom nk\binom{n+k}k\log(n-m+k+1).
 \label{eq:prevost-log}
\end{align}
Let \(d_n=\lcm(1,\ldots,n)\) and \(j_{n,m}=d_nJ_{n,m}\).  If
\(m/n\to\alpha\in(0,1)\), its saddle rate is
\begin{equation}\label{eq:prevost-rate}
 |J_{n,m}|=\exp\{(\log r(\alpha)+o(1))n\},
 \qquad
 r(\alpha)=
 \alpha^\alpha(1-\alpha)^{2-2\alpha}(2-\alpha)^{\alpha-2},
\end{equation}
up to a polynomial factor.  The unique optimum is
\begin{equation}\label{eq:prevost-optimum}
 \alpha_*=1-\frac1{\sqrt2},
 \qquad
 r(\alpha_*)=3-2\sqrt2,
 \qquad
 -\log r(\alpha_*)=\log(3+2\sqrt2).
\end{equation}

We first isolate an exact linear-algebra obstruction that is independent of
the saddle estimate.  Fix a finite set \(I\) of cells and collect the
normalized forms into
\begin{equation}\label{eq:prevost-vector}
 \mathbf j=\gamma\mathbf B-\mathbf A
 +\sum_{\ell\in\mathcal P_I}(\log\ell)\mathbf V_\ell\in\R^I,
\end{equation}
where \(\mathcal P_I\) is the finite set of primes occurring in the
logarithmic coordinates.  The vectors \(\mathbf B,\mathbf A,\mathbf V_\ell\)
are integral.  We write \(B(x)=\langle\mathbf B,x\rangle\), and similarly
for the other coordinate functionals.

\begin{theorem}[Residual-prime Hilbert saturation]
\label{thm:residual-prime-saturation}
Let \(\mathcal S\subset\mathcal P_I\), let \(a/b\in\Q\) be reduced with
\(b>0\), and put
\[
 E_{\mathcal S,a/b}
 =\bigcap_{\ell\notin\mathcal S}\ker V_\ell
  \cap\ker(aB-bA).
\]
If \(P\) is orthogonal projection onto this subspace, then
\begin{equation}\label{eq:residual-projection}
 \boxed{\displaystyle
 bP\mathbf j=(b\gamma-a)P\mathbf B
 +b\sum_{p\in\mathcal S}(\log p)P\mathbf V_p.}
\end{equation}

Let \(c=(c_p)_{p\in\mathcal S}\in\Z^{\mathcal S}\) be primitive and
nonzero.  Choose \(h_p\in\Z\) with \(\sum h_pc_p=1\), set
\[
 \mathbf T=\sum_{p\in\mathcal S}h_p\mathbf V_p,
 \qquad
 \alpha_c=\sum_{p\in\mathcal S}c_p\log p,
\]
write \(T(x)=\langle\mathbf T,x\rangle\), and define
\[
 E_{c,a/b}=E_{\mathcal S,a/b}
 \cap\bigcap_{p\in\mathcal S}\ker(V_p-c_pT).
\]
Denote its orthogonal projection by \(P_c\).
If \(\gamma=a/b\), then
\begin{equation}\label{eq:residual-direction}
 \boxed{P_c\mathbf j=\alpha_cP_c\mathbf T.}
\end{equation}
If \(\Lambda=E_{c,a/b}\cap\Z^I\) and \(T(\Lambda)=g\Z\), \(g>0\), then
the real quotient minimum
\[
 \lambda_\R=\inf\{\|x\|:x\in E_{c,a/b},\ T(x)=g\}
\]
satisfies
\begin{equation}\label{eq:residual-saturation}
 \boxed{\displaystyle
 \lambda_\R=\frac g{\|P_c\mathbf T\|},
 \qquad
 \lambda_\R\|P_c\mathbf j\|=g|\alpha_c|.}
\end{equation}
\end{theorem}

\begin{proof}
Project \eqref{eq:prevost-vector}.  All primes outside \(\mathcal S\)
vanish and \(aP\mathbf B=bP\mathbf A\), proving
\eqref{eq:residual-projection}.  The fixed-direction equations give
\(P_c\mathbf V_p=c_pP_c\mathbf T\), and
\eqref{eq:residual-direction} follows under \(\gamma=a/b\).
Finally, the restriction of \(T\) to \(E_{c,a/b}\) is represented by
\(P_c\mathbf T\).  Cauchy--Schwarz is sharp in its direction and gives
the first identity in \eqref{eq:residual-saturation}; the second follows
from \eqref{eq:residual-direction}.
\end{proof}

\begin{corollary}[Phase-blind quotient barrier]
\label{cor:phase-blind-quotient}
Under the hypotheses of the second part of
\cref{thm:residual-prime-saturation}, and under \(\gamma=a/b\),
\[
 \lambda_\R\|\mathbf j\|\ge g|\alpha_c|.
\]
Thus a strict inequality in the opposite direction cannot follow from this
real Euclidean quotient identity and Cauchy--Schwarz alone.  Integral
divisibility or genuinely lattice-theoretic information is not excluded.
Any strict reverse inequality would already imply \(\gamma\ne a/b\).
\end{corollary}

\begin{remark}[The one-prime shear]
Let \(P\) now denote orthogonal projection onto
\(\bigcap_{\ell\ne2}\ker V_\ell\), with no candidate rationality
hyperplane imposed, and suppose \(x=P\mathbf V_2\ne0\).  Put
\[
 z=P(\gamma\mathbf B-\mathbf A),
 \qquad \varepsilon=\|P\mathbf j\|.
\]
Then \(P\mathbf j=z+(\log2)x\), and hence
\begin{equation}\label{eq:log2-shear}
 \frac{\det\operatorname{Gram}(x,z)}{\|x\|^2}
 =\|P_{x^\perp}P\mathbf j\|^2\le\varepsilon^2,
 \qquad
 \frac{\langle x,z\rangle}{\|x\|^2}
 =-\log2+O\!\left(\frac\varepsilon{\|x\|}\right).
\end{equation}
Thus \(-\log2\) is the longitudinal shear coefficient, while
\(P_{x^\perp}P\mathbf j\) is the narrow transverse residual.  For an
external proxy \(\widetilde\gamma=a/b\ne\gamma\),
\[
 P(\widetilde\gamma\mathbf B-\mathbf A)
 =z+(\widetilde\gamma-\gamma)P\mathbf B.
\]
More explicitly, if an integral coefficient vector \(c\) kills the
odd-prime labels and has residual label \(k=V_2(c)\), then
\[
 (aB(c)-bA(c))+bk\log2
 =b\langle\mathbf j,c\rangle+(a-b\gamma)B(c).
\]
Thus a proxy form is not controlled by the analytic norm of \(P\mathbf j\)
unless \(a/b=\gamma\); finite-window diagnostics based on an external proxy
cannot be used as irrationality evidence.
\end{remark}

\subsection{The fixed-place Ridout deficit}

For a fixed finite set \(S\) of rational primes and \(0\ne r\in\Z\), write
\[
 r_S=\prod_{p\in S}p^{v_p(r)},
 \qquad r_{S^c}=|r|/r_S.
\]

\begin{proposition}[Fixed-\(S\) Ridout budget]
\label{prop:ridout-budget}
Let \(\xi\) be a real algebraic irrational number.  For every
\(\epsilon>0\), all but finitely many reduced fractions \(a/b\) with
\(b>0\), \(a\ne0\), and \(|\xi-a/b|<1\) satisfy
\begin{equation}\label{eq:ridout-linear-form}
 |b\xi-a|\ge\frac{a_S}{b_{S^c}b^\epsilon}.
\end{equation}
Consequently, suppose along a sequence that
\[
 \log b_N=hN+o(N),
 \quad
 \log(a_N)_S+\log(b_N)_S=\sigma N+o(N),
 \quad
 |b_N\xi-a_N|=e^{-(\beta+o(1))N}.
\]
Assume here that \(h>0\), that \(b_N\to\infty\), and that the reduced
fractions \(a_N/b_N\) take infinitely many distinct values.
If \(\beta>h-\sigma\), then \(\xi\) is not algebraic irrational.  If
\(\beta\le h-\sigma\), these rates alone do not enter the range excluded by
Ridout's theorem.
\end{proposition}

\begin{proof}
Ridout's theorem \cite{Ridout1958} says that
\[
 \prod_{p\in S}|ab|_p
 \min\{1,|\xi-a/b|\}<b^{-2-\epsilon}
\]
has only finitely many reduced solutions.  For a good approximant the left
side is \(|\xi-a/b|/(a_Sb_S)\).  Multiplying the complementary inequality
by \(b\) gives \eqref{eq:ridout-linear-form}; in the homogeneous-height
form of the theorem, \(|\xi-a/b|<1\) makes
\(\max(|a|,b)\ll_\xi b\), and applying the theorem with
\(\epsilon/2\) absorbs the fixed comparison constant for large \(b\).
Taking logarithms and then letting \(\epsilon\downarrow0\) proves the rate
assertion.
\end{proof}

\begin{theorem}[The Pr\'evost saddle--lcm ledger]
\label{thm:prevost-ridout-deficit}
For any integer sequence \(m_n\) with \(m_n/n\to\alpha_*\),
\[
 |d_nJ_{n,m_n}|=
 \exp\{-(\log(3+2\sqrt2)-1+o(1))n\}.
\]
For every fixed finite set \(S\),
\[
 \log(d_n)_S=|S|\log n+O_S(1)=o(n),
 \qquad
 \log d_n=n+o(n).
\]
Thus, in the favorable height-one ledger in which the lcm is the only
denominator cost and a fixed \(S\) contributes no additional exponential
divisibility, the analytic exponent falls short of the Ridout threshold by
\begin{equation}\label{eq:ridout-deficit}
 \boxed{\displaystyle
 \Delta_R=2-\log(3+2\sqrt2)
 =0.2372528259\ldots.}
\end{equation}
This is an exact saddle--denominator deficit, not yet a rational
approximation to \(\gamma\), because logarithmic coordinates remain in
\(J_{n,m_n}\).
\end{theorem}

\begin{proof}
Combine \eqref{eq:prevost-rate} with the prime-number theorem form
\(\log d_n=\psi(n)=n+o(n)\).  Differentiating \(\log r\) gives
\[
 (\log r)'(\alpha)
 =\log\frac{\alpha(2-\alpha)}{(1-\alpha)^2},
\]
which proves the unique optimum \eqref{eq:prevost-optimum}.  Finally
\(v_p(d_n)=\lfloor\log n/\log p\rfloor\), so every fixed \(S\) contributes
only \(o(n)\).  Comparing the resulting exponent
\(\log(3+2\sqrt2)-1\) with the height-one threshold (1) gives
\eqref{eq:ridout-deficit}.
\end{proof}

\begin{corollary}[Conditional rational-extraction budget]
Suppose an elimination of the logarithmic coordinates produces infinitely
many distinct reduced fractions \(a_n/b_n\) such that
\[
 \log b_n=n+o(n),\qquad
 \log(a_n)_S+\log(b_n)_S=o(n),
\]
and preserves the saddle rate
\[
 |b_n\gamma-a_n|
 =\exp\{-(\log(3+2\sqrt2)-1+o(1))n\}.
\]
Then \cref{prop:ridout-budget} does not exclude algebraic irrationality of
\(\gamma\).  Within this height-one template, crossing the Ridout threshold
requires an additional fixed-\(S\) valuation density greater than
\(\Delta_R\), an archimedean gain of the same size, or a corresponding
reduction in height.
\end{corollary}

\begin{remark}[Exact scope]
This is a budget theorem, not an exclusion of every linear combination of
Pr\'evost forms.  Further cancellation could increase the archimedean
exponent, and a new automatic divisibility theorem could increase the
\(S\)-part.  Neither input is presently available.  Enlarging a fixed finite
set \(S\) alone cannot help because \(\log(d_n)_S=o(n)\).
\end{remark}

\begin{proposition}[Bounded-degree logarithmic transfer]
\label{prop:bounded-degree-transfer}
After cancelling all odd-prime coordinates, write a resulting form as
\[
 J=B\gamma-A+k\log2,
 \qquad A,B,k\in\Z,
\]
and assume \(k\ne0\).  If \(\gamma\) is algebraic of degree \(d\), then
\[
 \beta=\frac{A-B\gamma}{k}
\]
has degree at most \(d\), naïve height at most
\(C_\gamma\max(1,|A|,|B|,|k|)^d\), and
\[
 |\log2-\beta|=\frac{|J|}{|k|}.
\]
Thus any such attack is bounded by the degree-\(d\) algebraic approximation
exponent of \(\log2\), with an additional factor \(d\) in the height
exponent.
\end{proposition}

\begin{proof}
If \(f\) is the minimal polynomial of \(\gamma\) and \(B\ne0\), then
\(\beta\) is a zero of
\[
 B^df\!\left(\frac{A-kX}{B}\right),
\]
and the affine transformation shows that its degree is exactly \(d\).
The displayed polynomial has coefficient height
\(O_\gamma(\max(1,|A|,|B|,|k|)^d)\); the standard factor-height bound
gives the same estimate, up to a constant depending on \(\gamma\), for the
minimal polynomial of \(\beta\).  The case \(B=0\) is rational, and the
approximation identity is immediate.
\end{proof}

\begin{remark}
The available bounds are
\[
 \mu(\log2)\le3.574553902525\ldots,
 \qquad
 \mu_2(\log2)\le12.841618132152\ldots;
\]
see \cite{Marcovecchio2009,Marcovecchio2025}.  Thus a rational branch must
beat the first exponent, while the degree-two height transfer makes
\(2\mu_2(\log2)=25.683236264304\ldots\) the corresponding coefficient-height
target.  These are target exponents once the growth of
\(\max(1,|A|,|B|,|k|)\) is known; no comparison with a projected Pr\'evost
family is claimed here without a separate height lemma.
\end{remark}

\section{Fixed-prime Rodrigues twists and universal saddle cancellation}
\label{sec:fixed-prime-rodrigues}

The deficit \eqref{eq:ridout-deficit} suggests inserting a factor
\(p^{cn}\), for one fixed prime \(p\), directly into the Pr\'evost kernel.
The following deformation preserves integrality, degree, and a linear number
of moments, and produces this factor in the endpoint coefficient.  Even if
the whole endpoint factor is favorably credited as nonarchimedean gain, the
real saddle cancels it in the precise sense of
\cref{thm:universal-fixed-p-cancellation}.

Fix a prime \(p\).  For \(n\ge1\), \(1\le\ell\le n\), put
\begin{equation}\label{eq:fixed-p-rodrigues}
 \mathcal P_{n,\ell,p}(x)=
 \frac{(-1)^\ell}{\ell!}\frac{d^\ell}{dx^\ell}
 \left[x^\ell(1-x)^\ell
 \bigl(1+(p-1)x\bigr)^{n-\ell}\right].
\end{equation}
Write \(\mathcal P_{n,\ell,p}(x)=\sum_{k=0}^nc_kx^k\).

\begin{proposition}[Arithmetic and orthogonality]
\label{prop:fixed-p-arithmetic}
One has \(\mathcal P_{n,\ell,p}\in\Z[x]\),
\(\deg\mathcal P_{n,\ell,p}=n\),
\[
 \int_0^1x^j\mathcal P_{n,\ell,p}(x)\,dx=0
 \qquad(0\le j<\ell),
\]
and
\begin{equation}\label{eq:fixed-p-endpoint}
 \boxed{\mathcal P_{n,\ell,p}(1)=p^{n-\ell}.}
\end{equation}
\end{proposition}

\begin{proof}
Expanding before differentiation gives
\[
 c_k=(-1)^\ell\binom{\ell+k}{\ell}
 \sum_{a+b=k}(-1)^a
 \binom\ell a\binom{n-\ell}b(p-1)^b,
\]
which proves integrality and the degree assertion.  Integration by parts
\(\ell\) times proves the moment identities; all boundary terms vanish.
At \(x=1\), only the term in which all derivatives fall on
\((1-x)^\ell\) survives, giving \eqref{eq:fixed-p-endpoint}.
\end{proof}

Put
\[
 W(x)=\frac1{\log x}+\frac1{1-x}
\]
and, for \(0\le m\le\ell\), define
\begin{equation}\label{eq:fixed-p-moment}
 \mathcal J_{n,\ell,m,p}=
 \int_0^1x^{\ell-m}\mathcal P_{n,\ell,p}(x)W(x)\,dx.
\end{equation}

\begin{proposition}[Exact logarithmic linear form]
\label{prop:fixed-p-linear-form}
Let
\[
 Q_{n,\ell,p}=p^{n-\ell},
 \qquad
 A_{n,\ell,p}=\sum_{k=0}^nc_kH_k,
\]
and
\[
 L_{n,\ell,m,p}=
 \sum_{k=0}^nc_k\log(\ell-m+k+1).
\]
Then
\begin{equation}\label{eq:fixed-p-exact-form}
 \boxed{\mathcal J_{n,\ell,m,p}
 =Q_{n,\ell,p}\gamma-A_{n,\ell,p}+L_{n,\ell,m,p},}
\end{equation}
and \(d_nA_{n,\ell,p}\in\Z\).
\end{proposition}

\begin{proof}
Use Euler's integral \(\gamma=\int_0^1W(x)\,dx\) and Frullani's identity
\[
 \int_0^1\frac{x^r-1}{\log x}\,dx=\log(r+1).
\]
The rational part is
\begin{align*}
 &\int_0^1\frac{\mathcal P(1)-x^{\ell-m}\mathcal P(x)}{1-x}\,dx\\
 &\quad=\int_0^1\frac{\mathcal P(1)-\mathcal P(x)}{1-x}\,dx
 +\int_0^1\mathcal P(x)\frac{1-x^{\ell-m}}{1-x}\,dx.
\end{align*}
The second integral vanishes by the moment conditions, and the first is
\(\sum c_kH_k\).  This proves \eqref{eq:fixed-p-exact-form}; its denominator
divides \(d_n\).
\end{proof}

Pr\'evost's positive Markov--Stieltjes representation is
\begin{equation}\label{eq:prevost-markov}
 W(u)=\int_0^1\frac{\omega(t)}{1-(1-u)t}\,dt,
 \qquad
 \omega(t)=\frac1{t\{\log^2(t^{-1}-1)+\pi^2\}}>0.
\end{equation}

\begin{proposition}[Positive double integral]
\label{prop:fixed-p-positive}
One has
\begin{align}
 (-1)^m\mathcal J_{n,\ell,m,p}
 &=\int_0^1\!\int_0^1
 \frac{u^\ell(1-u)^\ell
 (1+(p-1)u)^{n-\ell}
 (1-t)^{\ell-m}t^m}
 {(1-(1-u)t)^{\ell+1}}
 \omega(t)\,dt\,du>0.
 \label{eq:fixed-p-double}
\end{align}
\end{proposition}

\begin{proof}
Insert \eqref{eq:prevost-markov}, use Fubini, and integrate by parts
\(\ell\) times.  The polynomial part is annihilated, while
\[
 \frac{d^\ell}{du^\ell}
 \frac{u^{\ell-m}}{1-(1-u)t}
 =(-1)^m\ell!\frac{(1-t)^{\ell-m}t^m}
 {(1-(1-u)t)^{\ell+1}}.
\]
This gives \eqref{eq:fixed-p-double} and positivity.
\end{proof}

Let \(\theta=\ell/n\), \(\alpha=m/\ell\), and set
\begin{align}
 \phi_\alpha(u,t)&=
 \frac{u(1-u)(1-t)^{1-\alpha}t^\alpha}
 {1-(1-u)t},\label{eq:prevost-saddle-kernel}\\
 \Phi_{\theta,\alpha,p}(u,t)&=
 \phi_\alpha(u,t)^\theta
 (1+(p-1)u)^{1-\theta},\label{eq:fixed-p-saddle}\\
 R_{\theta,\alpha,p}&=-\log\max_{0\le u,t\le1}
 \Phi_{\theta,\alpha,p}(u,t).
\end{align}
We extend \(\phi_\alpha\), and hence
\(\Phi_{\theta,\alpha,p}\), continuously by zero at the boundary corner
\((u,t)=(0,1)\).  Indeed the quotient in
\eqref{eq:prevost-saddle-kernel} tends to zero there for
\(0<\alpha<1\).

For the finite parameters \(\theta_n=\ell/n\),
\(\alpha_n=m/\ell\), the positive double integral is
\[
 (-1)^m\mathcal J_{n,\ell,m,p}
 =\int_0^1\!\int_0^1
 \Phi_{\theta_n,\alpha_n,p}(u,t)^n
 \frac{\omega(t)}{1-(1-u)t}\,dt,du.
\]
The remaining measure has finite mass
\[
 \int_0^1\!\int_0^1
 \frac{\omega(t)}{1-(1-u)t}\,dt,du
 =\int_0^1W(u)\,du=\gamma.
\]
This gives the upper Laplace bound, while a small rectangle around an
interior maximizer gives the matching lower bound.  Consequently,
whenever \(\ell/n\to\theta\in(0,1]\) and
\(m/\ell\to\alpha\in(0,1)\),
\begin{equation}\label{eq:fixed-p-laplace}
 \lim_{n\to\infty}|\mathcal J_{n,\ell,m,p}|^{1/n}
 =e^{-R_{\theta,\alpha,p}}.
\end{equation}

\begin{theorem}[Universal fixed-prime saddle cancellation]
\label{thm:universal-fixed-p-cancellation}
For every prime \(p\), \(0<\theta\le1\), and \(0<\alpha<1\),
\begin{equation}\label{eq:universal-fixed-p-bound}
 \boxed{\displaystyle
 R_{\theta,\alpha,p}+(1-\theta)\log p
 \le\log(3+2\sqrt2).}
\end{equation}
Thus no member of \eqref{eq:fixed-p-rodrigues} improves the original
Pr\'evost saddle after the complete endpoint factor
\(p^{n-\ell}\) is favorably credited as fixed-\(p\) gain.
\end{theorem}

\begin{proof}
Since
\[
 e^{-\{R_{\theta,\alpha,p}+(1-\theta)\log p\}}
 =\max_{u,t}
 \left[
  \phi_\alpha(u,t)^\theta
  \left(\frac{1+(p-1)u}{p}\right)^{1-\theta}
 \right],
\]
it is enough to bound this normalized maximum from below.  Now
For \(0\le u\le1\),
\[
 \frac{1+(p-1)u}{p}=u+\frac{1-u}{p}\ge u.
\]
At \(t=1/2\), the dependence on \(\alpha\) disappears:
\[
 \phi_\alpha\!\left(u,\frac12\right)
 =\frac{u(1-u)}{1+u}.
\]
Consequently the normalized maximum is at least
\[
 \max_{0\le u\le1}
 u\left(\frac{1-u}{1+u}\right)^\theta
 \ge\max_{0\le u\le1}\frac{u(1-u)}{1+u}
 =3-2\sqrt2.
\]
Taking negative logarithms proves \eqref{eq:universal-fixed-p-bound}.
\end{proof}

\begin{corollary}[The fixed-place deficit survives]
\label{cor:fixed-p-deficit}
Even under the favorable assumptions that logarithmic coordinates are
removed at no exponential cost, \(d_n=e^{n+o(n)}\) is the only denominator
cost, and the full \(p^{n-\ell}\) factor survives primitive reduction, the
integerized exponent satisfies
\[
 \bigl(R_{\theta,\alpha,p}-1\bigr)+(1-\theta)\log p
 \le\log(3+2\sqrt2)-1<1.
\]
Thus, even after the entire endpoint factor is credited, this favorable
ledger does not cross the fixed-\(S\) Ridout threshold; its remaining gap is
at least \(\Delta_R\) from \eqref{eq:ridout-deficit}.
\end{corollary}

\begin{proposition}[Balanced factorial ratios have zero fixed-prime slope]
\label{prop:balanced-factorial-fixed-p}
Let
\[
 \mathcal R(n,k)=
 \frac{\prod_{i=1}^rL_i(n,k)!}{\prod_{j=1}^sM_j(n,k)!},
\]
where the integral linear forms are nonnegative and \(O(n)\) on the index
range, and suppose \(\sum_iL_i=\sum_jM_j\).  Then for every fixed prime
\(p\),
\[
 \lvert v_p(\mathcal R(n,k))\rvert=O(\log n)
\]
uniformly wherever the ratio is defined.  In particular its fixed-prime
valuation density is zero.
\end{proposition}

\begin{proof}
Legendre's formula gives a sum over \(p^\nu\).  At each level the linear
parts cancel by balance, leaving a bounded floor error; only
\(O(\log n)\) levels are nonzero.
\end{proof}

\begin{remark}
This covers the standard balanced Ap\'ery, Franel, Domb, and Christol
factorial ratios.  The valuation exponents supplied directly by the
digit-count and \(p\)-Lucas mechanisms of \cite{Delaygue2018} involve at
most \(1+\lfloor\log_p n\rfloor\) base-\(p\) digits and therefore provide
only an \(O(\log n)\) guaranteed gain.  This does not exclude additional
problem-specific cancellation inside a multisum.  If a factorial-balance
defect is \(\delta n+O(1)\), \(\delta\ne0\), Stirling's formula produces a
\(\delta n\log n\) term in the archimedean ledger.  A successful kernel
must therefore use information not already encoded by the balanced
factorial ratio or the positive endpoint weight twist analyzed here.
\end{remark}

\section{The Stokes-framing barrier}
\label{sec:stokes-framing}

The preceding barriers distinguish gamma-free directions from the Euler
extension.  The analogous distinction for an elementary differential
operator annihilating the Euler \(E\)-function is exact.

Put
\[
 \mathscr E(z)=\sum_{n\ge1}\frac{z^n}{n\,n!},
 \qquad F(x)=\mathscr E(-x)=-\Ein(x),
\]
and let \(\partial=\partial_x\).  Consider
\begin{equation}\label{eq:euler-E-operator}
 \mathcal L=(\partial+1)\circ\partial\circ x\circ\partial,
 \qquad
 \mathcal Ly=xy'''+(x+2)y''+y'.
\end{equation}

\begin{theorem}[Euler Stokes-framing barrier]
\label{thm:stokes-framing}
On \(\Omega=\C\setminus(-\infty,0]\), with principal branches,
\[
 \mathcal B_\infty=(1,\log x,E_1(x)),
 \qquad
 \mathcal B_0=(1,F(x),\log x)
\]
are solution bases of \(\mathcal L\), and
\begin{equation}\label{eq:C-gamma}
 \begin{pmatrix}1\\F(x)\\\log x\end{pmatrix}
 =C_\gamma
 \begin{pmatrix}1\\\log x\\E_1(x)\end{pmatrix},
 \qquad
 C_\gamma=
 \begin{pmatrix}
 1&0&0\\
 -\gamma&-1&-1\\
 0&1&0
 \end{pmatrix}.
\end{equation}
The complete upper/lower lateral transition across the negative real axis is
\begin{equation}\label{eq:stokes-transition}
 \begin{pmatrix}1\\\log_+x\\E_{1,+}(x)\end{pmatrix}
 =S
 \begin{pmatrix}1\\\log_-x\\E_{1,-}(x)\end{pmatrix},
 \qquad
 S=
 \begin{pmatrix}
 1&0&0\\
 2\pi i&1&0\\
 -2\pi i&0&1
 \end{pmatrix}.
\end{equation}
In particular \(S\) is independent of \(\gamma\), and the coefficient row
\(c_\gamma=(-\gamma,-1,-1)\) satisfies
\begin{equation}\label{eq:stokes-cocycle-blind}
 \boxed{c_\gamma S=c_\gamma}
\end{equation}
for every complex value of \(\gamma\).

For every \(t\in\Qbar\), \(F+t\) is an \(E\)-function solution of the same
operator.  With \(\mathcal B_0^{(t)}=(1,F+t,\log x)\), its connection
matrix is
\begin{equation}\label{eq:framing-shear}
 C_t=U_tC_\gamma,
 \qquad
 U_t=\begin{pmatrix}1&0&0\\t&1&0\\0&0&1\end{pmatrix}.
\end{equation}
Consequently, under \(\gamma\in\Qbar\), the choice \(t=\gamma\) removes
the Euler entry by an algebraic change of Frobenius framing without changing
the operator or any unframed Stokes conjugacy class.
\end{theorem}

\begin{proof}
Since \(xF'(x)=e^{-x}-1\),
\[
 (\partial+1)\partial(xF')=(\partial+1)(-e^{-x})=0.
\]
Moreover,
\[
 W(1,\log x,E_1(x))=\frac{e^{-x}}{x^2}\ne0,
\]
and \(\det C_\gamma=1\), so both displayed triples are indeed bases.
The connection identity
\[
 F(x)=-\gamma-\log x-E_1(x)
\]
proves \eqref{eq:C-gamma}.  For \(r>0\),
\[
 \log_+(-r)-\log_-(-r)=2\pi i,
 \qquad
 E_{1,+}(-r)-E_{1,-}(-r)=-2\pi i,
\]
which gives \eqref{eq:stokes-transition}.  Matrix multiplication proves
\eqref{eq:stokes-cocycle-blind}.  Adding the constant solution \(t\) adds
\(t\) times the first row of the connection matrix to its second row,
proving \eqref{eq:framing-shear}.
\end{proof}

\begin{corollary}[Exterior blindness]
\label{cor:stokes-exterior-blind}
One has
\[
 \det C_\gamma=1,
 \qquad \det S=1,
 \qquad \operatorname{tr}S=3.
\]
The displayed lateral transition and every invariant of its unframed
conjugacy class are independent of \(\gamma\).  Their determinant, trace,
and exterior-power characteristic data lose the off-diagonal Euler framing.
\end{corollary}

\begin{remark}[Strict factors and framed Stokes constants]
Fischler--Rivoal use ``Stokes constants'' for sector-dependent coefficients
of a distinguished \(E\)-function in a formal basis at infinity; in that
broader sense \(-\gamma\) is such a coefficient.  After the logarithmic
formal monodromy is separated, the strict factor is
\(E_1\mapsto E_1-2\pi i\) and is gamma-free.  Their arithmetic theorem is a
containment theorem for possible coefficients, not a numerical-to-motivic
separation theorem; see \cite{FischlerRivoal2016}.
\end{remark}

\begin{remark}[Class-group averaging does not erase \(\gamma\)]
\label{rem:Chowla-Selberg-correction}
Suppose a Kronecker--Chowla--Selberg family is written in class-group
components as \(A_C=\kappa\gamma+B_C\), with the same Euler coefficient in
each class.  Character projection gives
\[
 \sum_C\overline{\chi(C)}A_C
 =\kappa\gamma\sum_C\overline{\chi(C)}
 +\sum_C\overline{\chi(C)}B_C.
\]
Thus the trivial character, equivalently the class-group average, preserves
the Euler term.  It disappears only in nontrivial character components or
other coefficient-sum-zero directions.  In representation-theoretic
language, \(\gamma\) occupies the trivial isotypic component.  This corrects
the opposite averaging heuristic; it is not a new Chowla--Selberg theorem.
See \cite{GrossCS}.
\end{remark}

\section{The mixed-Gevrey aperture}
\label{sec:mixed-gevrey}

Fischler--Rivoal call a formal Laurent series mixed arithmetic-Gevrey when
its nonnegative part is an \(E\)-function and its negative part is an
arithmetic Gevrey series of order one
\cite{FischlerRivoalMixed2024}.  The Euler connection formula produces a
particularly transparent family in this class.

For \(\lambda\in\Qbar_{>0}\), set
\begin{equation}\label{eq:mixed-dilation-family}
 \widehat D_\lambda(z)=
 \sum_{n\ge0}(-1)^n n!\lambda^{-n-1}z^{-n-1},\qquad
 U_\lambda(z)=e^{\lambda z},\qquad
 W_\lambda(z)=e^{\lambda z}\Ein(\lambda z)-\widehat D_\lambda(z).
\end{equation}

\begin{proposition}[Dilation system and functional rank]
\label{prop:mixed-dilation-rank}
Direction zero is not anti-Stokes for \(\widehat D_\lambda\), and its
Borel sum is
\begin{equation}\label{eq:Dhat-Borel-sum}
 (\widehat D_\lambda)_0(z)=e^{\lambda z}E_1(\lambda z).
\end{equation}
Consequently
\begin{equation}\label{eq:Wlambda-sum}
 (W_\lambda)_0(z)=e^{\lambda z}
 (\gamma+\log\lambda+\log z),
\end{equation}
on the principal positive sector, and
\begin{equation}\label{eq:mixed-dilation-system}
 \binom{U_\lambda}{W_\lambda}'
 =\begin{pmatrix}\lambda&0\\z^{-1}&\lambda\end{pmatrix}
 \binom{U_\lambda}{W_\lambda}.
\end{equation}
For distinct \(\lambda_1,\ldots,\lambda_m\in\Qbar_{>0}\), the \(2m\)
mixed functions \(U_{\lambda_i},W_{\lambda_i}\) are linearly independent
over \(\C(z)\).
\end{proposition}

\begin{proof}
The Laplace integral of the formal negative part is
\[
 \int_0^\infty\frac{e^{-\lambda zt}}{1+t}\,dt
 =e^{\lambda z}E_1(\lambda z),
\]
which proves \eqref{eq:Dhat-Borel-sum}; its Borel singularity lies on the
negative direction.  Subtracting this identity from
\(e^{\lambda z}\Ein(\lambda z)\) and using
\eqref{eq:Ein-connection} proves \eqref{eq:Wlambda-sum}.  Direct
differentiation proves \eqref{eq:mixed-dilation-system}.

Suppose a rational-function linear relation among all the displayed
functions exists, and sum it in direction zero.  Analytic continuation once
around zero adds \(2\pi i e^{\lambda_i z}\) to
\((W_{\lambda_i})_0\) and fixes \(U_{\lambda_i}\).  Subtraction gives a
rational-function relation among exponentials with distinct exponents, so
all coefficients of the \(W\)-terms vanish.  The remaining exponential
relation then forces all \(U\)-coefficients to vanish.
\end{proof}

\begin{proposition}[Unconditional non-divisibility]
\label{prop:mixed-nondivisibility}
For every \(\lambda\in\Qbar_{>0}\) and \(a\in\Qbar\), the formal quotient
\begin{equation}\label{eq:mixed-quotient}
 \frac{W_\lambda-aU_\lambda}{z-1}
\end{equation}
is not a mixed arithmetic-Gevrey function.
\end{proposition}

\begin{proof}
Write the normalized positive part of the numerator as
\[
 F(z)=e^{\lambda z}\{\Ein(\lambda z)-a\}
\]
and its negative part as \(-\widehat D_\lambda(z)\), which has zero constant
term.  If the quotient were mixed, comparison of nonnegative and
nonpositive Laurent parts after multiplication by \(z-1\) would force
\(F(1)\in\Qbar\).  But
\[
 F(1)=e^\lambda\{\Ein(\lambda)-a\},
\]
which is transcendental because \(\Ein(\lambda)\) is algebraically
independent over \(K\supset\Qbar(e^\lambda)\) by
\cref{thm:input-independence}.  This is a contradiction.
\end{proof}

We can now state the exact conditional payoff.  Write
\(\mathbf E,\mathbf D\) for the rings of algebraic-point values of
\(E\)-functions and of directionally summed arithmetic Gevrey-one series,
respectively, in the notation of Fischler--Rivoal.

\begin{theorem}[Conditional simultaneous dilation theorem]
\label{thm:conditional-mixed-dilations}
Assume the intersection conjecture
\begin{equation}\label{eq:ED-intersection}
 \mathbf E\cap\mathbf D=\Qbar.
\end{equation}
For any distinct
\(\lambda_1,\ldots,\lambda_m\in\Qbar_{>0}\), the \(2m\) numbers
\begin{equation}\label{eq:mixed-dilation-values}
 \boxed{\displaystyle
 e^{\lambda_i},\qquad
 e^{\lambda_i}(\gamma+\log\lambda_i)
 \quad(1\le i\le m)}
\end{equation}
are linearly independent over \(\Qbar\).  In particular every
\(\gamma+\log\lambda_i\) is transcendental; the case \(\lambda=1\) gives
the conditional transcendence of \(\gamma\).
\end{theorem}

\begin{proof}
Assume that the values in \eqref{eq:mixed-dilation-values} satisfy a
linear relation over \(\Qbar\).  Using
\(W_\lambda=e^{\lambda z}\Ein(\lambda z)-\widehat D_\lambda\), move all
summed Gevrey terms to the right.  The left side is in \(\mathbf E\), the
right side is in \(\mathbf D\), so their common value is algebraic by
\eqref{eq:ED-intersection}.  The left side has the form
\[
 \sum_i a_i e^{\lambda_i}
 +\sum_i b_i e^{\lambda_i}\Ein(\lambda_i).
\]
The \(\Ein(\lambda_i)\) are algebraically independent over the field \(K\)
containing all \(e^{\lambda_i}\), so algebraicity forces every \(b_i=0\).
The Lindemann--Weierstrass theorem applied to
\(0,\lambda_1,\ldots,\lambda_m\) then forces every \(a_i=0\).  This proves
linear independence.  The final transcendence assertion follows because
an algebraic ratio would make the two values in its block linearly
dependent.
\end{proof}

The dilation construction has an all-order resonant extension.  Write
\(\stirling{n}{r}\) for the unsigned Stirling number of the first kind, and for
\(r\ge1\) put
\begin{equation}\label{eq:all-order-Dhat}
 \widehat D_{\lambda,r}(z)=
 \sum_{n\ge0}(-1)^n\stirling{n+1}{r}
 \lambda^{-n-1}z^{-n-1},
\end{equation}
and
\begin{equation}\label{eq:all-order-mixed-family}
 \mathcal W_{\lambda,0}(z)=e^{\lambda z},\qquad
 \mathcal W_{\lambda,r}(z)
 =e^{\lambda z}\Ein_r(\lambda z)-\widehat D_{\lambda,r}(z)
 \quad(r\ge1).
\end{equation}

\begin{proposition}[All-order mixed jet]
\label{prop:all-order-mixed-jet}
Every \(\mathcal W_{\lambda,r}\) is mixed arithmetic-Gevrey, direction zero
is non-anti-Stokes, and
\begin{equation}\label{eq:all-order-mixed-generating}
 \boxed{\displaystyle
 \sum_{r\ge0}(\mathcal W_{\lambda,r})_0(z)t^r
 =e^{\lambda z}\Gamma(1-t)(\lambda z)^t.}
\end{equation}
They satisfy the triangular differential system
\begin{equation}\label{eq:all-order-mixed-system}
 \mathcal W_{\lambda,0}'=\lambda\mathcal W_{\lambda,0},
 \qquad
 \mathcal W_{\lambda,r}'
 =\lambda\mathcal W_{\lambda,r}+z^{-1}\mathcal W_{\lambda,r-1}
 \quad(r\ge1).
\end{equation}
For distinct positive algebraic \(\lambda_i\), every finite family of the
functions \(\mathcal W_{\lambda_i,r}\) is linearly independent over
\(\C(z)\).
\end{proposition}

\begin{proof}
The normalized Borel transform of \eqref{eq:all-order-Dhat} is
\[
 \frac{(-1)^{r-1}}{\lambda(r-1)!}
 \frac{\log^{r-1}(1+\xi/\lambda)}{1+\xi/\lambda},
\]
a \(G\)-function whose singularities lie on the negative ray.  Hence the
formal series is arithmetic Gevrey of order one and direction zero is
admissible.  The elementary split
\[
 \Gamma(1-t)x^t
 =1+t\sum_{r\ge1}\Ein_r(x)t^{r-1}
   -t x^t\Gamma(-t,x)
\]
together with the large-
\(x\) expansion
\[
 e^x x^t\Gamma(-t,x)
 \sim\sum_{n\ge0}(-1)^n(t+1)_n x^{-n-1}
\]
proves \eqref{eq:all-order-mixed-generating}, since
\([t^{r-1}](t+1)_n=\stirling{n+1}{r}\).  Differentiating the generating
identity and comparing coefficients proves
\eqref{eq:all-order-mixed-system}.

The sectorial sum on the right of
\eqref{eq:all-order-mixed-generating} is \(e^{\lambda z}\) times a
polynomial in \(\log z\) of degree \(r\), with nonzero leading coefficient
\(1/r!\).  Repeated monodromy differences about zero strip off the highest
logarithmic degree.  At each step the linear independence of exponentials
with distinct exponents forces the corresponding rational coefficients to
vanish.  Descending induction on \(r\) proves the last assertion.
\end{proof}

\begin{theorem}[Conditional all-order Gamma-jet independence]
\label{thm:conditional-all-order-gamma}
Assume \eqref{eq:ED-intersection}.  Let
\(\lambda_1,\ldots,\lambda_m\in\Qbar_{>0}\) be distinct and choose
integers \(R_i\ge0\).  Then all numbers
\begin{equation}\label{eq:conditional-all-order-values}
 \boxed{\displaystyle
 e^{\lambda_i}[t^r]\{\Gamma(1-t)\lambda_i^t\},
 \qquad 1\le i\le m,\quad0\le r\le R_i,}
\end{equation}
are linearly independent over \(\Qbar\).  In particular
\begin{equation}\label{eq:conditional-gamma-derivative-LI}
 \boxed{1,\Gamma'(1),\Gamma''(1),\ldots}
\end{equation}
are linearly independent over \(\Qbar\).
\end{theorem}

\begin{proof}
In a putative relation, split every \(r\ge1\) value as the value of the
\(E\)-function \(e^{\lambda z}\Ein_r(\lambda z)\) minus the summed
Gevrey value \((\widehat D_{\lambda,r})_0\), and move all Gevrey terms to
one side.  By \eqref{eq:ED-intersection}, the common value is algebraic.
The other side is a linear form in the jointly algebraically independent
coordinates \(\Ein_r(\lambda_i)\) over \(K\), by
\eqref{eq:resonant-input}; therefore every coefficient with \(r\ge1\)
vanishes.  Lindemann--Weierstrass then removes the remaining \(r=0\)
exponential terms.  Formula \eqref{eq:all-order-mixed-generating} at
\(z=1\) proves \eqref{eq:conditional-all-order-values}.  Taking
\(\lambda=1\) and using
\([t^r]\Gamma(1-t)=(-1)^r\Gamma^{(r)}(1)/r!\) proves
\eqref{eq:conditional-gamma-derivative-LI}.
\end{proof}

\begin{remark}[Priority and scope]
Fischler--Rivoal already prove from \eqref{eq:ED-intersection} the
transcendence of Gamma derivatives at the relevant rational points; in
particular the one-point transcendence of \(\gamma\) is not claimed as new
here.  Their stronger mixed zero-removal Conjecture~3 implies
\eqref{eq:ED-intersection}; alternatively, their conditional linear
specialization theorem and \cref{prop:mixed-dilation-rank} give a second
proof of the displayed rank.  The contribution of
\cref{thm:conditional-mixed-dilations,thm:conditional-all-order-gamma} is
the explicit simultaneous dilation family, its resonant all-order lift, the
use of the companion multipoint theorem, and the full linear rank.  It
does not assert algebraic independence: products of mixed functions need
not be mixed, so the cited specialization theorem is deliberately linear.
The value \((W_1)_0(1)=e\gamma\) is the same Euler extension coordinate
that reappears in the framed motive below, but no implication between the
two conjectural routes is known.
\end{remark}

\section{The Euler--Tate matched-line target}
\label{sec:euler-tate}

The exponential-motive construction supplies the canonical framing absent
from the scalar Stokes conjugacy class.

\begin{theorem}[Euler--Mascheroni motive; Fres\'an--Jossen]
\label{thm:FJ-euler-motive}
There is a rank-two exponential motive \(M(\gamma)\) with adapted de Rham
and Betti bases \((d_0,d_1)\) and \((b_0,b_1)\) for which, with columns
indexed by the de Rham basis and rows by the Betti basis, the comparison
matrix is
\begin{equation}\label{eq:euler-period-matrix}
 P_\gamma=\begin{pmatrix}1&\gamma\\0&2\pi i\end{pmatrix}.
\end{equation}
Equivalently,
\[
 \operatorname{comp}(d_0)=b_0,
 \qquad
 \operatorname{comp}(d_1)=\gamma b_0+2\pi i\,b_1.
\]
It lies in a non-split exact sequence
\begin{equation}\label{eq:euler-motive-extension}
 0\longrightarrow\Q(0)\longrightarrow M(\gamma)
 \longrightarrow\Q(-1)\longrightarrow0,
\end{equation}
satisfies \(\operatorname{Hom}(M(\gamma),\Q(0))=0\), and has motivic
Galois group
\[
 G_{M(\gamma)}=
 \left\{\begin{pmatrix}1&b\\0&a\end{pmatrix}:
 a\in\mathbb G_m,\ b\in\mathbb G_a\right\}.
\]
\end{theorem}

\begin{proof}
This is the Euler--Mascheroni construction and Galois-group computation of
Fres\'an--Jossen \cite[\S12.8]{FresanJossen2024}.
\end{proof}

Its determinant is \(2\pi i\), the same as that of the semisimplification,
and is again blind to the extension coordinate.  The missing input can be
stated for this one object.

\begin{proposition}[Complete matched-line classification]
\label{prop:matched-line-classification}
Every de Rham line mapping isomorphically to the quotient is uniquely of
the form
\[
 W_{\mathrm{dR}}(x)=\Qbar(xd_0+d_1),\qquad x\in\Qbar,
\]
and every Betti line with the same property is uniquely of the form
\[
 W_{\mathrm B}(u)=\Qbar(ub_0+b_1),\qquad u\in\Qbar.
\]
The two lines are matched if and only if
\begin{equation}\label{eq:matched-line-equation}
 \boxed{x+\gamma=2\pi i\,u.}
\end{equation}
Consequently a matched pair exists if and only if
\begin{equation}\label{eq:matched-line-locus}
 \boxed{\gamma\in\Qbar+2\pi i\,\Qbar
              =\Qbar+\pi\Qbar.}
\end{equation}
\end{proposition}

\begin{proof}
The displayed parameterizations are the affine charts of lines transverse
to the respective subobjects.  Comparison sends
\[
 xd_0+d_1\longmapsto(x+\gamma)b_0+2\pi i\,b_1.
\]
This spans \(ub_0+b_1\) exactly when
\eqref{eq:matched-line-equation} holds.  Eliminating \(x,u\in\Qbar\)
gives \eqref{eq:matched-line-locus}; the two fields coincide because
\(2i\in\Qbar^\times\).
\end{proof}

\begin{problem}[Euler--Tate matched-line lifting]
\label{prob:euler-tate-fullness}
After extension to \(\Qbar\), suppose one-dimensional de Rham and Betti
subspaces both map isomorphically onto the corresponding realization of
\(\Q(-1)\), and comparison maps one line to the other.  Prove that the
matched pair is induced by a subobject
\[
 \Q(-1)_{\Qbar}\hookrightarrow M(\gamma)_{\Qbar}.
\]
\end{problem}

\begin{theorem}[Strengthened payoff of matched-line lifting]
\label{thm:matched-line-implies-gamma}
A positive solution of \cref{prob:euler-tate-fullness} implies that
\begin{equation}\label{eq:matched-line-payoff}
 \boxed{\gamma\notin\Qbar+\pi\Qbar.}
\end{equation}
Equivalently, \(1,\pi,\gamma\) are linearly independent over \(\Qbar\).
In particular \(\gamma\), \(\gamma/\pi\), and \(S_2\) are transcendental.
\end{theorem}

\begin{proof}
If \(\gamma\in\Qbar+\pi\Qbar\),
\cref{prop:matched-line-classification} supplies a matched pair.  The
proposed lifting would give a copy of \(\Q(-1)\) mapping nontrivially to the quotient
in \eqref{eq:euler-motive-extension}.  The composite is a nonzero scalar,
so rescaling gives a section and hence a splitting.  Moreover
\(\operatorname{Hom}(-,-)\) commutes here with extension from \(\Q\) to
\(\Qbar\), so the Hom-vanishing in \cref{thm:FJ-euler-motive} remains
zero.  This contradiction proves \eqref{eq:matched-line-payoff}.  Membership
in \(\Qbar+\pi\Qbar\) is exactly a nontrivial \(\Qbar\)-linear relation
among \(1,\pi,\gamma\).  Algebraicity of either \(\gamma\) or
\(\gamma/\pi\) would give such a relation.  Finally, if \(S_2\) were
algebraic, \(2S_2\gamma^2+\gamma-2=0\) would make \(\gamma\) algebraic.
\end{proof}

\begin{remark}[Rational and algebraic versions]
A matched-line lifting statement restricted to \(\Q\)-defined lines would
already exclude the corresponding rational matched locus.  The
\(\Qbar\)-coefficient formulation in \cref{prob:euler-tate-fullness} is the
version that yields the full linear-independence conclusion
\eqref{eq:matched-line-payoff}.
\end{remark}

\begin{remark}[Exact scope of current rapid-decay structures]
Snodgrass constructs cyclotomic rational structures on rapid-decay
cohomology under a hypothesis chosen so that Stokes constants do not enter
\cite[\S1.2]{Snodgrass2026}.  This supplies a realization and period pairing,
but not the numerical-to-motivic fullness asserted in
\cref{prob:euler-tate-fullness}.  The problem is a single-object rank-one
target, far more narrowly scoped than the full exponential period conjecture
but not presently proved.
\end{remark}

\begin{remark}[Three independent apertures]
The common-coefficient problem \cref{prob:common-coefficient} is quantitative
and uses one algebraic coefficient across infinitely many moving levels.
The matched-line problem is categorical and uses one canonical framed
extension.  The mixed-Gevrey route uses the intersection
\(\mathbf E\cap\mathbf D=\Qbar\).  Each would prove the transcendence of
\(\gamma\); matched-line lifting would additionally prove the
transcendence of \(\gamma/\pi\), while the mixed route gives an all-order
linear-independence family.  No implication among the three inputs is known.
\end{remark}

\section{Conclusions and remaining problems}
\label{sec:conclusion}

The completed results fall into the following thematic groups.
\begin{enumerate}[label=\textup{(\arabic*)}]
 \item Every divisor \(\Ein-c\) has a universal polynomial inverse-trace
 calculus.  Its full trace ring is exactly \(F[c^{-1}]\), generated by the
 two even traces \(T_2,T_4\), while every transversal moving selection
 preserves algebraic independence.  The only nontrivial collision of
 \((T_2,T_3)\) is the pair \(\{6-2\sqrt3,6+2\sqrt3\}\); at every order
 there are also algebraic levels with zero inverse trace and wholly
 transcendental spectral divisors.  Finite-rank multiplicative groups have
 defect at most their rank and admit cofinite jointly independent
 subfamilies; the complete gamma-free, log-free cancellation space is
 identified.  The tails \(E_1(\alpha)\), and jointly the four classical
 families \(E_1,\Ei,\Ci,\operatorname{Si}\), have defect at most one more
 than the multiplicative rank.  Regularized lower incomplete-Gamma jets
 have defect at most the rank, uniformly across all jet orders, while fixed
 upper-jet rows have one further possible defect.  Sine-integral graph
 ideals, Dickman log-Laplace relations, and characteristic-zero algebraic
 matroids give further realizations of this engine.  At every finite set of
 levels, the complete trace field is exactly the rational quotient by the
 one-dimensional gauge shift.
 \item The dyadic and fixed-power balanced levels approach \(\gamma\) through
 explicit algebraically independent finite-stage data.  Their order-two
 and order-four traces give positive algebraically independent spectral
 decompositions of \(S_2\) and \(S_4\); the two positive carrier traces
 recover \(\gamma\).  The order-two carrier also gives an order-zero
 Laguerre--P\'olya determinant.  Neither the dyadic level sequence nor its
 second trace is P-recursive.
 \item The gamma-free map \(\mathscr R\) has transcendental inverse images at
 every algebraic level in its range, giving the explicit family \(r_m\).
 Normalized finite \(\Ein\)-certificates are unique.  This turns all
 cross-index real Euler witnesses into one countable exclusion theorem and
 yields the relative five-class theorem.  The critical cancellation
 coefficients \(u_N(q)\) have algebraic defect at most one.
 \item For the 2026 Challenge recurrence, the terminating \({}_2F_2\)
 state is explicitly gauge-equivalent to the contragredient normalized
 official Meijer--\(G\) state.  The finite bilinear identity fixes the
 conserved pairing; a flat dual-number \(SL_3\), equivalently \(SL_6\),
 includes the numerator.  The Beta--log transform connects the Challenge to
 the cubic family and every fixed offset.  A single rational carrier yields
 the denominator, numerator, full zeta log-jets, and a monic contact
 polynomial.  The exact maxima mixture, Meijer stop-loss interpretation,
 Poisson--binomial law, zero-density limit, regularized Gamma determinant,
 strengthened divisor, factorial-six Casoratian, all-initial-value limit
 classification, and integer solution lattice complete the finite
 structure.
 \item The Euler--Pad\'e convergents have a finite residue criterion on the
 \(p\)-adic unit boundary and converge as a full sequence at
 \(t=-1\) for \(p=2,3,5,13\).  One lcm-indexed rational subsequence converges
 placewise at the real place and at every fixed prime.  The denominator
 interpolation \(q_p\) identifies unbounded valuations with a \(p\)-adic
 zero, without confusing that event with failure of convergence.
 \item The dyadic augmentation tail is the unique minimum in the
 cumulative-positive cone, while outside that cone three points already
 produce dense tail values.  For fixed \(q\), the exact denominator
 transition for consecutive augmentation-one triples occurs at
 \(e^{(q-1)q^N}/N\).  A moving-point lower bound covering the displayed
 fixed system cannot have sublinear point cost; at linear cost the example
 forces the stated height loss.  The two-base form
 \(Y_N(2)-Y_{M_N}(3)\) retains the same leading \(2^N\) logarithmic scale.
 The exact asymmetric Sondow formula
 replaces the invalid universal entropy-floor heuristic.
 \item In each nondegenerate fixed residual direction of the Pr\'evost grid,
 the real Hilbert quotient is exactly saturated.  For the uncombined
 optimally saddled form, and absent additional exponential divisibility or
 cancellation, the baseline fixed-place ledger leaves the deficit
 \(2-\log(3+2\sqrt2)\).  The bounded-degree reduction transfers a surviving
 one-prime problem to quantitative approximation of \(\log2\), but does not
 close it.
 The fixed-prime Rodrigues deformation produces integral polynomials,
 exact logarithmic linear forms, positivity, and the endpoint factor
 \(p^{(1-\theta)n}\).  Even after the whole factor is favorably credited as
 fixed-place gain, the universal saddle-cancellation inequality shows that
 this ledger does not improve the Pr\'evost baseline.  Balanced factorial
 ratios themselves have only \(O(\log n)\) fixed-prime valuation.
 \item The displayed unframed lateral transition of the elementary Euler
 differential operator is gamma-free.  The missing information is the
 canonical framing.  In the mixed arithmetic-Gevrey route, the intersection
 conjecture \(\mathbf E\cap\mathbf D=\Qbar\) yields simultaneous linear
 independence of the dilation family and of all derivatives
 \(1,\Gamma'(1),\Gamma''(1),\ldots\).  These conclusions remain
 conditional.  In the Euler--Mascheroni motive the same extension
 coordinate is the off-diagonal entry of the period matrix.  The explicit
 matched-line lifting problem \cref{prob:euler-tate-fullness} would force
 the \(\Qbar\)-linear independence of \(1,\pi,\gamma\), not merely the
 transcendence of \(\gamma\).
\end{enumerate}

The principal quantitative unresolved problem can now be stated without
spectral language.

\begin{problem}[Common-coefficient separation]
\label{prob:common-coefficient}
Prove that for every \(s\in\Qbar\),
\[
 \limsup_{N\to\infty}
 \frac{-\log|2sY_N^2+Y_N-2|}{2^N}<1,
\]
or establish any equivalent separation that exploits the same coefficient
\(s\) across all \(N\).
\end{problem}

By \cref{thm:critical-slope}, this problem is equivalent to the
transcendence of \(S_2\), hence to the transcendence of \(\gamma\).  It is
therefore a precise quantitative target, not a smaller lemma already within
current fixed-point \(E\)-function theory.  Independently,
\cref{prob:euler-tate-fullness} is a categorical target for the canonical
rank-two Euler extension; its \(\Qbar\)-version would additionally prove
the transcendence of \(\gamma/\pi\).  The mixed-Gevrey intersection
conjecture gives a logically independent conditional route and the stronger
all-order linear-independence package, but it is not presently known.

We emphasize the corresponding negative conclusions.  The algebraic independence of the
finite stages does not determine the arithmetic nature of their limit.  The
Laguerre--P\'olya determinant is an exact spectral carrier but is not proved
arithmetic or \(D\)-finite; the regularized Gamma determinant does not prove
the irrationality of any odd zeta value.  Certificate uniqueness shows that producing more
linear witnesses of the same normalized type creates mutually exclusive
alternatives rather than independent chances to prove transcendence.  The
specific Pr\'evost and Rodrigues families analyzed here do not pay their
complete height ledger, even under the favorable fixed-place accounting
stated above.  The \(p\)-adic convergence theorems prove neither the
irrationality of \(E_p(-1)\), nor Kurepa's conjecture, nor full convergence
for every prime.  The external relation classifications do not close a
famous transcendence conjecture.  Finally, unframed Stokes data forget the
Euler extension coordinate, and the mixed-Gevrey theorems that recover it
are explicitly conditional.
Accordingly, no unconditional theorem in this paper proves the irrationality or
transcendence of \(\gamma\), \(\delta\), \(S_2\), or any odd zeta value; the
exact Euler apertures above state what is still missing.

\appendix

\section{An \texorpdfstring{\(\Ein\)}{Ein}-linear realization of
characteristic-zero algebraic matroids}
\label{app:Ein-matroids}

This appendix is the only place where the paper uses the linearization
theorem of Rosen--Sidman--Theran.  It is kept separate so that the moving
trace and Ramanujan results do not appear to depend on it.

\begin{theorem}[\(\Ein\)-realization of algebraic matroids]
\label{thm:matroid-Ein-realization}
Let \(M\) be a finite characteristic-zero algebraic matroid of rank \(r\)
on a ground set \(E\).  There are a matrix
\(B=(b_{je})\in\Qbar^{r\times E}\) representing \(M\), distinct
\(\beta_1,\ldots,\beta_r\in\Qbar^\times\), and numbers
\begin{equation}\label{eq:matroid-Ein-realization}
 \xi_j=\Ein(\beta_j),\qquad
 W_e=\sum_{j=1}^r b_{je}\xi_j,
\end{equation}
such that, for every \(S\subseteq E\),
\begin{equation}\label{eq:matroid-rank-realization}
 \boxed{\displaystyle
 \trdeg_KK(W_e:e\in S)=r_M(S).}
\end{equation}
Moreover the complete relation ideal is
\begin{equation}\label{eq:matroid-linear-ideal}
 \boxed{\displaystyle
 \ker\!\left(K[X_e:e\in E]\longrightarrow K[W_e:e\in E]\right)
 =\left\langle\sum_{e\in E}c_eX_e:
 c\in\Qbar^E,\ Bc=0\right\rangle.}
\end{equation}
\end{theorem}

\begin{proof}
Extend a characteristic-zero field of algebraic realization to its
algebraic closure; the algebraic-dependence matroid is unchanged.  By
\cite[Thm.~1.10]{RosenSidmanTheran2025}, the resulting matroid is linearly
representable over that algebraically closed field.  Linear
representability of a fixed finite
matroid is described by the vanishing of its nonbasis minors and the
nonvanishing of its basis minors.  Introducing one auxiliary inverse for
the product of the basis minors converts these conditions into polynomial
equations over \(\Q\).  Since they have a solution in an algebraically
closed characteristic-zero field, the Nullstellensatz gives a solution in
\(\Qbar\).  Thus a representing matrix \(B\) may be chosen over \(\Qbar\).

Choose distinct nonzero algebraic \(\beta_j\).  By
\cref{thm:input-independence}, the \(\xi_j\) are algebraically independent
over \(K\).  If the columns of \(B\) indexed by \(S\) have rank \(q\),
choose \(q\) independent columns.  The corresponding \(W_e\) are
algebraically independent by \cref{lem:linear-images}, giving the lower
bound \(q\); every remaining column is a \(\Qbar\)-linear combination of
those columns, giving the upper bound.  This proves
\eqref{eq:matroid-rank-realization}.  Finally the symmetric algebra of the
column span is the polynomial ring on the \(X_e\) modulo the linear forms
coming from \(\ker B\), which proves \eqref{eq:matroid-linear-ideal}.
\end{proof}

\begin{remark}[Exact scope]
The construction realizes the rank function of \(M\) and the linearized
relation ideal \eqref{eq:matroid-linear-ideal}.  It does not reproduce the
nonlinear circuit polynomials of a chosen algebraic presentation of
\(M\), and it is not a claim that such a presentation becomes linear as an
algebraic variety.
\end{remark}

\section{Creative-telescoping certificates for the Euler recurrence}
\label{app:telescopers}

For \(a=n+3\), \(b=n+4\), define meromorphically
\[
 T_{n,k}=b\frac{(a!b!)^2}
 {\Gamma(k+1)^3\Gamma(a-k+1)\Gamma(b-k+1)},
 \qquad
 \mathcal L_{n,k}=3H_k-H_{a-k}-H_{b-k},
\]
where \(H_z=\psi(z+1)+\gamma\).  The regularized endpoint values are
\[
 T_{n,b}\mathcal L_{n,b}=a!,
 \qquad
 T_{n,k}\mathcal L_{n,k}=0\quad(k>b).
\]
Indeed, at \(k=b\) the simple zero of
\(1/\Gamma(a-k+1)\) cancels the simple pole of
\(-H_{a-k}\), with limiting product \(a!\).  For \(k>b\), the two
reciprocal-gamma zeros dominate the at-most-simple digamma pole, so the
regularized product vanishes.
Thus \eqref{eq:qn-closed} and \eqref{eq:pn-closed} are respectively
\(\sum_{k\ge0}T_{n,k}\) and
\(\sum_{k\ge0}T_{n,k}\mathcal L_{n,k}\).

Let
\[
 c_0=-A_n,
 \qquad c_1=B_n,
 \qquad c_2=-C_n,
 \qquad c_3=D_n,
\]
and \(\Delta_ku(k)=u(k+1)-u(k)\).  The certificates are
\begin{align}
 \sum_{j=0}^3c_jT_{n-j,k}
 &=\Delta_k\{T_{n,k}h_{n,k}\},
 \label{eq:q-telescoper}\\
 \sum_{j=0}^3c_jT_{n-j,k}\mathcal L_{n-j,k}
 &=\Delta_k\{T_{n,k}(h_{n,k}\mathcal L_{n,k}+g_{n,k})\}.
 \label{eq:p-telescoper}
\end{align}
Here
\[
 h_{n,k}=-\frac{k^3}{(n+3)^3(n+4)^3}\mathcal H(n,k),
 \qquad
 g_{n,k}=\frac{k^2}{(n+3)^3(n+4)^3}\mathcal G(n,k),
\]
where
\begin{align*}
\mathcal H(n,k)={}&
 (8n^3+75n^2+231n+232)k^3\\
&-(48n^4+562n^3+2436n^2+4626n+3248)k^2\\
&+(96n^5+1388n^4+7924n^3+22278n^2+30772n+16668)k\\
&-(64n^6+1120n^5+8071n^4+30590n^3+64148n^2+70347n+31388),
\end{align*}
and
\begin{align*}
\mathcal G(n,k)={}&
 (48n^3+450n^2+1386n+1392)k^3\\
&-(240n^4+2810n^3+12180n^2+23130n+16240)k^2\\
&+(384n^5+5552n^4+31696n^3+89112n^2+123088n+66672)k\\
&-(192n^6+3360n^5+24213n^4+91770n^3+192444n^2+211041n+94164).
\end{align*}

To verify \eqref{eq:q-telescoper}, divide by \(T_{n,k}\), use
\[
 \frac{T_{n,k+1}}{T_{n,k}}
 =\frac{(n+3-k)(n+4-k)}{(k+1)^3},
\]
and, for \(0\le j\le3\),
\[
 \frac{T_{n-j,k}}{T_{n,k}}
 =\frac{n-j+4}{n+4}
 \frac{(n+3-k)^{\underline j}(n+4-k)^{\underline j}}
 {\{(n+3)^{\underline j}(n+4)^{\underline j}\}^{2}},
\]
where \(x^{\underline j}=x(x-1)\cdots(x-j+1)\), and clear denominators.
For \eqref{eq:p-telescoper}, also use
\[
 \mathcal L_{n,k+1}-\mathcal L_{n,k}
 =\frac3{k+1}+\frac1{n+3-k}+\frac1{n+4-k}.
\]
The remaining shifts are reduced by
\[
 \mathcal L_{n-j,k}-\mathcal L_{n,k}
 =\sum_{r=0}^{j-1}
 \left(\frac1{n+3-k-r}+\frac1{n+4-k-r}\right).
\]
Both become polynomial identities in \(n,k\).  The lower boundary vanishes
because of the factors \(k^3,k^2\), and the reciprocal gamma factors give
the upper boundary.  Summing proves \eqref{eq:challenge-recurrence} for both
closed forms for every \(n\ge0\).  The accompanying source package contains
an exact symbolic verifier for these identities and for
\eqref{eq:cmf-det}--\eqref{eq:cmf-flatness}.

\section{Finite certificate for the integer initial lattice}
\label{app:integer-lattice-certificate}

The official dual transition and cyclic gauge may be taken as
\begin{equation}\label{eq:Rn-certificate}
 \mathsf R(n)=
 \begin{pmatrix}
 6n^2+21n+18&-8n-11&3\\
 3n^3+12n^2+16n+7&-3n^2-7n-4&n+1\\
 (n+1)^4&0&0
 \end{pmatrix}
\end{equation}
and
\begin{equation}\label{eq:Un-certificate}
\mathsf U(n)=
\begin{pmatrix}
1&6n^2+21n+18&15n^4+147n^3+520n^2+792n+440\\
0&3n^3+12n^2+16n+7&10n^5+96n^4+353n^3+625n^2+537n+179\\
0&(n+1)^4&6n^6+57n^5+213n^4+402n^3+408n^2+213n+45
\end{pmatrix}.
\end{equation}
Let the row-companion transfer of \eqref{eq:challenge-recurrence} be
\begin{equation}\label{eq:row-companion-certificate}
 \mathsf C_n^{\mathrm{comp}}=
 \begin{pmatrix}
 0&0&D_n/A_n\\
 1&0&-C_n/A_n\\
 0&1&B_n/A_n
 \end{pmatrix}.
\end{equation}
Direct multiplication gives the exact coboundary
\begin{equation}\label{eq:RU-companion-coboundary}
 \mathsf R(n)\mathsf U(n+1)
 =\mathsf U(n)\mathsf C_n^{\mathrm{comp}}.
\end{equation}
For a raw row \(c\), the integer recurrence state at level \(N\) is
\begin{equation}\label{eq:raw-state-certificate}
 c\,\mathsf R(0)\cdots\mathsf R(N-1)\mathsf U(N).
\end{equation}
Indeed \eqref{eq:RU-companion-coboundary} proves by induction that
\eqref{eq:raw-state-certificate} equals
\((u_{N-3},u_{N-2},u_{N-1})\) when \(c\mathsf U(0)\) is the initial row;
the empty product is used at \(N=0\).
Let
\[
 d_1=(8,-22,10),\qquad d_2=(0,2,-6),\qquad d_3=(0,-1,7).
\]
Direct multiplication gives
\begin{align}
 \frac{d_1}{8}\mathsf U(0)&=(1,0,4),&
 \frac{d_2}{8}\mathsf U(0)&=(0,1,11),&
 \frac{d_3}{8}\mathsf U(0)&=(0,0,17),
 \label{eq:lattice-basis-certificate}
\end{align}
and the finite congruence certificate
\begin{align}
 \mathsf R(0)\cdots\mathsf R(7)&\equiv0\pmod8,
 \label{eq:R8-zero}\\
 d_i\mathsf R(0)\cdots\mathsf R(N-1)\mathsf U(N)&\equiv0\pmod8
 \quad(1\le i\le3,\ 0\le N<8).
 \label{eq:finite-mod8-certificate}
\end{align}
For \(N\ge8\), \eqref{eq:R8-zero} supplies the factor eight; for
\(N<8\), \eqref{eq:finite-mod8-certificate} does.  Hence the three rows in
\eqref{eq:lattice-basis-certificate} generate integral solutions at all
levels, completing the sufficiency argument in
\cref{thm:ram-integer-lattice}.  The source verifier checks every displayed
matrix identity over \(\Z\).

\section{A numerical ledger}
\label{app:numerics}

The following values provide reproducibility checks for the indexing in
\eqref{eq:dN-def}.  They are not used in any proof.

\begin{table}[ht]
\centering
\caption{Positive increments and positive zeros of \(\Delta_2\).}
\begin{tabular}{@{}rcc@{}}
\toprule
\(N\)&\(d_N\)&\(r_N=\sqrt{2/d_N}\)\\
\midrule
0&1.474569389580200&1.1646150303\\
1&\(5.633945512124626\times10^{-1}\)&1.8841205074\\
2&\(9.586804440833174\times10^{-2}\)&4.5674948216\\
3&\(1.339700176289053\times10^{-3}\)&38.637693279\\
4&\(2.954643417466919\times10^{-7}\)&2601.7314103\\
5&\(2.050814843583590\times10^{-14}\)&\(9.8753335\times10^6\)\\
\bottomrule
\end{tabular}
\end{table}

The indexing is tied to the convention \(Y_N\) for \(N\ge1\):
\[
 d_0=\tau_1,\qquad d_N=\tau_{N+1}-\tau_N\quad(N\ge1).
\]
Starting a level sequence at \(N=0\) would give a different, though equally
valid, spectral decomposition.

\begin{thebibliography}{99}
\small
\raggedright

\bibitem{AdamczewskiFaverjon2025}
B.~Adamczewski and C.~Faverjon,
\newblock Algebraic independence measures for values of \(E\)-functions and
\(M\)-functions,
\newblock arXiv:2502.09999 (2025).
\newblock \url{https://arxiv.org/abs/2502.09999}.

\bibitem{DLMF}
NIST Digital Library of Mathematical Functions,
\newblock \S6.12, asymptotic expansions of exponential and logarithmic
integrals,
\newblock release 1.2.7 (2026).
\newblock \url{https://dlmf.nist.gov/6.12}.

\bibitem{Delaygue2018}
E.~Delaygue,
\newblock Arithmetic properties of Ap\'ery-like numbers,
\newblock \emph{Compos. Math.} \textbf{154} (2018), 249--274.
\newblock \url{https://doi.org/10.1112/S0010437X17007552}.

\bibitem{FischlerRivoal2016}
S.~Fischler and T.~Rivoal,
\newblock Arithmetic theory of \(E\)-operators,
\newblock \emph{J. \'Ec. polytech. Math.} \textbf{3} (2016), 31--65.
\newblock \url{https://doi.org/10.5802/jep.28}.

\bibitem{FischlerRivoal2024}
S.~Fischler and T.~Rivoal,
\newblock Values of \(E\)-functions are not Liouville numbers,
\newblock \emph{J. \'Ec. polytech. Math.} \textbf{11} (2024), 1--18.
\newblock \url{https://doi.org/10.5802/jep.249}.

\bibitem{FischlerRivoalMixed2024}
S.~Fischler and T.~Rivoal,
\newblock Relations between values of arithmetic Gevrey series, and
applications to values of the Gamma function,
\newblock \emph{J. Number Theory} \textbf{261} (2024), 36--54.
\newblock \url{https://arxiv.org/abs/2301.13518}.

\bibitem{FresanJossen2024}
J.~Fres\'an and P.~Jossen,
\newblock \emph{Exponential Motives},
\newblock monograph manuscript, version of April 2024.
\newblock \url{http://javier.fresan.perso.math.cnrs.fr/expmot.pdf}.

\bibitem{KEIO2026I}
KEIO,
\newblock Exponential-integral periods across arithmetic, Pad\'e, Hankel,
reciprocal geometry, and Frobenius framing,
\newblock Version 6.3, research manuscript, 108 pages (2026).

\bibitem{Lagarias2013}
J.~C.~Lagarias,
\newblock Euler's constant: Euler's work and modern developments,
\newblock \emph{Bull. Amer. Math. Soc.} \textbf{50} (2013), 527--628.
\newblock \url{https://doi.org/10.1090/S0273-0979-2013-01423-X}.

\bibitem{Levin1996}
B.~Ya.~Levin,
\newblock \emph{Lectures on Entire Functions},
\newblock Translations of Mathematical Monographs 150,
American Mathematical Society, Providence, RI, 1996.

\bibitem{Marcovecchio2009}
R.~Marcovecchio,
\newblock The Rhin--Viola method for \(\log 2\),
\newblock \emph{Acta Arith.} \textbf{139} (2009), 147--184.
\newblock \url{https://doi.org/10.4064/aa139-2-5}.

\bibitem{Marcovecchio2025}
R.~Marcovecchio,
\newblock Multiple Legendre polynomials in Diophantine approximation,
\newblock arXiv:2512.12890 (2025).
\newblock \url{https://arxiv.org/abs/2512.12890}.

\bibitem{Prevost2008}
M.~Pr\'evost,
\newblock A family of criteria for irrationality of Euler's constant,
\newblock \emph{J. Comput. Appl. Math.} \textbf{219} (2008), no.~2,
484--492.
\newblock \url{https://arxiv.org/abs/math/0507231}.

\bibitem{RamanujanChallenge2026}
M.~Shalyt, R.~Kalisch, C.~Schneider, H.~Barkan, E.~Leibtag,
J.~Campbell, S.~Weinbaum, T.~Monderer, A.~Narayanan, and I.~Kaminer,
\newblock The Ramanujan Challenge for AI,
\newblock arXiv:2607.09721 (2026).
\newblock \url{https://arxiv.org/abs/2607.09721}.

\bibitem{Ridout1958}
D.~Ridout,
\newblock The \(p\)-adic generalization of the Thue--Siegel--Roth theorem,
\newblock \emph{Mathematika} \textbf{5} (1958), 40--48.
\newblock \url{https://doi.org/10.1112/S0025579300001339}.

\bibitem{Snodgrass2026}
B.~Snodgrass,
\newblock Periods of \(E\)-operators,
\newblock arXiv:2608.06005v2 (2026).
\newblock \url{https://arxiv.org/abs/2608.06005}.

\bibitem{WolfsVanAssche2025}
W.~Van Assche and T.~Wolfs,
\newblock Rational approximation of Euler's constant using multiple
orthogonal polynomials,
\newblock \emph{Combinatorics and Number Theory} \textbf{14} (2025),
no.~2, 141--162.
\newblock \url{https://doi.org/10.2140/cnt.2025.14.141}.

\bibitem{Akhiezer1965}
N.~I.~Akhiezer,
\newblock \emph{The Classical Moment Problem and Some Related Questions in
Analysis},
\newblock Oliver \& Boyd, Edinburgh, 1965.

\bibitem{Delaygue2022}
\'E.~Delaygue,
\newblock A Lindemann--Weierstrass theorem for \(E\)-functions,
\newblock \emph{J. Reine Angew. Math.} \textbf{820} (2025), 75--85.
\newblock \url{https://doi.org/10.1515/crelle-2024-0090}.

\bibitem{ErnvallMatalaSeppala2022}
A.-M.~Ernvall-Hyt\"onen, T.~Matala-Aho, and L.~Sepp\"al\"a,
\newblock Euler's factorial series, Hardy integral, and continued fractions,
\newblock \emph{J. Number Theory} \textbf{244} (2023), 224--250.
\newblock \url{https://doi.org/10.1016/j.jnt.2022.09.007}.

\bibitem{GrahovacEtAl2024}
D.~Grahovac, A.~Kovtun, N.~N.~Leonenko, and A.~Pepelyshev,
\newblock Dickman type stochastic processes with short- and long-range
dependence,
\newblock \emph{Stochastics} \textbf{97} (2025), no.~6, 722--743.
\newblock \url{https://doi.org/10.1080/17442508.2025.2522789}.

\bibitem{GrossCS}
B.~H.~Gross,
\newblock On the periods of abelian integrals and a formula of Chowla and
Selberg,
\newblock \emph{Invent. Math.} \textbf{45} (1978), no.~2, 193--211.
\newblock \url{https://doi.org/10.1007/BF01390273}.

\bibitem{Petrov1975}
V.~V.~Petrov,
\newblock \emph{Sums of Independent Random Variables},
\newblock Ergebnisse der Mathematik und ihrer Grenzgebiete 82,
Springer, Berlin, 1975.

\bibitem{MurtySaha2015}
M.~R.~Murty and E.~Saha,
\newblock Transcendental values of the incomplete gamma function and
related questions,
\newblock \emph{Arch. Math.} \textbf{105} (2015), 271--283.
\newblock \url{https://doi.org/10.1007/s00013-015-0800-3}.

\bibitem{Pilehrood2013}
Kh.~Hessami Pilehrood and T.~Hessami Pilehrood,
\newblock On a continued fraction expansion for Euler's constant,
\newblock \emph{J. Number Theory} \textbf{133} (2013), no.~2, 769--786.
\newblock \url{https://doi.org/10.1016/j.jnt.2012.08.016}.

\bibitem{PilehroodBell2012}
Kh.~Hessami Pilehrood and T.~Hessami Pilehrood,
\newblock Rational approximations to values of Bell polynomials at points
involving Euler's constant and zeta values,
\newblock \emph{J. Aust. Math. Soc.} \textbf{92} (2012), no.~1, 71--98.
\newblock \url{https://doi.org/10.1017/S1446788712000134}.

\bibitem{PolyaSchur1914}
G.~P\'olya and I.~Schur,
\newblock \"Uber zwei Arten von Faktorenfolgen in der Theorie der
algebraischen Gleichungen,
\newblock \emph{J. Reine Angew. Math.} \textbf{144} (1914), 89--113.

\bibitem{Powers2023}
M.~R.~Powers,
\newblock A sequence of nested exponential random variables with connections
to two constants of Euler,
\newblock arXiv:2302.04605v4 (2023).
\newblock \url{https://arxiv.org/abs/2302.04605}.

\bibitem{Powers2025}
M.~R.~Powers,
\newblock Transcendence results for
\(\Gamma^{(n)}(1)\) and related sequences of generalized constants,
\newblock arXiv:2511.01849v7 (2026).
\newblock \url{https://arxiv.org/abs/2511.01849}.

\bibitem{RamanujanOfficial2026}
Ramanujan Machine Group,
\newblock \emph{Ramanujan Challenge Proofs},
\newblock official solution manuscript, 45 pages (2026).
\newblock
\url{https://ramanujanmachine.com/wp-content/uploads/2026/08/Ramanujan_Challenge_Proofs_02_08.pdf}.
\newblock Supplementary code:
\url{https://github.com/RamanujanMachine/RamanujanChallengeSolutions}.

\bibitem{RosenSidmanTheran2025}
Z.~Rosen, J.~Sidman, and L.~Theran,
\newblock Linearizing algebraic matroids,
\newblock arXiv:2507.07220v3 (2026).
\newblock \url{https://arxiv.org/abs/2507.07220}.

\bibitem{Sondow2003}
J.~Sondow,
\newblock Criteria for irrationality of Euler's constant,
\newblock \emph{Proc. Amer. Math. Soc.} \textbf{131} (2003), no.~11,
3335--3344.
\newblock \url{https://doi.org/10.1090/S0002-9939-03-07081-3}.

\bibitem{VanAsscheWolfs2023}
W.~Van Assche and T.~Wolfs,
\newblock Multiple orthogonal polynomials associated with the exponential
integral,
\newblock \emph{Stud. Appl. Math.} \textbf{151} (2023), no.~2, 411--449.
\newblock \url{https://doi.org/10.1111/sapm.12608}.

\end{thebibliography}

\end{document}
